\chapter{The key identity and the completion of the proof}\label{ch:14}

\providecommand{\chXIVsig}{\sigma^{\natural}}
\providecommand{\chXIVA}{V_0^{+}}
\providecommand{\chXIVB}{V_3^{+}}

This chapter documents the release of the library named in the front matter; every Lean quotation is checked by \texttt{scripts/textbook/sync.py} to occur verbatim (up to whitespace and comments) in the source files of that release. Checks that were run at an earlier state of the library and not re-run for this edition are marked as such.  It is the last mathematical chapter of the Moonshine part of the book, and it is the one in which the construction
of the Moonshine module is finished.  When it ends, the reader will have seen every step of the
proof of the Lean theorem
\begin{leancode}
theorem moonshine :
    ∃ V : VertexOperatorAlgebra Vnat,
      V.Y = Ynat VnatAxioms.wbarL ∧ V.vacuum = vacNat ∧ V.conformalVector = omegaNat ∧
      V.centralCharge = 24 ∧
      V.grading 1 = ⊥ ∧
      Module.finrank ℂ (V.grading 2) = 196884 ∧
      (∀ n : ℕ, (Module.finrank ℂ (V.grading (n : ℤ)) : ℂ) = PowerSeries.coeff n JShift) ∧
      V.IsSimple
\end{leancode}
(\leanpath{VOA/Moonshine/TopPairs.lean}, full name \lean{VOA.Moonshine.moonshine}), which has
\emph{no hypotheses at all}, and whose \texttt{\#print axioms} output is the three standard axioms
of Lean's foundations, \lean{propext}, \lean{Classical.choice} and \lean{Quot.sound}.  The last
section before the referee's checklist (\S\ref{ch14:sec:decode}) is a long, clause-by-clause
fidelity check of this statement — the most careful one in the book, because it is the statement
the whole library exists to prove.  Everything before it explains, on paper, why it is true, and
presents each intermediate statement with its own fidelity check.

\paragraph{The result in one sentence.}  In the wording recommended by the release audit of
2026-10-27 (\texttt{Notes/Textbook/referee/RELEASE\_AUDIT\_MOONSHINE.md}): \emph{we have formalized
in Lean~4 / Mathlib, with no hypotheses and only Lean's standard axioms, the Moonshine module
$\Vnat$: a simple vertex operator algebra structure on the Frenkel--Lepowsky--Meurman space
$V_\Leech^+\oplus(V_\Leech^T)^+$ built from the Leech lattice, with central charge $24$,
$\Vnat_1=0$, $\dim\Vnat_2=196884$, and graded dimension equal, in every degree, to
$J(q)=j(q)-744=q^{-1}+196884\,q+21493760\,q^2+\cdots$; the Monster symmetry $\Aut(\Vnat)=\mathbb M$ is
not part of the formalization.}  Three readings of this sentence are fixed once and for all.
(a) Although \lean{moonshine} is written as an existential, it is a \emph{construction}: the vertex
map, vacuum, conformal vector and central charge are pinned to explicit definitions, the grading is
then forced, and the witness is unique (\S\ref{ch14:sec:decode}).  (b) $J$ is the \emph{formal}
$q$-expansion \lean{JShift} $=E_4^3/(\Delta/q)-744q$; its identification with the analytic function
$j(\tau)-744$ on the upper half plane is classical but is not part of the Lean statement; it is proved separately in Lean
(\lean{qExpansion_qJ}, \lean{hasSum_JShift}, in \leanpath{VOA/Moonshine/JAnalyticQExpansion.lean}).
(c) ``The FLM space'' refers to the carrier; the twisted$\times$twisted block of the vertex
operator is the form-adjoint of the skew intertwiner (the Huang / Dolan--Goddard--Montague
presentation), not FLM's own explicit formula, and the isomorphism with FLM's $\Vnat$ is a paper
argument (uniqueness of simple-current extensions, Dong--Mason, Prop.~5.3), not a Lean theorem.

As everywhere in this book, our concern is the \emph{meaning} of the Lean statements.  Proofs are
explained mathematically; we do not describe how the Lean proof scripts are organised.

\section{Where we stand, and what remains}\label{ch14:sec:intro}

\subsection{The state of the construction at the start of this chapter}

Let us recall what the previous chapters have built.  The carrier of the Moonshine module is
\begin{equation*}
  \Vnat \;=\; V_{\Leech}^{+}\;\oplus\;(V_{\Leech}^{T})^{+},
\end{equation*}
where $V_\Leech^+$ is the fixed-point subalgebra of the Leech lattice vertex operator algebra under
the lift $\theta$ of $-1$ (Chapter~\ref{ch:04}) and $(V_\Leech^T)^+$ is the $\theta$-fixed part of
the twisted module $V_\Leech^T=S[\theta]\otimes T$ (Chapters~\ref{ch:06}--\ref{ch:08}).  In Lean,
$V_\Leech^+$ is the type \lean{SkewVP} (equivalently \lean{VPlus}; both are abbreviations of the
same subtype of the lattice VOA), $(V_\Leech^T)^+$ is the submodule \lean{TP} of
\lean{Tw := FockT 24 ⊗[ℂ] leechT}, and $\Vnat$ is the product type \lean{Vnat := SkewVP × TP}.

Chapter~\ref{ch:09} constructed the two ``off-diagonal'' fields: the skew intertwiner $W(t,z)$,
which sends $V_\Leech^+$ to $(V_\Leech^T)^+$ and is defined by skew-symmetry,
\begin{equation*}
  W(t,z)\varphi \;=\; e^{zL_{-1}}\,Y^{T}(\varphi,-z)\,t,
\end{equation*}
and its adjoint $\overline W(t,z)$, which goes the other way and is defined by transposition
through the two invariant bilinear forms $B_V$ (\lean{leechBV}) on $V_\Leech^+$ and
\lean{TPForm} on $(V_\Leech^T)^+$.  Chapter~\ref{ch:10} assembled the block vertex operator
\begin{equation}\label{ch14:eq:block}
  Y^{\natural}\bigl((u,t),z\bigr) \;=\;
  \begin{pmatrix} Y(u,z) & \overline W(t,z)\\ W(t,z) & Y^{T}(u,z)\end{pmatrix}
\end{equation}
(DGM (4.17); Lean \lean{Ynat wbarL}) and proved that it defines a vertex operator algebra with
central charge $24$ \emph{provided} one family of locality relations holds, the
``twisted$\times$twisted'' relation that Dolan--Goddard--Montague call (4.23):
\begin{equation}\label{ch14:eq:423}
  W(t_1,z_1)\,\overline W(t_2,z_2)\;\sim\;W(t_2,z_2)\,\overline W(t_1,z_1)
  \qquad\text{on }(V_\Leech^{T})^{+},
\end{equation}
where $\sim$ means equality after multiplication by a sufficiently high power of $(z_1-z_2)$.
In Lean this is the proposition \lean{TPBlockLocal leechWbar t₁ t₂}.  The companion relation
(4.22), the analogous statement on $V_\Leech^+$, was proved unconditionally
(\lean{leechWbar_VPBlockLocal}); simplicity was proved too, via nondegeneracy of an invariant form
(\lean{SimpleNat.vnatVOA_isSimple}).  The result was \lean{moonshine_of_frameLocality}: \emph{all}
targets, conditional on (4.23) for four explicit pairs of twisted vectors.

Chapters~\ref{ch:11}--\ref{ch:13} then built the tool with which FLM prove (4.23): the
\emph{triality}.  The Leech lattice contains a sublattice $\Leech_0$ of index two (the Golay frame
sublattice), and accordingly
\begin{equation*}
  V_\Leech^+ = \chXIVA\oplus\chXIVB,\qquad
  (V_\Leech^T)^+ = T_0^+\oplus T_1^+ ,
\end{equation*}
four $\chXIVA$-modules (Lean: \lean{frameSector 0}, \lean{frameSector 1}, \lean{TPframe 0},
\lean{TPframe 1}).  FLM's key insight is that there is an automorphism $\sigma$ of the subalgebra
$\chXIVA$ that does not extend to $V_\Leech^+$, but which permutes the three nontrivial sectors:
twisting $\chXIVB$ by $\sigma$ gives one twisted sector, and the other twisted sector is
$\sigma$-stable.  Chapter~\ref{ch:11} constructed $\sigma$ on $\chXIVA$ (\lean{v0Sigma});
Chapters~\ref{ch:12}--\ref{ch:13} constructed the twisted halves, the $\sigma$-twisted
$\chXIVA$-intertwiners
\begin{equation*}
  \Phi:\chXIVB\to T_s^+,\qquad \Theta:T_{s+1}^+\to T_{s+1}^+,
\end{equation*}
unconditionally (\lean{SigmaTop.topSigma_both}), using Zhu-algebra generation
(\lean{zhuGenerated_weightTwo}) to reduce their existence to finite top-level data.

\subsection{What this chapter does}

The idea of FLM is to \emph{conjugate}.  Assemble the sector maps into one linear involution
$\chXIVsig$ of $\Vnat$.  If $\chXIVsig$ were an automorphism of the (would-be) vertex algebra,
then for $v\in\chXIVB$,
\begin{equation*}
  \chXIVsig\,Y^\natural(v,z)\,\chXIVsig \;=\; Y^\natural(\chXIVsig v,z),
\end{equation*}
and the right-hand side is a twisted field $Y^\natural((0,\Phi v),z)$.  The left-hand side is local
with itself because the untwisted field $Y^\natural(v,z)$ is (proved in Chapter~\ref{ch:10}).  So
the twisted field $Y^\natural((0,x),z)$, $x=\Phi v$, would be local with itself, which contains
(4.23) for the pair $(x,x)$; and one generating twisted vector suffices.

Of course, we cannot assume that $\chXIVsig$ is an automorphism; that is what is to be proved.
The honest content is an \emph{identification of two fields}: the conjugated field
$L_v=\chXIVsig Y^\natural(v,z)\chXIVsig$ and the block field $R_v=Y^\natural(\chXIVsig v,z)$.
The chapter proceeds as follows.

\begin{enumerate}[label=(\arabic*)]
\item \S\ref{ch14:sec:conj}: the conjugation set-up.  Both $L_v$ and $R_v$ are local with all
  $\chXIVA$-fields and create the same state from the vacuum, so their difference is determined by
  its values on the cyclic vectors of the four sectors.  The remaining agreements are three
  ``three-point identities'' (\lean{SigmaKey}); if they hold, then (4.23) holds for all pairs and
  $\Vnat$ is constructed (\lean{moonshine_of_sigma}).
\item \S\ref{ch14:sec:sigmakey}: adjunction against the invariant forms reduces the three
  identities to two \emph{isometry} statements and one \emph{self-coupling} identity
  \lean{KeySkew}; the isometries reduce further to finite top-level statements
  (\lean{moonshine_of_top}).
\item \S\ref{ch14:sec:assemble}: the explicit triality data of Chapter~\ref{ch:13} are assembled into
  \lean{mkData c ε}, the isometries are proved (with the normalisation $\mu=+1$ \emph{proved}), and
  the only remaining hypothesis is \lean{KeySkew0 ε} for one sign $\varepsilon$
  (\lean{moonshine_of_keySkew0}).
\item \S\ref{ch14:sec:keyskew}: a lemma of Li type shows that \lean{KeySkew} is controlled by the
  top-output components of a ``block defect'' (\lean{keySkew_of_keyTopV}); that residual is still
  infinite.
\item \S\ref{ch14:sec:fz}: Frenkel--Zhu bimodule theory makes it finite.  The bimodule
  $A(\chXIVB)$ is generated by its top level (\lean{bimodGen_V3}), and the transposed top-level
  identities need only the commutator formula, so the intertwining Jacobi identity of $W$ — which
  the library does not have — is never needed (\lean{keySkew_of_FZ}).
\item \S\ref{ch14:sec:toppairs}: the last finite identity, $\Theta\,o(v)\Phi b=o(b)\Phi v$ on the
  $49152$-dimensional top level, is proved by an explicit sign computation in Leech coordinates
  (\lean{topPairs_one}, with $\varepsilon=+1$).  This yields \lean{moonshine}.
\item \S\ref{ch14:sec:decode}: the final theorem, read clause by clause; then the referee's
  checklist.
\end{enumerate}

\begin{remark}[The chain of reductions at a glance]\label{ch14:rem:chain}
Each implication below is a Lean theorem presented in this chapter; the top line is the goal, the
bottom lines are what is finally established by direct argument.
\begin{align*}
  \text{(4.23) for all pairs}
  &\;\Longleftarrow\; \lean{SigmaKey}
   && (\lean{tpBlockLocal_of_sigma})\\
  \lean{SigmaKey}
  &\;\Longleftarrow\; \lean{PhiTop}\wedge\lean{ThetaTop}\wedge\lean{KeySkew}
   && (\lean{sigmaKey_of_top})\\
  \lean{PhiTop}\wedge\lean{ThetaTop}
  &\;\Longleftarrow\; \text{explicit data }\lean{mkData}
   && (\lean{phiIso_mkData},\ \lean{thetaTop_mkData})\\
  \lean{KeySkew}
  &\;\Longleftarrow\; \lean{KeyTopV}
   && (\lean{keySkew_of_keyTopV})\\
  \lean{KeyTopV}
  &\;\Longleftarrow\; \lean{BimodGenV3}\wedge\lean{TopPairs}
   && (\lean{keyTopV_of_FZ})\\
  \lean{BimodGenV3}
  &\;\Longleftarrow\; \text{lattice combinatorics}
   && (\lean{fzGenerated_MV3})\\
  \lean{TopPairs}
  &\;\Longleftarrow\; \text{sign arithmetic in }\Leech
   && (\lean{topPairs_one}).
\end{align*}
\end{remark}

\subsection{Notation used throughout the chapter}\label{ch14:sec:notation}

We fix the following abbreviations, which match the Lean files and the route notes
\leanpath{Notes/SigmaKey.md}, \leanpath{Notes/KeySkew.md}.
\begin{itemize}
\item $A=\chXIVA$ (Lean \lean{V0Sub.carrier}, equal to \lean{frameSector 0}); as a VOA in its own
  right it is \lean{V0P}.  $B=\chXIVB$ (\lean{frameSector 1}).  $C=T_s^+$ and $D'=T_{s+1}^+$
  (\lean{TPframe s}, \lean{TPframe (s+1)}), where $s\in\Z/2$ labels the twisted sector that the
  triality pairs with $\chXIVB$.  The label is part of the data because the library's definitions do
  not single out which of the two twisted frame sectors is FLM's $T_0^+$
  (\lean{golayE_not_determined}, Chapter~\ref{ch:12}).
\item For $u\in A$ and $n\in\Z$, \lean{actV u n} and \lean{actT u n} are the modes $u_{(n)}$ of $u$
  acting on $V_\Leech^+$ and on \lean{TP}.  Recall the library convention
  \lean{mode f n = f (-n - 1)} (Definition~\ref{l01:def:mode}): \lean{mode (Y u) n} is the mathematician's
  $u_{(n)}$, whereas the raw coefficient \lean{Y u k} is the coefficient of $z^{k}$ in
  $Y(u,z)=\sum_n u_{(n)}z^{-n-1}$, i.e.\ $u_{(-k-1)}$.  The library records this as
  \lean{Y_eq_actV}: \lean{leechVOAPlus.Y u k x = actV u (-k - 1) x}.  \emph{Every} statement of the
  form \lean{∀ p, … Y v p …} below is therefore a statement about all coefficients of a series, and
  the reindexing $p\leftrightarrow -p-1$ never changes its content.
\item $\sigma=$ \lean{v0Sigma} $\in\Aut(\chXIVA)$; \lean{sig u} is its value on $u$.
\item The cyclic vectors: $x_\beta=e^\beta+e^{-\beta}\in\chXIVB$ for $\beta\in\Leech\setminus\Leech_0$
  (\lean{frameX0P β}); $x_0(t)=h_0(-\tfrac12)\mathbb 1\otimes t\in T^+$ for $t\in T$
  (\lean{x0T t}), which lies in $T_s^+$ when $t\in T_s$ (\lean{x0T_mem}) and has weight $2$.
\item $o(v)$ is the zero mode $v_{(\wt v-1)}$ of a homogeneous $v$; it preserves weight.
\item The weight-$2$ ``tops'': $\Omega_{V_3}=(\chXIVB)_2$ (\lean{OmegaV3}, of dimension $49152$) and
  $\Omega_T(s)=(T_s^+)_2$ (\lean{OmegaT s}).  Both $\chXIVB$ and $T^+$ have \emph{no} states of weight
  below $2$.
\end{itemize}

\section{\texorpdfstring{The would-be triality $\chXIVsig$ and conjugation of fields}{The would-be triality sigma-natural and conjugation of fields}}\label{ch14:sec:conj}

\subsection{The physical picture}

In the language of two-dimensional conformal field theory, $\Vnat$ is the chiral algebra of a
$\Z/2$-orbifold of 24 free bosons compactified on the Leech torus.  The untwisted sector
$V_\Leech^+$ contains the $\theta$-invariant operators of the original theory; the twisted sector
$(V_\Leech^T)^+$ contains the $\theta$-invariant \emph{twist fields}, created by inserting a branch
cut around which the bosons change sign.  The operator product of two twist fields lands back in
the untwisted sector, and the content of (4.23) is that these operator products are mutually local:
the twist fields of $\Vnat$ behave like ordinary, single-valued local operators.  In general an
orbifold of a lattice theory by $-1$ produces a consistent chiral algebra only after a
subtle sign analysis; FLM's triality makes this analysis unnecessary by exhibiting a
\emph{symmetry} of the would-be theory that turns twist fields into untwisted fields, for which
locality is already known.  The physicist's slogan is: ``the twisted sector is not special; there is
a hidden symmetry that rotates it into the untwisted sector.''  This section turns the slogan into
a precise reduction.

\subsection{One twisted vector generates the twisted sector}

The first ingredient is a cyclicity statement: a single twisted vector of the simplest kind
generates all of $(V_\Leech^T)^+$ under the modes of $V_\Leech^+$.

\begin{definition}[Mode span]\label{ch14:def:modespan}
Let $g=(g_i)_{i\in\kappa}$ be a family of vectors of $T^+=(V_\Leech^T)^+$.  We say that $g$
\emph{mode-spans} $T^+$ if
\begin{equation*}
  T^{+}=\operatorname{span}_\C\{\,u_{(n)}g_i\;:\;u\in V_\Leech^+,\ n\in\Z,\ i\in\kappa\,\}.
\end{equation*}
\end{definition}

\begin{leancode}
def ModeSpan : Prop :=
  Submodule.span ℂ {x : TP | ∃ (u : SkewVP) (n : ℤ) (i : κ), x = mode (plusModule.YM u) n (g i)} = ⊤
\end{leancode}
(\leanpath{VOA/Moonshine/M8h0Gap.lean}, \lean{VOA.Moonshine.M8h0.ModeSpan}, with $g:\kappa\to$
\lean{TP} a section variable.)

\begin{dictionary}
\textbf{Decoding.} \lean{plusModule} is the $V_\Leech^+$-module structure on \lean{TP} constructed in
Chapter~\ref{ch:08}; \lean{plusModule.YM u} is the formal series $Y^T(u,z)$ acting on \lean{TP}, and
\lean{mode (plusModule.YM u) n} is its $n$-th mode $u_{(n)}$ by the convention
\lean{mode f n = f (-n-1)}.  \lean{Submodule.span ℂ S} is the smallest complex subspace containing
$S$, and \lean{= ⊤} says it is the whole space \lean{TP}.

\textbf{Same statement.}  The Lean set is exactly the set of vectors $u_{(n)}g_i$.  No hypothesis.
\end{dictionary}
\begin{physics}
A family of twist fields generates the twisted sector if all twisted states arise from them by OPE with untwisted operators: the twisted sector is an irreducible ``module'' in the CFT sense, built over a few ground states.
\end{physics}


\begin{theorem}[One vector generates $T^+$]\label{ch14:thm:modespan}
For every nonzero $t_0\in T$, the single vector $x_0(t_0)=h_0(-\tfrac12)\mathbb 1\otimes t_0$
mode-spans $(V_\Leech^T)^+$.
\end{theorem}

\begin{leancode}
theorem modeSpan_x0T {t₀ : TLeechT} (h0 : t₀ ≠ 0) : ModeSpan (fun _ : Unit => x0T t₀)
\end{leancode}
(\leanpath{VOA/Moonshine/M8hKeySpan.lean}, \lean{VOA.Moonshine.M8hKey.modeSpan_x0T}.)

\begin{dictionary}
\textbf{Decoding.}  \lean{TLeechT} is the $2^{12}$-dimensional irreducible module $T$ of the
twisted group algebra of $\Leech$ (Chapter~\ref{ch:08}, \lean{finrank_leechT = 4096}).
\lean{x0T t₀} is the vector $h_0(-\tfrac12)\mathbb 1\otimes t_0$ regarded as an element of
\lean{TP}; here $h_0$ is the first orthonormal coordinate direction of $\C\otimes\Leech$ and
$h_0(-\tfrac12)$ the corresponding half-integral twisted Heisenberg creation operator.  The family
is indexed by the one-point type \lean{Unit}, so ``mode-spans'' means: the $u_{(n)}x_0(t_0)$ span.

\textbf{Hypothesis.}  Only $t_0\ne0$, which is obviously necessary.  Note that $t_0$ is \emph{not}
required to lie in a frame sector $T_s$: the theorem uses all of $V_\Leech^+$, not only $\chXIVA$.

\textbf{Same statement.}  Yes.

\textbf{Could it be vacuous?}  No: $T\ne0$, and the conclusion is a genuine spanning statement
about a nonzero space.
\end{dictionary}
\begin{physics}
One twist field suffices: inserting untwisted operators (oscillators $\partial X$ and vertex operators $e^{i\alpha\cdot X}$) next to a single ground-state twist field produces every twisted state. Physically the $e^\alpha$ act on the $2^{12}$ ground states as the Clifford-type algebra of the zero modes, which acts irreducibly.
\end{physics}


\paragraph{Proof on paper.}  Let $Z\subseteq V_\Leech^T$ be the span of all $Y^T(u)_{(k/2)}x_0$,
with $u$ running over the whole lattice VOA $V_\Leech$ (twisted modes are half-integral).  $Z$ is
stable under every twisted mode.  First, the annihilation operator $h_0(\tfrac12)$ removes the
oscillator: $h_0(\tfrac12)x_0=\tfrac12\,\mathbb 1\otimes t_0$, so $\mathbb 1\otimes t_0\in Z$.  Next,
the twisted vertex operator of $e^\alpha$ acts on $\mathbb 1\otimes t$ with a lowest coefficient
proportional to $\mathbb 1\otimes e_\alpha t$, so $\mathbb 1\otimes e_\alpha t_0\in Z$ for all
$\alpha\in\Leech$.  Since $T$ is a \emph{simple} module for the twisted group algebra
(Chapter~\ref{ch:08}), the $e_\alpha t_0$ span $T$; hence $\mathbb 1\otimes T\subseteq Z$.  The
Heisenberg creation operators then produce $S[\theta]\otimes T$, so $Z=V_\Leech^T$.  Finally,
project to the $\theta$-invariant parts: the $\theta$-fixed part of a mode $u_{(n)}x_0$ is the mode
of the $\theta$-fixed part of $u$ applied to the $\theta$-fixed $x_0$, so $T^+$ is spanned by modes of
$V_\Leech^+$ applied to $x_0$. \qed

\begin{status}
Unconditional.  \texttt{\#print axioms}: \lean{[propext, Classical.choice, Quot.sound]} (the file is
imported by \lean{VOA.Moonshine.moonshine}, whose axiom list is this).
\end{status}

\subsection{Conjugating a field by an involution}

\begin{definition}[Conjugate field]\label{ch14:def:conjF}
Let $V$ be a complex vector space, $\sigma\in\End V$ and $A(z)=\sum_k A_k z^k$ a formal
distribution with coefficients in $\End V$.  The \emph{conjugate} of $A$ by $\sigma$ is
$(\sigma A\sigma)(z)=\sum_k \sigma A_k\sigma\, z^k$.
\end{definition}
\begin{leancode}
def conjF (σ : End ℂ V) (A : FormalDist (End ℂ V)) : FormalDist (End ℂ V) :=
  fun k => σ * A k * σ
\end{leancode}
(\leanpath{VOA/Moonshine/M8hKeySigma.lean}, \lean{VOA.Moonshine.M8hKey.conjF}.)

\begin{dictionary}
\lean{FormalDist X} is the type of functions $\Z\to X$, read as formal series
$\sum_{k\in\Z}A_kz^k$ (Definition~\ref{l01:def:formaldist}); \lean{A k} is the coefficient of $z^k$.  The product
\lean{σ * A k * σ} is composition in the algebra \lean{End ℂ V}.  Exactly the textbook definition.
\end{dictionary}
\begin{physics}
Conjugating an operator by a symmetry, $\mathcal O\mapsto g\,\mathcal O\,g^{-1}$: the transformed field. For a genuine symmetry of the CFT, the transformed field of $\phi(z)$ is the field of the transformed state.
\end{physics}


\begin{proposition}\label{ch14:prop:conj}
Let $\sigma^2=1$.  (a) If $A$ and $B$ are mutually local, so are $\sigma A\sigma$ and $\sigma B\sigma$.
(b) If $\sigma$ commutes with an operator $T$ and $A$ is $T$-covariant, $[T,A(z)]=\partial_zA(z)$,
then so is $\sigma A\sigma$.
\end{proposition}
\begin{leancode}
theorem local_conjF {σ : End ℂ V} (hσ : σ * σ = 1) {A B : FormalDist (End ℂ V)}
    (h : Local A B) : Local (conjF σ A) (conjF σ B)

theorem dCovariant_conjF {σ T : End ℂ V} (hT : σ * T = T * σ)
    {A : FormalDist (End ℂ V)} (h : DCovariant T A) : DCovariant T (conjF σ A)
\end{leancode}
(\leanpath{VOA/Moonshine/M8hKeySigma.lean}.)

\begin{dictionary}
\textbf{Decoding.}  \lean{Local A B} (Definition~\ref{l01:def:bracket}) means: there is $N\in\N$ with
$(z_1-z_2)^N[A(z_1),B(z_2)]=0$ coefficientwise.  \lean{DCovariant T A} means
$[T,A_k]=$ the coefficient of $z^k$ in $\partial_z A$, i.e.\ translation covariance with respect to
$T$ (which will be $L_{-1}$).

\textbf{Hypotheses.}  $\sigma^2=1$ in (a); in (b), $\sigma T=T\sigma$.  Note (b) does not need
$\sigma^2=1$.  Both are the textbook hypotheses.
\end{dictionary}
\begin{physics}
Locality and translation covariance are preserved by any operator that commutes with them in the obvious sense: a change of basis on the Hilbert space does not create new singularities in OPEs.
\end{physics}


\paragraph{Proof on paper.}  (a) $[\sigma A_m\sigma,\sigma B_n\sigma]=\sigma[A_m,B_n]\sigma$ because
$\sigma\sigma=1$ in the middle; multiplying by $(z_1-z_2)^N$ commutes with the linear map
$X\mapsto\sigma X\sigma$.  (b) $[T,\sigma A_k\sigma]=\sigma[T,A_k]\sigma$ since $\sigma$ commutes
with $T$. \qed

\begin{status}Unconditional; standard axioms only.\end{status}

\subsection{The sector decomposition and the triality data}

The splittings $V_\Leech^+=\chXIVA\oplus\chXIVB$ and $T^+=T_0^+\oplus T_1^+$ were established in
Chapter~\ref{ch:11} (\lean{frameSector_isCompl}, \lean{TPframe_isCompl}).  Consequently a linear map
out of $\Vnat$ is determined by its four restrictions (\lean{vnat_ext}); this ``four-sector
extensionality'' is used constantly below.

\begin{definition}[$\sigma$-twisted intertwiners]\label{ch14:def:twint}
Let $M,N$ be $\chXIVA$-modules (here: subspaces of $V_\Leech^+$ or $T^+$ with the restricted mode
action) and $X\subseteq M$ a subspace.  A linear map $\Phi:M\to N$ is a \emph{$\sigma$-twisted
intertwiner on $X$} if
\begin{equation*}
  \Phi(u_{(n)}x)=(\sigma u)_{(n)}\Phi(x)\qquad\text{for all }u\in\chXIVA,\ n\in\Z,\ x\in X.
\end{equation*}
\end{definition}
\begin{leancode}
def IsTwInt (ρ : A → ℤ → M →ₗ[ℂ] M) (τ : A → ℤ → N →ₗ[ℂ] N) (g : A → A) (S : Set A)
    (X : Submodule ℂ M) (Φ : M →ₗ[ℂ] N) : Prop :=
  ∀ u ∈ S, ∀ (n : ℤ), ∀ x ∈ X, Φ (ρ u n x) = τ (g u) n (Φ x)

def IsSigmaV3T (Φ : VPlus →ₗ[ℂ] TP) : Prop :=
  IsTwInt actV actT sig Set.univ (frameSector 1) Φ
def IsSigmaTT (s : ZMod 2) (Φ : TP →ₗ[ℂ] TP) : Prop :=
  IsTwInt actT actT sig Set.univ (TPframe s) Φ
def IsSigmaTV (s : ZMod 2) (Φ : TP →ₗ[ℂ] VPlus) : Prop :=
  IsTwInt actT actV sig Set.univ (TPframe s) Φ
\end{leancode}
(\leanpath{VOA/Moonshine/TrialityTwGeneric.lean}, \leanpath{VOA/Moonshine/TrialityTwSigma.lean}.)

\begin{dictionary}
\textbf{Decoding.}  \lean{ρ u n}, \lean{τ u n} are the $n$-th modes of $u$ on source and target;
\lean{g = sig} is $\sigma$; \lean{S = Set.univ} means ``all $u\in\chXIVA$''.  The map $\Phi$ is defined
on the whole ambient space ($V_\Leech^+$ or \lean{TP}), but the intertwining condition is only
imposed for $x$ in the indicated sector $X$; this is harmless because only the restriction to $X$
matters later.  \lean{M →ₗ[ℂ] N} is the type of $\C$-linear maps.

\textbf{Same statement.}  Yes.  Note the twist is on the \emph{target} side:
$\Phi\circ u_{(n)}=(\sigma u)_{(n)}\circ\Phi$; with $\sigma^2=1$ this is equivalent to the other
convention.
\end{dictionary}
\begin{physics}
A map between sectors that commutes with the chiral algebra $\chXIVA$ up to the outer automorphism $\sigma$: in CFT terms, $\Phi$ realises the action of a symmetry of the subtheory $\chXIVA$ that permutes its superselection sectors.
\end{physics}


\begin{definition}[Triality data]\label{ch14:def:sigmadata}
\emph{Triality data} consist of a label $s\in\Z/2$ and three linear maps
\begin{equation*}
  \Phi:V_\Leech^+\to T^+,\qquad \Psi:T^+\to V_\Leech^+,\qquad \Theta:T^+\to T^+
\end{equation*}
such that $\Phi$ is a $\sigma$-twisted intertwiner on $\chXIVB$, $\Psi$ on $T_s^+$, $\Theta$ on
$T_{s+1}^+$; $\Phi(\chXIVB)\subseteq T_s^+$, $\Psi(T_s^+)\subseteq\chXIVB$,
$\Theta(T_{s+1}^+)\subseteq T_{s+1}^+$; and $\Psi\Phi=\mathrm{id}$ on $\chXIVB$, $\Phi\Psi=\mathrm{id}$
on $T_s^+$, $\Theta^2=\mathrm{id}$ on $T_{s+1}^+$.
\end{definition}
\begin{leancode}
structure SigmaData where
  /-- FLM's `T₀`-label: `σ(V₃⁺) = TPframe s`. -/
  s : ZMod 2
  /-- `σ : V₃⁺ → T_s⁺`. -/
  Φ : VPlus →ₗ[ℂ] TP
  /-- `σ : T_s⁺ → V₃⁺`. -/
  Ψ : TP →ₗ[ℂ] VPlus
  /-- `σ : T_{s+1}⁺ → T_{s+1}⁺`. -/
  Θ : TP →ₗ[ℂ] TP
  isSigma_Φ : IsSigmaV3T Φ
  isSigma_Ψ : IsSigmaTV s Ψ
  isSigma_Θ : IsSigmaTT (s + 1) Θ
  Φ_mem : ∀ x ∈ frameSector 1, Φ x ∈ TPframe s
  Ψ_mem : ∀ t ∈ TPframe s, Ψ t ∈ frameSector 1
  Θ_mem : ∀ t ∈ TPframe (s + 1), Θ t ∈ TPframe (s + 1)
  ΨΦ : ∀ x ∈ frameSector 1, Ψ (Φ x) = x
  ΦΨ : ∀ t ∈ TPframe s, Φ (Ψ t) = t
  ΘΘ : ∀ t ∈ TPframe (s + 1), Θ (Θ t) = t
\end{leancode}
(\leanpath{VOA/Moonshine/M8hKeySigma.lean}, \lean{VOA.Moonshine.M8hKey.SigmaData}.)

\begin{dictionary}
\textbf{Decoding.}  A Lean \lean{structure} is a record: a term of type \lean{SigmaData} is a tuple
of the data fields (\lean{s}, \lean{Φ}, \lean{Ψ}, \lean{Θ}) together with proofs of the listed
properties.  \lean{ZMod 2} is $\Z/2$.  Every field corresponds to one clause of
Definition~\ref{ch14:def:sigmadata}; nothing more is required.

\textbf{What is \emph{not} required.}  \lean{SigmaData} does not ask that $\Phi$, $\Theta$ preserve
the invariant forms, nor any compatibility with the fields $W$, $\overline W$.  It is thus a
\emph{weak} notion: the free normalisations $\Phi\mapsto c\Phi$, $\Psi\mapsto c^{-1}\Psi$ ($c\ne0$) and
$\Theta\mapsto-\Theta$ preserve it.  These freedoms are pinned later.

\textbf{Could it be vacuous?}  A structure whose fields are contradictory would make every theorem
``for all $D$: \ldots'' empty.  This one is inhabited: \S\ref{ch14:sec:assemble} constructs an
explicit term \lean{mkData c ε}.
\end{dictionary}
\begin{physics}
The candidate symmetry of the full orbifold theory, defined sector by sector: it is FLM's triality, which mixes the untwisted odd sector with a twisted sector. The physicist's expectation (to be proved) is that it is a symmetry of all OPEs, not only of those of $\chXIVA$.
\end{physics}


\begin{definition}[The would-be triality $\chXIVsig$]\label{ch14:def:signat}
Given triality data $D$, let $\chXIVsig=\chXIVsig_D\in\End\Vnat$ be the operator acting by
\begin{center}
\begin{tabular}{lll}
\toprule
sector & action of $\chXIVsig$ & lands in \\ \midrule
$\chXIVA$ & $\sigma$ & $\chXIVA$\\
$\chXIVB$ & $\Phi$ & $T_s^+$\\
$T_s^+$ & $\Psi$ & $\chXIVB$\\
$T_{s+1}^+$ & $\Theta$ & $T_{s+1}^+$\\ \bottomrule
\end{tabular}
\end{center}
extended linearly using the four-sector decomposition.
\end{definition}
\begin{leancode}
noncomputable def sig0 : VPlus →ₗ[ℂ] VPlus :=
  V0Sub.carrier.subtype ∘ₗ v0Sigma.toLinearEquiv.toLinearMap ∘ₗ
    V0Sub.carrier.projectionOnto (frameSector 1) V0_isCompl

noncomputable def sigNat : End ℂ Vnat :=
  blockOp sig0
    (D.Ψ ∘ₗ (TPframe D.s).projection (TPframe (D.s + 1)) (TPframe_isCompl' D.s))
    (D.Φ ∘ₗ (frameSector 1).projection V0Sub.carrier V0_isCompl.symm)
    (D.Θ ∘ₗ (TPframe (D.s + 1)).projection (TPframe D.s) (TPframe_isCompl' D.s).symm)
\end{leancode}
(\leanpath{VOA/Moonshine/M8hKeySigma.lean}, \lean{VOA.Moonshine.M8hKey.sigNat}.)

\begin{dictionary}
\textbf{Decoding.}  \lean{blockOp a c d b} is the $2\times2$ block operator
$\begin{pmatrix}a&c\\d&b\end{pmatrix}$ on $V_\Leech^+\times T^+$ (Chapter~\ref{ch:10}).
\lean{sig0} projects onto $\chXIVA$ along $\chXIVB$, applies $\sigma$, and includes back; it is $\sigma$ on
$\chXIVA$ and $0$ on $\chXIVB$.  The top-right block is $\Psi$ composed with the projection of $T^+$ onto
$T_s^+$ along $T_{s+1}^+$; the bottom-left block is $\Phi$ after projecting $V_\Leech^+$ onto $\chXIVB$;
the bottom-right is $\Theta$ after projecting onto $T_{s+1}^+$.  The Lean theorems
\lean{sigNat_V0}, \lean{sigNat_V3}, \lean{sigNat_Ts}, \lean{sigNat_Ts1} state exactly the four rows of the
table.  Same definition as the table.
\end{dictionary}
\begin{physics}
$\chXIVsig$ is the would-be symmetry operator on the full Hilbert space $\Vnat$; in the Moonshine CFT it is the involution that, together with the obvious symmetries, generates the Monster (heuristic; not stated in Lean).
\end{physics}


\begin{theorem}[$\chXIVsig$ is an involution compatible with $\chXIVA$]\label{ch14:thm:signat}
For triality data $D$:
\begin{enumerate}[label=(\alph*)]
\item $(\chXIVsig)^2=1$;
\item for every $u\in\chXIVA$ and $k\in\Z$: $\chXIVsig\,Y^\natural(u)_k=Y^\natural(\sigma u)_k\,\chXIVsig$,
  where $Y^\natural(u)$ is the diagonal field $\operatorname{diag}(Y(u,z),Y^T(u,z))$;
  consequently $\chXIVsig Y^\natural(u)\chXIVsig=Y^\natural(\sigma u)$;
\item $\chXIVsig$ commutes with the translation operator $L_{-1}$ of $\Vnat$, and fixes the vacuum.
\end{enumerate}
\end{theorem}
\begin{leancode}
theorem sigNat_sq : sigNat D * sigNat D = 1

theorem sigNat_YnatU_V0 (u : V0Sub.carrier) (k : ℤ) :
    sigNat D * YnatU (u : VPlus) k = YnatU ((sig u : V0Sub.carrier) : VPlus) k * sigNat D

theorem conjF_YnatU_V0 (u : V0Sub.carrier) :
    conjF (sigNat D) (YnatU (u : VPlus)) = YnatU ((sig u : V0Sub.carrier) : VPlus)

theorem sigNat_Tnat : sigNat D * Tnat = Tnat * sigNat D

theorem sigNat_vac : sigNat D vacNat = vacNat
\end{leancode}
(\leanpath{VOA/Moonshine/M8hKeySigma.lean}.)

\begin{dictionary}
\textbf{Decoding.}  \lean{YnatU u} is the field $Y^\natural((u,0),z)$, the diagonal block
$\operatorname{diag}(Y(u,z),Y^T(u,z))$ (Chapter~\ref{ch:10}, \lean{YnatU_apply}).  \lean{YnatU u k}
is its coefficient of $z^k$.  \lean{Tnat} is the translation operator $L_{-1}\oplus L_{-1}^T$ of $\Vnat$;
\lean{vacNat} is $(\mathbb 1,0)$.  The coercion \lean{(u : VPlus)} views $u\in\chXIVA$ as an element of
$V_\Leech^+$.

\textbf{Hypotheses.}  Only the triality data $D$ (a section variable).  The statements are exactly
(a)--(c).
\end{dictionary}
\begin{physics}
The candidate symmetry is an involution, preserves the vacuum and the stress tensor, and is a symmetry of the subalgebra $\chXIVA$. What remains is to show it respects the OPEs involving the other three sectors.
\end{physics}


\paragraph{Proof on paper.}  Check each identity on each of the four sectors.  (a) On $\chXIVA$,
$\sigma^2=1$ (Chapter~\ref{ch:11}); on $\chXIVB$, $\chXIVsig$ goes to $T_s^+$ by $\Phi$ and back by $\Psi$,
and $\Psi\Phi=\mathrm{id}$; on $T_s^+$, $\Phi\Psi=\mathrm{id}$; on $T_{s+1}^+$, $\Theta^2=\mathrm{id}$.
(b) The modes of $u\in\chXIVA$ preserve each of the four sectors.  On $\chXIVA$ the identity is
$\sigma(u_{(n)}a)=(\sigma u)_{(n)}\sigma a$, i.e.\ $\sigma\in\Aut\chXIVA$; on the other three it is the
twisted-intertwiner property of $\Phi$, $\Psi$, $\Theta$.  The conjugation form follows from (a).
(c) Take $u=\omega$: $\sigma\omega=\omega$, and $L_{-1}$ is the coefficient of $Y^\natural(\omega)$ that
gives $\omega_{(0)}$.  The vacuum lies in $\chXIVA$ and $\sigma\mathbb 1=\mathbb 1$. \qed

\begin{status}Unconditional given $D$; standard axioms only.\end{status}

\subsection{\texorpdfstring{Uniqueness using $\chXIVA$-locality only}{Uniqueness using V0+-locality only}}

The heart of the conjugation argument is a uniqueness lemma.  The library already had a
``Goddard uniqueness'' theorem from Chapter~\ref{ch:10}, which we now quote because it is the wrong
tool here, and it is instructive to see why.

\begin{leancode}
theorem YnatT_eq_of_creates (hspan : ModeSpan g) (x : TP) (Ψ : FormalDist (End ℂ Vnat))
    (hcov : DCovariant Tnat Ψ) (h0 : Ψ 0 vacNat = ((0 : SkewVP), x))
    (hneg : ∀ q : ℤ, q < 0 → Ψ q vacNat = 0) (hU : ∀ v : SkewVP, Local Ψ (YnatU v))
    (hT : ∀ j, Local Ψ (YnatT leechWbar (g j))) : YnatT leechWbar x = Ψ
\end{leancode}
(\leanpath{VOA/Moonshine/M8h0Gap.lean}.)  It says: a translation-covariant field $\Psi$ that creates
$(0,x)$ from the vacuum and is local with \emph{every} untwisted field $Y^\natural(v)$,
$v\in V_\Leech^+$, and with the twisted generators, is the block field $Y^\natural((0,x))$.

\begin{warning}[title={Trap / dead end: The circularity of \lean{YnatT_eq_of_creates} in this context}]
We want to identify $L_v=\chXIVsig Y^\natural(v)\chXIVsig$, for $v\in\chXIVB$, with
$Y^\natural(\chXIVsig v)$.  To apply \lean{YnatT_eq_of_creates} with $\Psi=L_v$ we would need $L_v$ to be
local with $Y^\natural(b)$ for $b\in\chXIVB$.  Conjugating by $\chXIVsig$, this is the statement that
$Y^\natural(v)$ is local with $\chXIVsig Y^\natural(b)\chXIVsig=L_b$ — which is (part of) the very
identity $L_b=Y^\natural(\chXIVsig b)$ to be proved, now for $b$.  Using the theorem would be circular.
The same circularity reappears if one tries to use $\chXIVB$-locality of the difference field
(recorded in \leanpath{Notes/SigmaKey.md} \S7).  The remedy is a uniqueness lemma that asks for
locality with $\chXIVA$-fields only, because $\chXIVsig$ transports those for free
(Theorem~\ref{ch14:thm:signat}(b)).
\end{warning}

\begin{theorem}[Circularity-free uniqueness]\label{ch14:thm:V0local}
Let $F$ be a formal distribution with coefficients in $\End\Vnat$ that is local with
$Y^\natural(u)$ for every $u\in\chXIVA$.  Fix $\beta\in\Leech\setminus\Leech_0$ and nonzero
$t_0\in T_0$, $t_1\in T_1$.  If every coefficient of $F$ kills the four vectors
\begin{equation*}
  \mathbb 1,\qquad (x_\beta,0),\qquad (0,x_0(t_0)),\qquad (0,x_0(t_1)),
\end{equation*}
then $F=0$.
\end{theorem}
\begin{leancode}
theorem eq_zero_of_V0local (F : FormalDist (End ℂ Vnat))
    (hloc : ∀ u : V0Sub.carrier, Local F (YnatU (u : VPlus)))
    {β : leech.carrier} (hβ : β ∉ lambda0) (t : ZMod 2 → TLeechT)
    (ht : ∀ s, t s ∈ leechFrameT s) (ht0 : ∀ s, t s ≠ 0)
    (hvac : ∀ p, F p vacNat = 0) (h3 : ∀ p, F p (frameX0P β, 0) = 0)
    (hT : ∀ s p, F p (0, x0T (t s)) = 0) : F = 0
\end{leancode}
(\leanpath{VOA/Moonshine/M8hKeySigma.lean}, \lean{VOA.Moonshine.M8hKey.eq_zero_of_V0local}.)

\begin{dictionary}
\textbf{Decoding.}  \lean{leech.carrier} is the Leech lattice $\Leech$, \lean{lambda0} its
frame sublattice $\Leech_0$ (Chapter~\ref{ch:11}); \lean{leechFrameT s} is the frame eigenspace
$T_s\subseteq T$; the two vectors $t_0,t_1$ are packaged as a function \lean{t : ZMod 2 → TLeechT}.
\lean{F p v = 0} for all $p$ says every coefficient of $F$ kills $v$.

\textbf{Hypotheses.}  (i) $\chXIVA$-locality of $F$ — only $\chXIVA$, which is the point; (ii)
$\beta\notin\Leech_0$, so that $x_\beta$ lies in $\chXIVB$ and is $\chXIVA$-cyclic there (Chapter~\ref{ch:11},
\lean{cyc_V3}); (iii) $t_s\in T_s$ nonzero, so that $x_0(t_s)$ is $\chXIVA$-cyclic in $T_s^+$
(\lean{cyc_T}); (iv) the four annihilation hypotheses.  No translation covariance is needed.

\textbf{Same as textbook.}  This is the familiar ``a field local with a generating set and
annihilating the cyclic vectors is zero'', in its four-sector form.
\end{dictionary}
\begin{physics}
A field that is local with a chiral algebra and annihilates the lowest state of each superselection sector vanishes: the Reeh--Schlieder-type principle that states generated by local operators from a cyclic vector exhaust the sector.
\end{physics}


\paragraph{Proof on paper.}  Let $K=\bigcap_p\ker F_p$.  Locality plus annihilation give the
key fact: if $F$ is local with $Y^\natural(u)$ and all $F_p$ kill $y$, then all $F_p$ kill every
$u_{(q)}y$.  (Write $(z_1-z_2)^N[F(z_1),Y(u,z_2)]=0$ and apply to $y$; since $F(z_1)y=0$ one gets
$(z_1-z_2)^NF(z_1)Y(u,z_2)y=0$; the series $Y(u,z_2)y$ has only finitely many negative powers of
$z_2$, so the factor $(z_1-z_2)^N$ can be cancelled, giving $F(z_1)Y(u,z_2)y=0$.)  Thus $K$ is
$\chXIVA$-stable.  It contains $\mathbb 1$, hence $u_{(-1)}\mathbb 1=u$ for all $u\in\chXIVA$; it
contains the $\chXIVA$-cyclic vectors of the other three sectors, hence the three sectors.  So
$K=\Vnat$. \qed

\begin{status}Unconditional; standard axioms only.  Uses no form of (4.23).\end{status}

\subsection{\texorpdfstring{The identification of $L_v$ with $R_v$, and the key identities}{The identification of L_v with R_v, and the key identities}}

\begin{theorem}[Agreement on cyclic vectors suffices]\label{ch14:thm:conjV3}
Let $D$ be triality data, $v\in\chXIVB$, and $\beta,t$ as in Theorem~\ref{ch14:thm:V0local}.  If
$L_v=\chXIVsig Y^\natural(v)\chXIVsig$ and $R_v=Y^\natural((0,\Phi v))$ agree on $(x_\beta,0)$ and on
$(0,x_0(t_s))$ for both $s$, then $L_v=R_v$.
\end{theorem}
\begin{leancode}
theorem conjF_V3_eq {v : VPlus} (hv : v ∈ frameSector 1) {β : leech.carrier} (hβ : β ∉ lambda0)
    (t : ZMod 2 → TLeechT) (ht : ∀ s, t s ∈ leechFrameT s) (ht0 : ∀ s, t s ≠ 0)
    (h3 : ∀ p, conjF (sigNat D) (YnatU v) p (frameX0P β, 0) = YnatT leechWbar (D.Φ v) p (frameX0P β, 0))
    (hT : ∀ s p, conjF (sigNat D) (YnatU v) p (0, x0T (t s)) = YnatT leechWbar (D.Φ v) p (0, x0T (t s))) :
    conjF (sigNat D) (YnatU v) = YnatT leechWbar (D.Φ v)
\end{leancode}
(\leanpath{VOA/Moonshine/M8hKey.lean}.)

\begin{dictionary}
\textbf{Decoding.}  \lean{YnatT leechWbar x} is the purely twisted block field
$Y^\natural((0,x),z)=\begin{pmatrix}0&\overline W(x,z)\\W(x,z)&0\end{pmatrix}$ built with the genuine
adjoint \lean{leechWbar} (Chapter~\ref{ch:09}).  \lean{YnatU v} for $v\in\chXIVB$ is the untwisted block
field of $(v,0)$.  Note that $R_v=Y^\natural(\chXIVsig(v,0))$ since $\chXIVsig(v,0)=(0,\Phi v)$.

\textbf{Hypotheses.}  The two agreement hypotheses are for all coefficients $p$; there is no
hypothesis on the vacuum — that agreement is \emph{proved}.
\end{dictionary}
\begin{physics}
Two fields with the same OPEs with the even chiral algebra $\chXIVA$ and the same action on the
lowest state of every sector are the same field.  This is how a physicist identifies the transform
of an operator under a symmetry: compare its Ward identities and its matrix elements on primaries.
\end{physics}


\paragraph{Proof on paper.}  Apply Theorem~\ref{ch14:thm:V0local} to $F=L_v-R_v$.  Locality with
$\chXIVA$: $R_v$ is local with $Y^\natural(u)$ by the module/intertwiner relations (4.20)--(4.21)
proved in Chapter~\ref{ch:10}; $L_v$ is, because $Y^\natural(v)$ is local with $Y^\natural(\sigma u)$
(Chapter~\ref{ch:10}) and conjugation preserves locality and turns $Y^\natural(\sigma u)$ into
$Y^\natural(u)$ (Theorem~\ref{ch14:thm:signat}).  The vacuum: both fields are $L_{-1}$-covariant
(conjugation preserves covariance since $\chXIVsig L_{-1}=L_{-1}\chXIVsig$), both have vanishing
coefficients of $z^{q}$, $q<0$, on $\mathbb 1$, and both create $(0,\Phi v)$ at $z^0$ ($L_v$ because $\chXIVsig$ fixes $\mathbb 1$ and maps $(v,0)$ to $(0,\Phi v)$).
A covariant field is determined on $\mathbb 1$ by these data: $F(z)\mathbb 1=e^{zL_{-1}}F_0\mathbb 1$.
So $F$ kills $\mathbb 1$, and the remaining three vectors are the hypotheses. \qed

\begin{definition}[The FLM three-point identities on cyclic vectors]\label{ch14:def:sigmakey}
Let $D$ be triality data, $\beta\in\Leech\setminus\Leech_0$, and nonzero $t_s\in T_s$.  Write
$x=x_0(t_s)\in T_s^+$ (with $s=D.s$), $v=\Psi x\in\chXIVB$, $x_3=x_\beta$, $x'=x_0(t_{s+1})$.
\lean{SigmaKey D β t} is the conjunction of
\begin{align*}
 &\text{(key\_V3)} & \Theta\bigl(Y^T(v,z)\,\Phi x_3\bigr)&=W(x,z)\,x_3 &&\text{in }T_{s+1}^+,\\
 &\text{(key\_Ts)} & \sigma\bigl(Y(v,z)\,v\bigr)&=\overline W(x,z)\,x &&\text{in }\chXIVA,\\
 &\text{(key\_Ts1)} & \Psi\bigl(Y^T(v,z)\,\Theta x'\bigr)&=\overline W(x,z)\,x' &&\text{in }\chXIVB,
\end{align*}
coefficientwise, together with the side conditions on $\beta$ and $t$.
\end{definition}
\begin{leancode}
structure SigmaKey (β : leech.carrier) (t : ZMod 2 → TLeechT) : Prop where
  hβ : β ∉ lambda0
  ht : ∀ s, t s ∈ leechFrameT s
  ht0 : ∀ s, t s ≠ 0
  key_V3 : ∀ p, D.Θ (plusModule.YM (D.Ψ (x0T (t D.s))) p (D.Φ (frameX0P β)))
    = leechWsk (x0T (t D.s)) p (frameX0P β)
  key_Ts : ∀ p, sig0 (leechVOAPlus.Y (D.Ψ (x0T (t D.s))) p (D.Ψ (x0T (t D.s))))
    = leechWbar (x0T (t D.s)) p (x0T (t D.s))
  key_Ts1 : ∀ p, D.Ψ (plusModule.YM (D.Ψ (x0T (t D.s))) p (D.Θ (x0T (t (D.s + 1)))))
    = leechWbar (x0T (t D.s)) p (x0T (t (D.s + 1)))
\end{leancode}
(\leanpath{VOA/Moonshine/M8hKey.lean}, \lean{VOA.Moonshine.M8hKey.SigmaKey}.)

\begin{dictionary}
\textbf{Decoding.}  \lean{leechWsk x p φ} is the coefficient of $z^p$ in $W(x,z)\varphi$;
\lean{leechWbar x p y} that of $\overline W(x,z)y$.  \lean{sig0} is $\sigma$ on $\chXIVA$ (and $0$ on $\chXIVB$);
the product $Y(v,z)v$ of two elements of $\chXIVB$ lies in $\chXIVA$ by the fusion rules of the frame
grading, so \lean{sig0} acts as $\sigma$ there.

\textbf{Why these three, and why they are the agreements.}  The Lean theorems
\begin{leancode}
theorem key_V3_iff {v a : VPlus} (hv : v ∈ frameSector 1) (ha : a ∈ frameSector 1) (p : ℤ) :
    conjF (sigNat D) (YnatU v) p (a, 0) = YnatT leechWbar (D.Φ v) p (a, 0)
      ↔ D.Θ (plusModule.YM v p (D.Φ a)) = leechWsk (D.Φ v) p a
theorem key_Ts_iff {v : VPlus} (hv : v ∈ frameSector 1) {y : TP} (hy : y ∈ TPframe D.s) (p : ℤ) :
    conjF (sigNat D) (YnatU v) p (0, y) = YnatT leechWbar (D.Φ v) p (0, y)
      ↔ sig0 (leechVOAPlus.Y v p (D.Ψ y)) = leechWbar (D.Φ v) p y
theorem key_Ts1_iff {v : VPlus} (hv : v ∈ frameSector 1) {y : TP} (hy : y ∈ TPframe (D.s + 1)) (p : ℤ) :
    conjF (sigNat D) (YnatU v) p (0, y) = YnatT leechWbar (D.Φ v) p (0, y)
      ↔ D.Ψ (plusModule.YM v p (D.Θ y)) = leechWbar (D.Φ v) p y
\end{leancode}
show that, at $v=\Psi x$ (so $\Phi v=x$) and on the cyclic vectors, the three clauses are
\emph{exactly} the three agreement hypotheses of Theorem~\ref{ch14:thm:conjV3}.  To see it on
paper: $\chXIVsig$ sends $(a,0)$, $a\in\chXIVB$, to $(0,\Phi a)$; $Y^\natural(v)$ maps it to $(0,Y^T(v)\Phi a)$
with $Y^T(v)\Phi a\in T_{s+1}^+$; $\chXIVsig$ applies $\Theta$.  The block field $R_v$ sends $(a,0)$ to
$(0,W(\Phi v)a)$.  The other two rows are similar.

\textbf{Hypothesis status.}  \lean{SigmaKey} is a \emph{genuine open input} at this stage: each clause
compares two different intertwining operators and does not follow from the intertwining property of
$\Phi,\Psi,\Theta$ alone.  It also pins the normalisations that \lean{SigmaData} leaves free.
\end{dictionary}
\begin{physics}
Three specific three-point functions (OPE coefficients) must agree in the original and the triality-rotated pictures. If $\chXIVsig$ is a symmetry they do; these are FLM's ``key identities''.
\end{physics}


\begin{theorem}[(4.23) from the key identities]\label{ch14:thm:tpblock}
If \lean{SigmaKey D β t} holds, then (i) $\chXIVsig Y^\natural(\Psi x)\chXIVsig=Y^\natural((0,x))$ for
$x=x_0(t_s)$; (ii) (4.23) holds for \emph{all} pairs $t_1,t_2\in T^+$; (iii) $\Vnat$ carries all the
structure listed in the final theorem.
\end{theorem}
\begin{leancode}
theorem SigmaKey.conj_eq {β : leech.carrier} {t : ZMod 2 → TLeechT} (h : SigmaKey D β t) :
    conjF (sigNat D) (YnatU (D.Ψ (x0T (t D.s)))) = YnatT leechWbar (x0T (t D.s))

theorem tpBlockLocal_of_sigma {β : leech.carrier} {t : ZMod 2 → TLeechT} (h : SigmaKey D β t)
    (t₁ t₂ : TP) : TPBlockLocal leechWbar t₁ t₂

theorem moonshine_of_sigma {β : leech.carrier} {t : ZMod 2 → TLeechT} (h : SigmaKey D β t) :
    ∃ V : VertexOperatorAlgebra Vnat,
      V.Y = Ynat VnatAxioms.wbarL ∧ V.vacuum = vacNat ∧ V.conformalVector = omegaNat ∧
      V.centralCharge = 24 ∧
      V.grading 1 = ⊥ ∧
      Module.finrank ℂ (V.grading 2) = 196884 ∧
      (∀ n : ℕ, (Module.finrank ℂ (V.grading (n : ℤ)) : ℂ) = PowerSeries.coeff n JShift) ∧
      V.IsSimple
\end{leancode}
(\leanpath{VOA/Moonshine/M8hKey.lean}.)

\begin{dictionary}
\textbf{Decoding.}  \lean{TPBlockLocal leechWbar t₁ t₂} is (4.23):
\noindent File \leanpath{VOA/Moonshine/TTBlock.lean}.
\begin{leancode}
def TPBlockLocal (Wbar : TP → FormalDist (TP →ₗ[ℂ] SkewVP)) (t₁ t₂ : TP) : Prop :=
  ∃ N : ℕ, mulBySubPow N
    ((fun m n => leechWsk t₁ m ∘ₗ Wbar t₂ n - leechWsk t₂ n ∘ₗ Wbar t₁ m) :
      FormalDist₂ (End ℂ TP)) = 0
\end{leancode}
i.e.\ $(z_1-z_2)^N\bigl(W(t_1,z_1)\overline W(t_2,z_2)-W(t_2,z_2)\overline W(t_1,z_1)\bigr)=0$, with
\lean{mulBySubPow N} multiplication by $(z_1-z_2)^N$ on two-variable formal series.  The conclusion of
\lean{moonshine_of_sigma} is literally the conclusion of the final theorem, decoded in
\S\ref{ch14:sec:decode}.

\textbf{Hypotheses.}  Only $D$ (implicit section variable) and $h$.  Since \lean{SigmaKey} contains the side
conditions on $\beta$ and $t$, nothing else is assumed.
\end{dictionary}
\begin{physics}
If the candidate symmetry maps one untwisted field to the corresponding twist field, mutual locality of twist fields (single-valuedness of their OPEs, the delicate point of any orbifold) is inherited from the untwisted sector. Once one twist field is local with itself, Dong's lemma propagates locality to all.
\end{physics}


\paragraph{Proof on paper.}  (i) By Theorem~\ref{ch14:thm:conjV3} with $v=\Psi x$, whose hypotheses are
the clauses of \lean{SigmaKey} (by the three \lean{iff} theorems), and $\Phi\Psi x=x$.  (ii) The untwisted
field $Y^\natural(\Psi x)$ is local with itself (Chapter~\ref{ch:10}); conjugating, by (i),
$Y^\natural((0,x))$ is local with itself.  Self-locality of a twisted block field is the conjunction of
(4.22) and (4.23) for the pair $(x,x)$ (\lean{YnatT_local_iff}), so (4.23) holds for $(x,x)$.  By
Theorem~\ref{ch14:thm:modespan}, $x$ mode-spans $T^+$, and the propagation theorem of
Chapter~\ref{ch:10} (\lean{TPBlockLocal_all_of_gen}: (4.23) on a mode-spanning family implies (4.23)
for all pairs; its proof is Dong's lemma plus the Goddard uniqueness theorem, now used legitimately,
because the twisted generator's locality is supplied) gives (4.23) for all pairs.  (iii) Feed (ii)
into \lean{moonshine_of_frameLocality} (Chapter~\ref{ch:10}). \qed

\begin{status}
Conditional on \lean{SigmaKey D β t} (and on the data $D$).  \texttt{\#print axioms} \lean{VOA.Moonshine.M8hKey.moonshine_of_sigma}, re-run for this chapter with
\lean{lake env lean -DmaxSynthPendingDepth=3}: \lean{[propext, Classical.choice, Quot.sound]}.
No step assumes (4.23) for any pair.
\end{status}

\section{Reducing the key identities: isometries and one self-coupling}\label{ch14:sec:sigmakey}

The three clauses of \lean{SigmaKey} are not independent.  The difference field
$F_v=L_v-R_v$ (for $v\in\chXIVB$) has four off-diagonal blocks between the four sectors, and the
invariant bilinear forms pair them up:
\begin{center}
\begin{tabular}{llll}
\toprule
block of $F_v$ & content & how it is handled \\ \midrule
$A\to C$ & creation & automatic (vacuum + $\chXIVA$-locality)\\
$C\to A$ & \lean{key_Ts} & adjoint of $A\to C$, given that $\Phi$ is an isometry\\
$B\to D'$ & \lean{key_V3} & the genuine content: \lean{KeySkew}\\
$D'\to B$ & \lean{key_Ts1} & adjoint of $B\to D'$, given that $\Phi,\Theta$ are isometries\\
\bottomrule
\end{tabular}
\end{center}
This section makes this table precise.

\subsection{Automorphisms preserve invariant forms}

Recall from Definition~\ref{l04:def:invform} that an \emph{invariant bilinear form} on a VOA $W$ is a bilinear
form $(\cdot|\cdot)$ with $(L_na|b)=(a|L_{-n}b)$, and that it is \emph{fully invariant} if for every
homogeneous $u$ of weight $k$
\begin{equation}\label{ch14:eq:fhl}
  (u_{(n)}a\,|\,b)=(-1)^k\sum_{j\ge0}\frac{1}{j!}\,\bigl(a\,\big|\,(L_1^ju)_{(2k-n-2-j)}b\bigr)
\end{equation}
(FHL's invariance, the adjoint of $Y(u,z)$ being $Y(e^{zL_1}(-z^{-2})^{L_0}u,z^{-1})$).  It is
\emph{normalised} if $(\mathbb 1|\mathbb 1)=1$.  For a VOA of CFT type with $L_1V_1=0$ there is
exactly one normalised fully invariant form (\lean{invariantForm_unique_FHL}, Theorem~\ref{l04:thm:canonical}).

\begin{theorem}[Automorphisms are isometries]\label{ch14:thm:autiso}
Let $W$ be a VOA with $\mathbb 1\ne0$ and $L_1W_1=0$, $B$ a fully invariant normalised form, and
$\alpha\in\Aut W$.  Then $B(\alpha a,\alpha b)=B(a,b)$ for all $a,b$.
\end{theorem}
\begin{leancode}
theorem aut_isometry {W : VertexOperatorAlgebra 𝒱} (hne : W.vacuum ≠ 0)
    (hqp : ∀ v, v ∈ W.grading 1 → W.L 1 v = 0) {B : InvariantForm W}
    (hB : B.IsFullyInvariant) (hn : B.Normalized) (α : Aut W) (a b : 𝒱) :
    B.form (α a) (α b) = B.form a b
\end{leancode}
(\leanpath{VOA/Moonshine/SigmaKeyForm.lean}, \lean{VOA.Moonshine.SigmaKey.aut_isometry}.)

\begin{dictionary}
\textbf{Decoding.}  \lean{InvariantForm W} bundles a bilinear form \lean{B.form} with the Virasoro
adjointness; \lean{IsFullyInvariant} is \eqref{ch14:eq:fhl}; \lean{Normalized} is
$(\mathbb 1|\mathbb 1)=1$.  \lean{Aut W} is the group of VOA automorphisms (linear bijections
fixing $\mathbb 1$, $\omega$ and intertwining all modes).  Hypotheses \lean{hne}, \lean{hqp} are exactly
those of the uniqueness theorem; \lean{hqp} is ``$L_1$ kills $W_1$''.  Note that CFT type is part
of the library's definition of a VOA (Definition~\ref{l03:def:VOA}: $W_0=\C\mathbb 1$, $W_n=0$ for $n<0$).
Same statement as the textbook one.
\end{dictionary}
\begin{physics}
A symmetry preserves the two-point function: $\langle g a|g b\rangle=\langle a|b\rangle$. Here this follows from uniqueness of the vacuum-normalised two-point function in a CFT with no currents annihilated by $L_1$.
\end{physics}


\paragraph{Proof on paper.}  The pull-back $(a,b)\mapsto B(\alpha a,\alpha b)$ is again fully
invariant (since $\alpha$ commutes with every $L_n$ and every mode) and normalised (since
$\alpha\mathbb 1=\mathbb 1$); by uniqueness it equals $B$. \qed

\begin{corollary}\label{ch14:cor:orth}
(a) For the canonical form $B_V$ on $V_\Leech^+$: $B_V(\chXIVA,\chXIVB)=0$.
(b) $\sigma$ is an isometry of $B_V$ restricted to $\chXIVA$.
\end{corollary}
\begin{leancode}
theorem leechBV_frame_orth {a b : VPlus} (ha : a ∈ frameSector 0) (hb : b ∈ frameSector 1) :
    leechBV a b = 0

theorem leechBV_sig (a b : V0Sub.carrier) :
    leechBV ((sig a : V0Sub.carrier) : VPlus) (sig b : V0Sub.carrier) = leechBV (a : VPlus) b
\end{leancode}
(\leanpath{VOA/Moonshine/SigmaKeyForm.lean}.)

\begin{dictionary}
\lean{leechBV} is the canonical (Li/FHL) normalised invariant form of $V_\Leech^+$ (Chapter~\ref{ch:09},
\lean{leechBVForm}).  The frame sectors are the $\pm1$ eigenspaces of the frame automorphism
$\pi$ (\lean{plusFrameAut}), which acts by $+1$ on $\chXIVA$ and $-1$ on $\chXIVB$.  (b) is stated for
elements of $\chXIVA$, coerced into $V_\Leech^+$.  Both exactly the intended statements.
\end{dictionary}
\begin{physics}
States in different eigenspaces of a $\Z/2$ symmetry have vanishing two-point functions (charge conservation); and the triality restricted to $\chXIVA$ is unitary for its two-point function.
\end{physics}


\paragraph{Proof on paper.}  (a) Apply Theorem~\ref{ch14:thm:autiso} to $\pi$:
$B_V(a,b)=B_V(\pi a,\pi b)=-B_V(a,b)$.  (b) Apply it to the VOA $\chXIVA$, whose restricted form is
fully invariant and normalised, and to $\sigma$.  ($(\chXIVA)_1=0$, so $L_1$ kills it trivially.) \qed

\begin{status}Unconditional; standard axioms only.\end{status}

\subsection{The three named conditions}

\begin{definition}[Isometry conditions and the self-coupling identity]\label{ch14:def:iso}
For triality data $D$ and $t$ as before, with $x=x_0(t_s)$:
\begin{itemize}
\item \lean{PhiIso D}: $\langle\Phi a,\Phi b\rangle_T=B_V(a,b)$ for all $a,b\in\chXIVB$;
\item \lean{ThetaIso D}: $\langle\Theta y,\Theta y'\rangle_T=\langle y,y'\rangle_T$ for all $y,y'\in T_{s+1}^+$;
\item \lean{KeySkew D t}: $\Theta\bigl(Y^T(\Psi x,z)x\bigr)=W(x,z)\Psi x$, coefficientwise;
\end{itemize}
here $\langle\cdot,\cdot\rangle_T$ is \lean{TPForm}.
\end{definition}
\begin{leancode}
def PhiIso : Prop :=
  ∀ a ∈ frameSector 1, ∀ b ∈ frameSector 1, TPForm (D.Φ a) (D.Φ b) = leechBV a b

def ThetaIso : Prop :=
  ∀ y ∈ TPframe (D.s + 1), ∀ y' ∈ TPframe (D.s + 1), TPForm (D.Θ y) (D.Θ y') = TPForm y y'

def KeySkew (t : ZMod 2 → TLeechT) : Prop :=
  ∀ p, D.Θ (plusModule.YM (D.Ψ (x0T (t D.s))) p (x0T (t D.s)))
    = leechWsk (x0T (t D.s)) p (D.Ψ (x0T (t D.s)))
\end{leancode}
(\leanpath{VOA/Moonshine/SigmaKeyAdj.lean}.)

\begin{dictionary}
\textbf{Decoding.}  \lean{TPForm} is the symmetric nondegenerate invariant form on $T^+$ of
Chapter~\ref{ch:09}.  In \lean{KeySkew}, $v=\Psi x$, the left side is $\Theta(Y^T(v,z)x)$ and the right
side $W(x,z)v$; by the definition of $W$ this is $e^{zL_{-1}}Y^T(v,-z)x$.  So \lean{KeySkew} says: the
series $J(z)=Y^T(v,z)x$ is mapped by $\Theta$ to its ``skew-symmetric partner''
$e^{zL_{-1}}J(-z)$.  It is a statement about one vector-valued series.

\textbf{Status as hypotheses.}  These are \emph{named Props}, i.e.\ statements to be proved; they are
not assumed anywhere in the final theorem.  Note \lean{PhiIso} is not a normalisation that can always be
arranged: it requires $\Phi$ to be a similitude in the first place.
\end{dictionary}
\begin{physics}
\lean{PhiIso}, \lean{ThetaIso}: the candidate symmetry preserves two-point functions between sectors. \lean{KeySkew}: the OPE of one odd untwisted operator with one twist field is symmetric under exchange after applying the symmetry — a single crossing-type relation.
\end{physics}


\subsection{The creation block and the adjunction chains}

\begin{proposition}[Creation block]\label{ch14:prop:create}
For triality data $D$, $v\in\chXIVB$, $a\in\chXIVA$ and $m\in\Z$:
$\Phi\bigl(v_{(m)}\sigma a\bigr)=W(\Phi v)_{(m)}a$.
\end{proposition}
\begin{leancode}
theorem phi_create {v : VPlus} (hv : v ∈ frameSector 1) (a : V0Sub.carrier) (m : ℤ) :
    D.Φ (leechVOAPlus.smul v m ((sig a : V0Sub.carrier) : VPlus))
      = mode (leechWsk (D.Φ v)) m (a : VPlus)
\end{leancode}
(\leanpath{VOA/Moonshine/SigmaKeyAdj.lean}.)

\begin{dictionary}
\lean{leechVOAPlus.smul v m w} is the $m$-th product $v_{(m)}w$ in $V_\Leech^+$ (Definition~\ref{l03:def:smulT}:
\lean{smul} is the mode action, already in mode indexing).  \lean{mode (leechWsk x) m} is
$W(x)_{(m)}$.  Unconditional given $D$; this is \emph{not} an assumption.
\end{dictionary}
\begin{physics}
The state-operator map intertwines the rotated and unrotated pictures on the even sector: creating a state with the rotated odd operator equals creating it with the twist field. This is automatic because both are fixed by their action on the vacuum.
\end{physics}


\paragraph{Proof on paper.}  The difference field $F_v$ is local with all $\chXIVA$-fields and
kills the vacuum (proof of Theorem~\ref{ch14:thm:conjV3}); by the locality argument in the proof of
Theorem~\ref{ch14:thm:V0local} it kills $u_{(-1)}\mathbb 1=u$ for every $u\in\chXIVA$.  Now compute
both fields on $(a,0)$: $\chXIVsig(a,0)=(\sigma a,0)$; the untwisted field of $v$ maps it to
$(v_{(m)}\sigma a,0)$ with $v_{(m)}\sigma a\in\chXIVB$ (frame grading: $B\cdot A\subseteq B$); the second
$\chXIVsig$ maps this to $(0,\Phi(v_{(m)}\sigma a))$.  On the other side $R_v(a,0)=(0,W(\Phi v)_{(m)}a)$.
Equality of the second components is the claim. \qed

\begin{theorem}[\lean{key\_Ts} and \lean{key\_Ts1} by adjunction]\label{ch14:thm:keyTs}
Let $D$ be triality data, $x\in T_s^+$ of weight $2$, $v=\Psi x$.
\begin{enumerate}[label=(\alph*)]
\item If \lean{PhiIso D}, then $\sigma\bigl(v_{(n)}\Psi y\bigr)=\overline W(x)_{(n)}y$ for all
  $y\in T_s^+$, $n\in\Z$.
\item If \lean{PhiIso D}, \lean{ThetaIso D}, and the whole block $B\to D'$ vanishes, i.e.\
  $\Theta\bigl(v_{(m)}\Phi b\bigr)=W(x)_{(m)}b$ for all $b\in\chXIVB$, $m$, then
  $\Psi\bigl(v_{(n)}\Theta y\bigr)=\overline W(x)_{(n)}y$ for all $y\in T_{s+1}^+$.
\end{enumerate}
\end{theorem}
\begin{leancode}
theorem keyTs_gen (hΦ : PhiIso D) {x : TP} (hx : x ∈ TPframe D.s)
    (hx2 : x ∈ plusModule.weight ((2 : ℤ) : ℂ)) {y : TP} (hy : y ∈ TPframe D.s) (n : ℤ) :
    sig0 (leechVOAPlus.smul (D.Ψ x) n (D.Ψ y)) = leechWbarMode n x y

theorem keyTs1_gen (hΦ : PhiIso D) (hΘ : ThetaIso D) {x : TP} (hx : x ∈ TPframe D.s)
    (hx2 : x ∈ plusModule.weight ((2 : ℤ) : ℂ))
    (hV3 : ∀ b ∈ frameSector 1, ∀ m : ℤ,
      D.Θ (mode (plusModule.YM (D.Ψ x)) m (D.Φ b)) = mode (leechWsk x) m b)
    {y : TP} (hy : y ∈ TPframe (D.s + 1)) (n : ℤ) :
    D.Ψ (mode (plusModule.YM (D.Ψ x)) n (D.Θ y)) = leechWbarMode n x y
\end{leancode}
(\leanpath{VOA/Moonshine/SigmaKeyAdj.lean}.)

\begin{dictionary}
\textbf{Decoding.}  \lean{plusModule.weight w} is the $L_0$-eigenspace of eigenvalue $w$ in $T^+$
(weights of an ordinary module are complex numbers in the library, hence the cast
\lean{((2 : ℤ) : ℂ)}).  \lean{leechWbarMode n x} is $\overline W(x)_{(n)}$.  The hypothesis \lean{hV3}
in (b) is exactly ``the $B\to D'$ block vanishes for all $b$''.  Weight $2$ for $x$ is needed: it makes
$x$ and $v$ quasi-primary, which collapses the FHL sum \eqref{ch14:eq:fhl} to one term.
\end{dictionary}
\begin{physics}
Three-point functions $\langle a|\phi_1(z)|b\rangle$ can be read with any leg as the outgoing state; if the symmetry preserves two-point functions, OPE coefficients in one channel determine those in the adjoint channel.
\end{physics}


\paragraph{Proof on paper.}  (a) Both sides lie in $\chXIVA$, and $B_V$ is nondegenerate on $V_\Leech^+$
with $\chXIVA\perp\chXIVB$; so it suffices to pair with $a\in\chXIVA$:
\begin{align*}
B_V\bigl(\overline W(x)_{(n)}y,\,a\bigr)
 &=\langle y,\,W(x)_{(2-n)}a\rangle_T &&\text{(definition of }\overline W\text{, }L_1x=0)\\
 &=\langle y,\,\Phi(v_{(2-n)}\sigma a)\rangle_T &&\text{(creation block)}\\
 &=B_V\bigl(\Psi y,\,v_{(2-n)}\sigma a\bigr) &&\text{(\lean{PhiIso}, }y=\Phi\Psi y)\\
 &=B_V\bigl(v_{(n)}\Psi y,\,\sigma a\bigr) &&\text{(invariance, }v\text{ quasi-primary of weight 2)}\\
 &=B_V\bigl(\sigma(v_{(n)}\Psi y),\,a\bigr) &&(\sigma\text{ isometry and involution of }\chXIVA).
\end{align*}
(b) Similarly, pairing with $b\in\chXIVB$ and using \lean{ThetaIso}, invariance of \lean{TPForm} and
\lean{PhiIso}:
$B_V(\overline W(x)_{(n)}y,b)=\langle y,\Theta(v_{(2-n)}\Phi b)\rangle_T=\langle\Theta y,v_{(2-n)}\Phi b\rangle_T
=\langle v_{(n)}\Theta y,\Phi b\rangle_T=B_V(\Psi(v_{(n)}\Theta y),b)$; pairing with $\chXIVA$ gives zero on both
sides by sector orthogonality. \qed

\begin{theorem}[\lean{SigmaKey} from the isometries and \lean{KeySkew}]\label{ch14:thm:skew}
If \lean{PhiIso D}, \lean{ThetaIso D} and \lean{KeySkew D t} hold (with $\beta\notin\Leech_0$, $t_s\in T_s$
nonzero), then \lean{SigmaKey D β t}; consequently all targets of the final theorem hold.
\end{theorem}
\begin{leancode}
theorem sigmaKey_of_skew {β : leech.carrier} (hβ : β ∉ lambda0) {t : ZMod 2 → TLeechT}
    (ht : ∀ s, t s ∈ leechFrameT s) (ht0 : ∀ s, t s ≠ 0)
    (hΦ : PhiIso D) (hΘ : ThetaIso D) (hk : KeySkew D t) : SigmaKey D β t
\end{leancode}
(\leanpath{VOA/Moonshine/SigmaKeyAdj.lean}; the corollary is \lean{moonshine_of_skew}.)

\begin{dictionary}
All hypotheses are visible: the side conditions on $\beta,t$ and the three named Props.  The conclusion
is the structure of Definition~\ref{ch14:def:sigmakey}.
\end{dictionary}
\begin{physics}
All three key identities follow from unitarity of the symmetry on two sectors and one self-coupling OPE; this is the reduction a physicist would make by using conjugation and crossing.
\end{physics}


\paragraph{Proof on paper.}  The key extra observation is that $v=\Psi x$ is itself a
$\chXIVA$-cyclic vector of $\chXIVB$: $\Psi$ is a bijection $T_s^+\to\chXIVB$ intertwining up to $\sigma$,
and $x$ is cyclic in $T_s^+$.  \lean{KeySkew} says precisely that $F_v$ kills $(v,0)$ (by
\lean{key_V3_iff} with $a=v$, $\Phi v=x$).  By the locality argument, $F_v$ then kills all of $\chXIVB$:
the block $B\to D'$ vanishes identically.  This gives \lean{key_V3} (at $x_\beta$), and through
Theorem~\ref{ch14:thm:keyTs} also \lean{key_Ts} and \lean{key_Ts1}. \qed

\subsection{The isometries are top-level statements}

\begin{theorem}[Top-level reduction of the isometry conditions]\label{ch14:thm:isotop}
Let $z_0\in\chXIVB$ be $\chXIVA$-cyclic and homogeneous of weight $w$.  If
$\langle\Phi z_0,\Phi b\rangle_T=B_V(z_0,b)$ for all $b\in(\chXIVB)_w$, then \lean{PhiIso D}.  Likewise,
if $x'\in T_{s+1}^+$ is cyclic of weight $w$ and $\langle\Theta x',\Theta y\rangle_T=\langle x',y\rangle_T$
for all $y\in(T_{s+1}^+)_w$, then \lean{ThetaIso D}.
\end{theorem}
\begin{leancode}
theorem phiIso_of_top {z₀ : VPlus} (hz₀ : z₀ ∈ frameSector 1) {w : ℤ}
    (hzw : z₀ ∈ leechVOAPlus.grading w)
    (hcyc : Submodule.span ℂ {y | ∃ u ∈ (Set.univ : Set V0Sub.carrier), ∃ n : ℤ, y = actV u n z₀}
      = frameSector 1)
    (htop : ∀ b ∈ frameSector 1, b ∈ leechVOAPlus.grading w →
      TPForm (D.Φ z₀) (D.Φ b) = leechBV z₀ b) : PhiIso D

theorem thetaIso_of_top {x' : TP} (hx' : x' ∈ TPframe (D.s + 1)) {w : ℂ}
    (hxw : x' ∈ plusModule.weight w)
    (hcyc : Submodule.span ℂ {y | ∃ u ∈ (Set.univ : Set V0Sub.carrier), ∃ n : ℤ, y = actT u n x'}
      = TPframe (D.s + 1))
    (htop : ∀ y ∈ TPframe (D.s + 1), y ∈ plusModule.weight w →
      TPForm (D.Θ x') (D.Θ y) = TPForm x' y) : ThetaIso D
\end{leancode}
(\leanpath{VOA/Moonshine/SigmaKeyTop.lean}.)

\begin{dictionary}
\textbf{Decoding.}  \lean{hcyc} says $z_0$ is $\chXIVA$-cyclic: the modes of all $u\in\chXIVA$ applied to
$z_0$ span $\chXIVB$.  \lean{leechVOAPlus.grading w} is the weight-$w$ subspace.  \lean{htop} is one ``row''
of the Gram matrices, on a single finite-dimensional weight space.  Same as the mathematical statement.
\end{dictionary}
\begin{physics}
A two-point function of a sector generated by its lowest state is determined by the two-point function of that lowest state with other lowest states (Ward identities move all modes to one side).
\end{physics}


\paragraph{Proof on paper.}  The defect $F(a,b)=\langle\Phi a,\Phi b\rangle_T-B_V(a,b)$ is a symmetric
bilinear form on $\chXIVB$ satisfying the same FHL invariance \eqref{ch14:eq:fhl} for $u\in\chXIVA$ (both
forms do; for the first one transports $u$ through $\Phi$ as $\sigma u$, and $\sigma$ commutes with
$L_1$).  Invariance moves every mode from the left slot to the right, so by cyclicity $F=0$ as soon as
$F(z_0,\cdot)=0$; by symmetry and cyclicity again it suffices that $F(z_0,c)=0$ for homogeneous $c$;
both forms vanish between different weights; so only $c$ of weight $w$ matter. \qed

\begin{definition}[The finite isometry residuals]\label{ch14:def:phitop}
\lean{PhiTop D t}: $\langle x,\Phi b\rangle_T=B_V(\Psi x,b)$ for all $b\in(\chXIVB)_2$, with $x=x_0(t_s)$.
\lean{ThetaTop D t}: $\langle\Theta x',\Theta y\rangle_T=\langle x',y\rangle_T$ for all $y\in(T_{s+1}^+)_2$, with
$x'=x_0(t_{s+1})$.
\end{definition}
\begin{leancode}
def PhiTop (t : ZMod 2 → TLeechT) : Prop :=
  ∀ b ∈ frameSector 1, b ∈ leechVOAPlus.grading 2 →
    TPForm (x0T (t D.s)) (D.Φ b) = leechBV (D.Ψ (x0T (t D.s))) b

def ThetaTop (t : ZMod 2 → TLeechT) : Prop :=
  ∀ y ∈ TPframe (D.s + 1), y ∈ plusModule.weight (2 : ℂ) →
    TPForm (D.Θ (x0T (t (D.s + 1)))) (D.Θ y) = TPForm (x0T (t (D.s + 1))) y
\end{leancode}
\begin{physics}
Normalisations of the lowest states: the norms of the $49152$ lowest odd untwisted states and twisted states must match under the symmetry.
\end{physics}


\begin{theorem}[$\Vnat$ from two finite isometries and \lean{KeySkew}]\label{ch14:thm:top}
\begin{leancode}
theorem sigmaKey_of_top {β : leech.carrier} (hβ : β ∉ lambda0) {t : ZMod 2 → TLeechT}
    (ht : ∀ s, t s ∈ leechFrameT s) (ht0 : ∀ s, t s ≠ 0)
    (hΦ : PhiTop D t) (hΘ : ThetaTop D t) (hk : KeySkew D t) : SigmaKey D β t

theorem moonshine_of_top {β : leech.carrier} (hβ : β ∉ lambda0) {t : ZMod 2 → TLeechT}
    (ht : ∀ s, t s ∈ leechFrameT s) (ht0 : ∀ s, t s ≠ 0)
    (hΦ : PhiTop D t) (hΘ : ThetaTop D t) (hk : KeySkew D t) :
    ∃ V : VertexOperatorAlgebra Vnat,
      V.Y = Ynat VnatAxioms.wbarL ∧ V.vacuum = vacNat ∧ V.conformalVector = omegaNat ∧
      V.centralCharge = 24 ∧
      V.grading 1 = ⊥ ∧
      Module.finrank ℂ (V.grading 2) = 196884 ∧
      (∀ n : ℕ, (Module.finrank ℂ (V.grading (n : ℤ)) : ℂ) = PowerSeries.coeff n JShift) ∧
      V.IsSimple
\end{leancode}
\end{theorem}
(\leanpath{VOA/Moonshine/SigmaKeyTop.lean}.)

\begin{dictionary}
In \lean{PhiTop} the cyclic vector is $z_0=\Psi x$, which has weight $2$ and satisfies $\Phi z_0=x$;
thus \lean{PhiTop} is the \lean{htop} of Theorem~\ref{ch14:thm:isotop} for this $z_0$.  \lean{PhiTop} and
\lean{ThetaTop} are \emph{finite}: they involve only the $49152$-dimensional weight-$2$ spaces.  The
hypotheses of \lean{moonshine_of_top} are listed in full; nothing else enters.
\end{dictionary}
\begin{physics}
The consistency of the orbifold now rests on finitely many normalisations and one OPE identity.
\end{physics}


\paragraph{Proof on paper.}  Theorem~\ref{ch14:thm:isotop} turns \lean{PhiTop}, \lean{ThetaTop} into
\lean{PhiIso}, \lean{ThetaIso}; then Theorem~\ref{ch14:thm:skew} and Theorem~\ref{ch14:thm:tpblock}. \qed

\begin{status}
Conditional on $D$, \lean{PhiTop D t}, \lean{ThetaTop D t}, \lean{KeySkew D t}; standard axioms only.
\end{status}

\subsection{Normalisation: what is forced and what is free}\label{ch14:sec:norm}

\lean{SigmaData} leaves free $\Phi\mapsto c\Phi$, $\Psi\mapsto c^{-1}\Psi$ ($c\ne0$) and $\Theta\mapsto-\Theta$.
How do the residuals react?
\begin{itemize}
\item \lean{PhiTop} scales as $c$ on the left and $c^{-1}$ on the right.  If $\Phi$ is a
  \emph{similitude}, $\langle\Phi a,\Phi b\rangle_T=\kappa_0B_V(a,b)$, then \lean{PhiTop} holds iff
  $c^2\kappa_0=1$: the scalar is determined up to sign.  The remaining sign is a genuine symmetry:
  $(\Phi,\Psi,\Theta)\mapsto(-\Phi,-\Psi,\Theta)$ is $\chXIVsig$ composed with the automorphism $-1$ on
  $\chXIVB\oplus T_s^+$, and changes none of the clauses.
\item \lean{ThetaTop} is insensitive to the sign of $\Theta$.  Since $\Theta^2=1$ and the invariant form on
  $T_{s+1}^+$ is unique up to a scalar, $\langle\Theta\cdot,\Theta\cdot\rangle_T=\mu\langle\cdot,\cdot\rangle_T$
  with $\mu=\pm1$.  \lean{ThetaTop} asserts $\mu=+1$, which is a \emph{condition}, not a normalisation.
  If $\mu=-1$, both eigenspaces of $\Theta$ would be isotropic and \lean{SigmaKey} unsatisfiable.
\item \lean{KeySkew} flips under $\Theta\mapsto-\Theta$; at most one sign can work.
\end{itemize}

\begin{warning}[title={Evidence: Numerical de-risking of the normalisation, and its limits}]
Before any of this was formalised, the project tested the conditions numerically in an independent
model of the weight-$2$ spaces built from the Niemeier lattice $N(A_1^{24})$ (route notes
\leanpath{Notes/SigmaKey.md} \S6, scripts in \leanpath{scripts/moonshine/numerics/sigmakey/}).  Findings,
\emph{evidence only}:
(N1) $\Theta_J^2\propto1$, and after normalisation $\mu=1.000000000000$ ($|\mu-1|<4\cdot10^{-12}$);
independently $\operatorname{tr}\Theta_J=-112\ne0$ excludes $\mu=-1$.
(N2) $\Phi$ is a similitude: the ratio $\langle\Phi v,\Phi v\rangle/B_V(v,v)$ is constant to spread
$6.7\cdot10^{-12}$, off-diagonal entries $<3.8\cdot10^{-12}$.
(N3) The top level of \lean{key_V3}, $\Theta\,o(v)\Phi b=\varepsilon\,o(b)\Phi v$, holds on $170$ pairs with one
global $\varepsilon$.  \emph{Caveat}: with $\Theta$ replaced by the identity the vectors were also
proportional (they all lay in one eigenspace), so this test fixes the sign but discriminates weakly.
The model's sector identifications and cocycle were not checked against the library's.
In the Lean development, $\mu=+1$ was eventually \emph{proved} (\S\ref{ch14:sec:assemble}), and the
sign was proved to be $\varepsilon=+1$ for the library's $\Theta_0$ (\S\ref{ch14:sec:toppairs}).
\end{warning}

\begin{warning}[title={Trap / dead end: Dead ends recorded at this stage}]
(1) ``An intertwiner is determined by its top component'' is false without the full $A(M)$-bimodule
theory: the top-level part of \lean{KeySkew} is necessary but \emph{not} sufficient; the iterate
formula moves descendants of the input into non-top outputs.  Do not formalise ``top-level
\lean{key_V3} $\Rightarrow$ \lean{key_V3}''.
(2) Using $\chXIVB$-locality of $L_v$ is circular (see the warning before
Theorem~\ref{ch14:thm:V0local}).
(3) (4.22) gives no new information: it constrains only products of already matched blocks or the
$B\to D'\to B$ composite, which is symmetric under the unknown scalar.
(4) Skew-symmetry alone does not decide \lean{KeySkew}: it says that
$J(v,z)b=\Theta Y^T(v,z)\Phi b$ is a fixed point of the involution $J\mapsto\Theta e^{zL_{-1}}J(\cdot,-z)$
(swapping $v,b$); without a one-dimensional space of intertwiners, $J$ could have components in both
eigenspaces.
\end{warning}

\section{The assembled triality data}\label{ch14:sec:assemble}

We now produce a concrete term of type \lean{SigmaData} and prove \lean{PhiTop} and \lean{ThetaTop} for
it.  The raw material is Chapter~\ref{ch:13}: a seed vector, explicit top-level maps, and their unique
twisted extensions.

\subsection{The seed and the explicit top maps}

Recall (Chapter~\ref{ch:13}) the operators on $T$: for $j\ne k$ there are lattice vectors
$\gamma^\pm_{jk}\in\Leech_0$ of norm $4$ (Lean \lean{gP j k}, \lean{gM j k}), with the corresponding
normalised operators $A^\pm_{jk}=\kappa^\pm_{jk}e_{\gamma^\pm_{jk}}$ on $T$ (\lean{AP}, \lean{AM}), and
$Z_{0k}=2A^-_{0k}$ (\lean{ZZ 0 k}).  These $Z_{0k}$, $k\ne0$, commute and square to $1$.

\begin{definition}[Seed]\label{ch14:def:seed}
A vector $t\in T$ is a \emph{seed} if $A^+_{jk}t=-\tfrac12t$ and $A^-_{jk}t=\tfrac12t$ for all $j\ne k$.
\end{definition}
\begin{leancode}
def IsSeed (t : TLeechT) : Prop :=
  ∀ j k : Fin 24, j ≠ k → AP j k t = -(1 / 2 : ℂ) • t ∧ AM j k t = (1 / 2 : ℂ) • t
\end{leancode}
\begin{physics}
A seed is a twisted ground state with prescribed eigenvalues under certain lattice-shift symmetries of the ground-state space: a ``highest-weight'' reference state for the finite Heisenberg group $2^{1+24}$ acting on the $2^{12}$ ground states.
\end{physics}

(\leanpath{VOA/Moonshine/SigmaTop2Eigen.lean}.)  Chapter~\ref{ch:13} proved that a nonzero seed exists,
and lies in a frame sector (\lean{seed_in_frame seedExists}).  We fix one: $t_0\in T_{s_0}$
(Lean \lean{t0}, \lean{s0}, chosen by \lean{Classical.choice}; the theorems below hold for every such
choice).

The two explicit weight-$2$ maps of Chapter~\ref{ch:13} are:
\begin{itemize}
\item $\varphi_\star:(\chXIVB)_2\to(T^+)_2$, $\varphi_\star(x_n)=c_\star(n)\,h_{w_{j(n)}}(-\tfrac12)\otimes e_nt_0$
  on the basis $x_n=e^n+e^{-n}$, $n\in S_1$ (a fixed half of the norm-$4$ vectors of
  $\Leech\setminus\Leech_0$), where $j(n)$ is the coordinate of the entry $\pm3$ of $n$, $w_j$ the $j$-th frame
  vector, and $c_\star(n)=i^{n_0}\cdot n_{j(n)}/3$, a fourth root of unity times $\pm1$; here
  $n_0,\dots,n_{23}$ are the integer coordinates of $n$ in the scaling where
  $\langle\mu,\nu\rangle=\tfrac18\sum_k\mu_k\nu_k$ (Lean \lean{lv}), in which a norm-$4$ vector of
  $\Leech\setminus\Leech_0$ has shape $(\mp3,(\pm1)^{23})$ (Lean \lean{phiV = phiV3 (tauStar t0)},
  \lean{cStar});
\item $\varphi_B:(T_{s_0}^+)_2\to(T_{s_0}^+)_2$, $\varphi_B(h_{w_m}(-\tfrac12)\otimes t)=h_{w_m}(-\tfrac12)\otimes\psi_mt$,
  where $\psi_0=B=\prod_{k\ne0}(a+bZ_{0k})$ (restricted to $T_{s_0}$), $\psi_m=Z_{0m}B$, with
  $a=\tfrac12(1+\omega)$, $b=\tfrac12(1-\omega)$, $\omega=e^{i\pi/4}$ (Lean \lean{phiB = phiTT (psiOf (Bop s0))}).
\end{itemize}
On a joint eigenvector $z$ of the $Z_{0k}$ with eigenvalues $u_k=\pm1$, $B$ acts by $\omega^{|N|}$ with
$N=\{k:u_k=-1\}$, because $a+b=1$ and $a-b=\omega$.

\begin{definition}[The unnormalised intertwiners]\label{ch14:def:Phi0}
$\Phi_0:\chXIVB\to T_{s_0+1}^+$ and $\Theta_0:T_{s_0}^+\to T_{s_0}^+$ are the unique $\sigma$-twisted
intertwiners extending $\varphi_\star$ and $\varphi_B$ (Chapter~\ref{ch:13}, \lean{sigmaV3T_exists_of_top},
\lean{sigmaTT_exists_of_top}).
\end{definition}
\begin{leancode}
noncomputable def Φ0 : VPlus →ₗ[ℂ] TP := (sigmaV3T_exists_of_top topV).choose
theorem Φ0_isSigma : IsSigmaV3T Φ0
theorem Φ0_top : ∀ x ∈ OmegaV3, Φ0 x = phiV x

noncomputable def Θ0 : TP →ₗ[ℂ] TP := (sigmaTT_exists_of_top topB t0_mem t0_ne).choose
theorem Θ0_isSigma : IsSigmaTT s0 Θ0
theorem Θ0_top : ∀ x ∈ OmegaT s0, Θ0 x = phiB x
\end{leancode}
(\leanpath{VOA/Moonshine/SigmaAssemble.lean}.)

\begin{dictionary}
\textbf{Decoding.}  \lean{.choose} picks a witness of an existence theorem; the accompanying theorems
record its properties.  Because the extension is \emph{unique} (an intertwiner is determined on a
cyclic vector, \lean{IsSigmaV3T.eqOn}), the choice is immaterial.  \lean{topV : TopSigmaV3 (s0 + 1) phiV} and
\lean{topB : TopSigmaTT s0 s0 phiB} are the top-level existence certificates of Chapter~\ref{ch:13}.

\textbf{Sector bookkeeping.}  $\Phi_0$ lands in $T_{s_0+1}^+$, so the triality label is $s=s_0+1$ and the
$\Theta$-sector is $T_{s+1}^+=T_{s_0}^+$, which is where $\Theta_0$ acts (the Lean identity
\lean{s0_add_one_add_one : s0 + 1 + 1 = s0}).
\end{dictionary}
\begin{physics}
$\Phi_0,\Theta_0$ are the triality maps defined by their action on the lowest states and extended by the chiral algebra $\chXIVA$ — exactly how one constructs a symmetry operator in CFT from its action on primaries.
\end{physics}


\subsection{\texorpdfstring{Properties of $\Theta_0$: involution and isometry, with no rescaling}{Properties of Theta0: involution and isometry, with no rescaling}}

\begin{theorem}\label{ch14:thm:psi}
(a) For a seed $t\in T_s$ and every $m$, $\psi_m^2=1$ on $T_s$, and $\psi_m$ is an isometry of the
invariant form $\langle\cdot,\cdot\rangle$ on $T$.  (b) Hence $\varphi_B$ is an isometric involution of
$(T_{s_0}^+)_2$, and (c) $\Theta_0^2=1$ on $T_{s_0}^+$ and $\Theta_0$ preserves \lean{TPForm} on the top.
\end{theorem}
\begin{leancode}
theorem psi_psi {s : ZMod 2} {t : TLeechT} (hts : t ∈ leechFrameT s) (ht0 : t ≠ 0) (ht : IsSeed t)
    (m : Fin 24) {x : TLeechT} (hx : x ∈ leechFrameT s) :
    psiOf (Bop s) m (psiOf (Bop s) m x) = x

theorem TForm_psi_psi {s : ZMod 2} {t : TLeechT} (hts : t ∈ leechFrameT s) (ht0 : t ≠ 0)
    (ht : IsSeed t) (m : Fin 24) {x y : TLeechT} (hx : x ∈ leechFrameT s) (hy : y ∈ leechFrameT s) :
    TForm (psiOf (Bop s) m x) (psiOf (Bop s) m y) = TForm x y

theorem Θ0_Θ0 : ∀ y ∈ TPframe s0, Θ0 (Θ0 y) = y

theorem Θ0_isom {x y : TP} (hx : x ∈ OmegaT s0) (hy : y ∈ OmegaT s0) :
    TPForm (Θ0 x) (Θ0 y) = TPForm x y
\end{leancode}
(\leanpath{VOA/Moonshine/SigmaAssembleTop.lean}, \leanpath{VOA/Moonshine/SigmaAssemble.lean}.)

\begin{dictionary}
\lean{TForm} is the symmetric invariant form on $T$ of Chapter~\ref{ch:08} (it satisfies
$\langle e_\alpha t,t'\rangle=(-1)^{|\alpha|^2/2}\langle t,e_\alpha t'\rangle$, \lean{TForm_eT}).
\lean{Fin 24} indexes the 24 coordinates.  The seed hypothesis is genuinely needed in (a): it is what makes
every joint eigenvector character satisfy the Golay constraint below.  Statements identical to (a)--(c).
\end{dictionary}
\begin{physics}
The triality acts on twisted ground states by a discrete Fourier-type transform whose eigenvalues are eighth roots of unity; on the physical (Golay-constrained) states it squares to one and is unitary.
\end{physics}


\paragraph{Proof on paper.}  Decompose $T_s$ into joint eigenvectors $z$ of the $Z_{0k}$.  On $z$ with
character $N$, $\psi_m$ acts by $u_m\omega^{|N|}$, so $\psi_m^2$ acts by $\omega^{2|N|}$.  The Golay constraint
of Chapter~\ref{ch:13} (\lean{golay_constraint}, \lean{four_dvd_card_of_orth}) says that the characters that
occur in $T_s$ — given the seed — satisfy $4\mid|N|$; then $\omega^{2|N|}=\omega^{8q}=1$.  Isometry: each $Z_{0k}$
is symmetric for the form (since $|\gamma^-|^2=4$), joint eigenvectors of different characters are
orthogonal, hence $\psi_m$ is symmetric; symmetric and involutive means isometric.  (b) The blocks
$h_{w_m}(-\tfrac12)\otimes T$ are mutually orthogonal for \lean{TPForm} with form $-4\langle\cdot,\cdot\rangle$
on the diagonal, and $\varphi_B$ acts blockwise by $\psi_m$.  (c) $\Theta_0^2$ is a ($\sigma^2=1$)-twisted,
i.e.\ honest, intertwiner which is the identity on the cyclic vector $x_0(t_0)$, hence the identity. \qed

\begin{status}
Unconditional; standard axioms only.  This \emph{proves} $\mu=+1$ (cf.\ \S\ref{ch14:sec:norm}): no
rescaling of $\Theta_0$ was needed, consistent with the numerics.
\end{status}

\subsection{\texorpdfstring{$\Phi_0$ is a nonzero similitude}{Phi0 is a nonzero similitude}}

\begin{theorem}[Similitude]\label{ch14:thm:simil}
Let $z_0=x_{n_0}$ for a fixed $n_0\in S_1$.  Then $B_V(z_0,z_0)\ne0$, and with
$\kappa_0=\langle\Phi_0z_0,\Phi_0z_0\rangle_T/B_V(z_0,z_0)$,
\begin{equation*}
  \langle\Phi_0a,\Phi_0b\rangle_T=\kappa_0\,B_V(a,b)\qquad\text{for all }a,b\in\chXIVB,
\end{equation*}
and $\kappa_0\ne0$.  Consequently $\Phi_0:\chXIVB\to T_{s_0+1}^+$ is bijective, and its inverse $\Psi_0$ (extended
linearly to $T^+$) is a $\sigma$-twisted intertwiner on $T_{s_0+1}^+$.
\end{theorem}
\begin{leancode}
noncomputable def κ0 : ℂ := TPForm (Φ0 z0) (Φ0 z0) / leechBV z0 z0

theorem Φ0_simil : ∀ a ∈ frameSector 1, ∀ b ∈ frameSector 1,
    TPForm (Φ0 a) (Φ0 b) = κ0 * leechBV a b

theorem κ0_ne : κ0 ≠ 0

theorem Φ0_inj {a b : VPlus} (ha : a ∈ frameSector 1) (hb : b ∈ frameSector 1)
    (h : Φ0 a = Φ0 b) : a = b
theorem Φ0_surj : ∀ y ∈ TPframe (s0 + 1), ∃ b ∈ frameSector 1, Φ0 b = y

noncomputable def Ψ0 : TP →ₗ[ℂ] VPlus := exists_Ψ0.choose
theorem Φ0_Ψ0 {y : TP} (hy : y ∈ TPframe (s0 + 1)) : Φ0 (Ψ0 y) = y
theorem Ψ0_Φ0 {x : VPlus} (hx : x ∈ frameSector 1) : Ψ0 (Φ0 x) = x
theorem Ψ0_isSigma : IsSigmaTV (s0 + 1) Ψ0
\end{leancode}
(\leanpath{VOA/Moonshine/SigmaAssemble.lean}.)

\begin{dictionary}
\textbf{Decoding.}  \lean{z0 := frameX0P (nvec 0)}.  Division in \lean{κ0} is honest field division in $\C$;
the theorem \lean{bv_z0 : leechBV z0 z0 ≠ 0} guarantees the denominator is nonzero, so \lean{κ0} is not
a junk value (in Lean, $x/0=0$, which would have made \lean{Φ0_simil} suspicious had the denominator
vanished).  \lean{Ψ0} is obtained by extending the inverse of the linear bijection
$\chXIVB\to T_{s_0+1}^+$ arbitrarily to a complement; only its values on $T_{s_0+1}^+$ ever matter.
\end{dictionary}
\begin{physics}
The triality map from odd untwisted states to twisted states preserves two-point functions up to one overall constant, which can be absorbed in the normalisation of the twisted sector.
\end{physics}


\paragraph{Proof on paper.}  Both $(a,b)\mapsto\langle\Phi_0a,\Phi_0b\rangle_T$ and $B_V$ are invariant
symmetric forms on $\chXIVB$ for which the zero mode $o(n(-1)^2\mathbb 1)$ ($n\in S_1$) is symmetric.  On the
top basis $x_m$ this zero mode has eigenvalue $\langle n,m\rangle^2$, and $\langle n,n\rangle^2=16\ne\langle n,m\rangle^2$
for $m\ne n$ in $S_1$ (since $\langle n,m\rangle=\pm4$ would force $n\mp m$ to have norm $0$).  Hence the basis is
orthogonal for both forms and each Gram matrix is determined by its row at $z_0$ restricted to the
top; by the top-level reduction (Theorem~\ref{ch14:thm:isotop}, scaled by $\kappa_0$), proportionality on
one row propagates to all of $\chXIVB$.  $B_V(z_0,z_0)\ne0$ because otherwise $z_0$ would be in the radical
of $B_V$ (by the same propagation with $\Phi=0$), contradicting nondegeneracy.  If $\kappa_0=0$, then
\lean{TPForm} would vanish on $\Phi_0(\chXIVB)=T_{s_0+1}^+$ (surjectivity: $\Phi_0$ hits a cyclic vector), where it
is nondegenerate.  Injectivity follows from $\kappa_0\ne0$ and nondegeneracy of $B_V$. \qed

\begin{status}Unconditional; standard axioms only.\end{status}

\subsection{\texorpdfstring{The triality data $D(c,\varepsilon)$}{The triality data D(c,epsilon)}}

\begin{definition}[$D(c,\varepsilon)$]\label{ch14:def:mkData}
For $c\in\C^\times$ and $\varepsilon$ with $\varepsilon^2=1$, the triality data
$D(c,\varepsilon)$ have $s=s_0+1$, $\Phi=c\Phi_0$, $\Psi=c^{-1}\Psi_0$, $\Theta=\varepsilon\Theta_0$.
\end{definition}
\begin{leancode}
noncomputable def mkData (c ε : ℂ) (hc : c ≠ 0) (hε : ε * ε = 1) : SigmaData where
  s := s0 + 1
  Φ := c • Φ0
  Ψ := c⁻¹ • Ψ0
  Θ := ε • Θ0
\end{leancode}
(\leanpath{VOA/Moonshine/SigmaAssemble.lean}; the proof fields are omitted — they are the properties
established above, rescaled.)

\begin{dictionary}
This is the whole normalisation freedom of \S\ref{ch14:sec:norm} made explicit.  The hypothesis
$\varepsilon\varepsilon=1$ allows exactly $\varepsilon=\pm1$.  Its existence shows that \lean{SigmaData}
is inhabited, so none of the earlier universally quantified results over $D$ is vacuous.
\end{dictionary}
\begin{physics}
The ambiguity in normalising a symmetry operator on sectors: an overall rescaling of a sector and a sign on a self-paired sector.
\end{physics}


\begin{theorem}[Isometries for the normalised data]\label{ch14:thm:isomk}
(a) There is $c$ with $c^2\kappa_0=1$.  (b) For such $c$ and either $\varepsilon$, \lean{PhiIso} holds for
$D(c,\varepsilon)$.  (c) For every $c\ne0$ and $\varepsilon$, \lean{ThetaTop} holds for $D(c,\varepsilon)$ with the
frame vectors $t_{s_0}=t_0$, $t_{s_0+1}=e_{m_0}t_0$.
\end{theorem}
\begin{leancode}
theorem exists_norm : ∃ c : ℂ, c * c * κ0 = 1

theorem phiIso_mkData {c ε : ℂ} (hc : c * c * κ0 = 1) (hε : ε * ε = 1) :
    PhiIso (mkData c ε (c_ne_of_norm hc) hε)

noncomputable def tF (s' : ZMod 2) : TLeechT := if s' = s0 then t0 else eL (mvec 0) t0

theorem thetaTop_mkData (c ε : ℂ) (hc : c ≠ 0) (hε : ε * ε = 1) :
    ThetaTop (mkData c ε hc hε) tF
\end{leancode}
(\leanpath{VOA/Moonshine/SigmaAssemble.lean}.)

\begin{dictionary}
\lean{eL (mvec 0)} is the operator $e_{m_0}$ for a fixed $m_0\in\Leech\setminus\Leech_0$, which swaps the frame
sectors of $T$ (\lean{tF_mem}); \lean{tF_ne} records $t_{s}\ne0$.  The theorem \lean{phiTop_of_phiIso}
converts \lean{PhiIso} into \lean{PhiTop} for any $D$.  These are exactly the needed residuals, now
\emph{proved}.
\end{dictionary}
\begin{physics}
With the right normalisation the triality is unitary on the relevant sectors; the possible obstruction $\mu=-1$ (an anti-unitary-like sign) does not occur.
\end{physics}


\paragraph{Proof on paper.}  (a) $\C$ is algebraically closed and $\kappa_0\ne0$.  (b)
$\langle c\Phi_0a,c\Phi_0b\rangle=c^2\kappa_0B_V(a,b)$.  (c) The rescaling by $\varepsilon$ squares away,
and $\Theta_0$ is an isometry on the top (Theorem~\ref{ch14:thm:psi}). \qed

\subsection{The one remaining hypothesis}

\begin{definition}[$c$-free self-coupling]\label{ch14:def:keyskew0}
With $x_K=x_0(t_{s_0+1})\in T_{s_0+1}^+$:
\lean{KeySkew0 ε} $:\iff$ $\varepsilon\,\Theta_0\bigl(Y^T(\Psi_0x_K,z)x_K\bigr)=W(x_K,z)\Psi_0x_K$.
\end{definition}
\begin{leancode}
noncomputable def xK : TP := x0T (tF (s0 + 1))

def KeySkew0 (ε : ℂ) : Prop :=
  ∀ p, ε • Θ0 (plusModule.YM (Ψ0 xK) p xK) = leechWsk xK p (Ψ0 xK)

theorem keySkew_mkData_iff (c ε : ℂ) (hc : c ≠ 0) (hε : ε * ε = 1) :
    KeySkew (mkData c ε hc hε) tF ↔ KeySkew0 ε

theorem moonshine_of_keySkew0 (h : KeySkew0 1 ∨ KeySkew0 (-1)) :
    ∃ V : VertexOperatorAlgebra Vnat,
      V.Y = Ynat VnatAxioms.wbarL ∧ V.vacuum = vacNat ∧ V.conformalVector = omegaNat ∧
      V.centralCharge = 24 ∧
      V.grading 1 = ⊥ ∧
      Module.finrank ℂ (V.grading 2) = 196884 ∧
      (∀ n : ℕ, (Module.finrank ℂ (V.grading (n : ℤ)) : ℂ) = PowerSeries.coeff n JShift) ∧
      V.IsSimple
\end{leancode}
(\leanpath{VOA/Moonshine/SigmaAssemble.lean}.)

\begin{dictionary}
\textbf{Why $c$ drops out.}  With $\Psi=c^{-1}\Psi_0$, both sides of \lean{KeySkew} are linear in
$\Psi x$, hence scale by $c^{-1}$; \lean{keySkew_mkData_iff} cancels it.  The disjunction in
\lean{moonshine_of_keySkew0} is honest: at this stage the sign was not determined in Lean.

\textbf{Hypotheses.}  Only \lean{h}.  The conclusion is the final theorem's.
\end{dictionary}
\begin{physics}
The single remaining OPE identity, now free of normalisation ambiguity except for one sign.
\end{physics}


\begin{status}
\lean{moonshine_of_keySkew0}: conditional on \lean{KeySkew0 1 ∨ KeySkew0 (-1)}; \texttt{\#print axioms}
(re-run): \lean{[propext, Classical.choice, Quot.sound]}.  \lean{mkData}, \lean{phiIso_mkData},
\lean{thetaTop_mkData}: unconditional.
\end{status}

\section{The self-coupling identity reduced to top-level output data}\label{ch14:sec:keyskew}

\subsection{The block defect}

\lean{KeySkew} is one entry of a larger object.  For $v,b\in\chXIVB$ the two candidate $B\to D'$ blocks
are $b\mapsto\Theta(Y^T(v,z)\Phi b)$ (from $\chXIVsig Y^\natural(v)\chXIVsig$) and
$b\mapsto W(\Phi v,z)b$ (from $Y^\natural(\chXIVsig v)$).

\begin{definition}[Block defect]\label{ch14:def:Dq}
For triality data $D$, $v,b\in V_\Leech^+$ and $q\in\Z$,
\begin{equation*}
  D_q(v,b)=\Theta\bigl(v_{(q)}\Phi b\bigr)-W(\Phi v)_{(q)}b\in T^+,
  \qquad
  W(t)_{(q)}\varphi=\sum_{i\ge0}\frac{(-1)^{q+i+1}}{i!}L_{-1}^i\bigl(\varphi_{(q+i)}t\bigr).
\end{equation*}
\end{definition}
\begin{leancode}
noncomputable def Dq (v b : VPlus) (q : ℤ) : TP :=
  D.Θ (mode (plusModule.YM v) q (D.Φ b)) - mode (leechWsk (D.Φ v)) q b
\end{leancode}
File \leanpath{VOA/Moonshine/KeySkewBlock.lean}.

\begin{dictionary}
\lean{mode (plusModule.YM v) q} is $v_{(q)}$ on $T^+$; \lean{mode (leechWsk x) q} is $W(x)_{(q)}$, whose
expansion in the display is the library theorem \lean{mode_leechWsk} (Chapter~\ref{ch:09}).  The
definition is made for all $v,b\in V_\Leech^+$ but used only for $v,b\in\chXIVB$, where
$D_q(v,b)\in T_{s+1}^+$ (\lean{Dq_mem}).  \lean{KeySkew D t} is exactly $D_q(\Psi x,\Psi x)=0$ for all $q$
(since $\Phi\Psi x=x$).
\end{dictionary}

\begin{physics}
$D_q(v,b)$ measures the failure of the triality to be a symmetry of one particular operator product:
the OPE of the untwisted operator $v$ with the untwisted operator $b$, both in the ``odd'' frame
sector $\chXIVB$, \emph{after} rotating by $\chXIVsig$, compared with the OPE of the twist field
$\Phi v$ with $b$.  In a genuine CFT with a symmetry $\chXIVsig$ these OPE coefficients agree, so
the defect is a quantity that a physicist expects to vanish identically; proving that it does is
the content of the rest of the chapter.
\end{physics}

Coefficient by coefficient, writing $v=\Psi x$, the output weight is $w=3-q$, and \lean{KeySkew} reads
\begin{align*}
 w=2:&\quad \Theta(o(v)x)=o(v)x,\\
 w=3:&\quad \Theta(v_{(0)}x)=L_{-1}(v_{(1)}x)-v_{(0)}x,\\
 w=4:&\quad \Theta(v_{(-1)}x)=v_{(-1)}x-L_{-1}v_{(0)}x+\tfrac12L_{-1}^2v_{(1)}x .
\end{align*}
Equivalently, with $K(z)=e^{-zL_{-1}/2}Y^T(v,z)x$, \lean{KeySkew} says $\Theta K_w=(-1)^wK_w$ on each weight
component (an analysis remark from \leanpath{Notes/KeySkew.md}, not formalised).

\begin{warning}[title={Evidence: Numerical evidence at weights 3 and 4}]
In the $N(A_1^{24})$ model (\leanpath{Notes/KeySkew.md} \S1, \leanpath{scripts/moonshine/numerics/keyskew/}),
the general block identity $\Theta Y^T(v,z)\Phi b=e^{zL_{-1}}Y^T(b,-z)\Phi v$ was tested through
lowering modes: at weight 3 on about $5300$ nonzero identities (residual $1.1\cdot10^{-11}$) and at
weight 4 on about $3300$ ($1.1\cdot10^{-11}$).  Controls fail decisively: $-\Theta$ gives residual $2.0$,
$\Theta\to\mathrm{Id}$ at least $0.32$, dropping $L_{-1}$ at least $0.75$.  This is evidence, not proof;
the tests see a $323$-dimensional slice per fibre, not the whole weight space.
\end{warning}

\subsection{The commutator formula and Li's lemma}

\begin{proposition}[Commutator formula for the defect]\label{ch14:prop:Dqcomm}
For $v,b\in\chXIVB$, $u\in\chXIVA$, $m,q\in\Z$ there is $J\in\N$ with
\begin{equation*}
  u_{(m)}D_q(v,b)=D_q(v,u_{(m)}b)+\sum_{r<J}\binom{m}{r}\,D_{m+q-r}\bigl((\sigma u)_{(r)}v,\,b\bigr).
\end{equation*}
\end{proposition}
\begin{leancode}
theorem Dq_comm {v b : VPlus} (hv : v ∈ frameSector 1) (hb : b ∈ frameSector 1)
    (u : V0Sub.carrier) (m q : ℤ) :
    ∃ J : ℕ, actT u m (Dq D v b q) = Dq D v (actV u m b) q
      + ∑ r ∈ Finset.range J, gBinom m r • Dq D (actV (sig u) r v) b (m + q - r)
\end{leancode}
File \leanpath{VOA/Moonshine/KeySkewBlock.lean}.

\begin{dictionary}
\lean{gBinom m r} is the generalised binomial coefficient $\binom mr=m(m-1)\cdots(m-r+1)/r!$, valid for
negative $m$.  The existential $J$ is any bound beyond which $(\sigma u)_{(r)}v=0$ (truncation).  This is
the textbook intertwining commutator formula, applied to both terms of the defect.
\end{dictionary}

\begin{physics}
This is the Ward identity for the $\chXIVA$-symmetry: moving a mode of a symmetry current past an
OPE coefficient produces the OPE coefficients of the transformed operators.  Both candidate blocks
satisfy the same Ward identity (one because $\Phi,\Theta$ intertwine up to $\sigma$, the other because
$W$ is built from the module vertex operator), hence so does their difference.
\end{physics}

\paragraph{Proof on paper.}  For the first term: $u_{(m)}\Theta=\Theta(\sigma u)_{(m)}$, then the module
commutator formula $[(\sigma u)_{(m)},v_{(q)}]=\sum_r\binom mr((\sigma u)_{(r)}v)_{(m+q-r)}$ on $T^+$, and
$(\sigma u)_{(m)}\Phi b=\Phi(u_{(m)}b)$.  For the second term: the commutator formula for the skew
intertwiner $W$, proved in Chapter~\ref{ch:09} (\lean{leechWsk_comm}), together with
$u_{(r)}\Phi v=\Phi((\sigma u)_{(r)}v)$.  Subtract. \qed

\begin{lemma}[Only the $i=0$ term of $W$ survives against weight 2]\label{ch14:lem:Wtop}
For $y\in(T^+)_2$, $t\in T^+$, $\varphi\in V_\Leech^+$:
$\langle y,W(t)_{(q)}\varphi\rangle_T=(-1)^{q+1}\langle y,\varphi_{(q)}t\rangle_T$.
\end{lemma}
\begin{leancode}
theorem TPForm_W_top {y : TP} (hy : y ∈ plusModule.weight (2 : ℂ)) (t : TP) (q : ℤ) (φ : VPlus) :
    TPForm y (mode (leechWsk t) q φ) = negOnePow (q + 1) * TPForm y (mode (plusModule.YM φ) q t)
\end{leancode}
File \leanpath{VOA/Moonshine/KeySkewBlock.lean}.

\begin{dictionary}
\lean{negOnePow k} is $(-1)^k$ for $k\in\Z$.  No hypothesis on $t$ or $\varphi$.
\end{dictionary}

\begin{physics}
The terms $L_{-1}^i(\cdots)$, $i\ge1$, are derivatives (descendants); they pair to zero with a
primary state of lowest weight, since $L_1$ annihilates $y$.  In CFT language: a two-point function
with a lowest-weight state only sees the primary part of the OPE.
\end{physics}

\paragraph{Proof.}  $\langle y,L_{-1}^iX\rangle=\langle L_1^iy,X\rangle=0$ for $i\ge1$, since $L_1y$ has weight $1$
and $T^+$ has nothing below weight $2$. \qed

\begin{theorem}[Li's lemma for the block defect]\label{ch14:thm:li}
Let $x'$ be a $\chXIVA$-cyclic vector of $T_{s+1}^+$.  If $\langle x',D_q(v,b)\rangle_T=0$ for all
$v,b\in\chXIVB$ and all $q$, then $D_q(v,b)=0$ for all $v,b\in\chXIVB$, $q\in\Z$.
\end{theorem}
\begin{leancode}
theorem Dq_eq_zero_of_pair {x' : TP}
    (hcyc : Submodule.span ℂ {y | ∃ u ∈ (Set.univ : Set V0Sub.carrier), ∃ n : ℤ, y = actT u n x'}
      = TPframe (D.s + 1))
    (hpair : ∀ v ∈ frameSector 1, ∀ b ∈ frameSector 1, ∀ q : ℤ, TPForm x' (Dq D v b q) = 0)
    {v b : VPlus} (hv : v ∈ frameSector 1) (hb : b ∈ frameSector 1) (q : ℤ) : Dq D v b q = 0
\end{leancode}
File \leanpath{VOA/Moonshine/KeySkewBlock.lean}.

\begin{dictionary}
The conclusion is a vector equation; the hypothesis is a family of scalar equations against one
vector.  \lean{hcyc} is cyclicity of $x'$.  Same as stated.
\end{dictionary}

\begin{physics}
A correlation function $\langle x'|\,\mathcal O\,\rangle$ with all possible insertions of symmetry
currents determines the state $\mathcal O$ when $x'$ generates the sector: the
``state--operator'' principle that states are determined by their overlaps with a complete set of
descendants of one cyclic state.  This is the mechanism of Li's uniqueness theorem for
intertwining operators.
\end{physics}

\paragraph{Proof on paper.}  The set of $y\in T_{s+1}^+$ with $\langle y,D_q(v,b)\rangle=0$ for all $v,b,q$
contains $x'$ and is stable under the modes of $\chXIVA$: by FHL invariance \eqref{ch14:eq:fhl},
$\langle u_{(n)}x',D\rangle$ is a combination of $\langle x',(L_1^ju)_{(\cdot)}D\rangle$, and by
Proposition~\ref{ch14:prop:Dqcomm} each $u'_{(m)}D_q(v,b)$ is a combination of defects, which pair to zero
with $x'$.  By cyclicity it is all of $T_{s+1}^+$.  The defect also pairs to zero with $T_s^+$ (the two
frame sectors are orthogonal), so it lies in the radical of the nondegenerate \lean{TPForm}. \qed

\begin{theorem}[Input reduction]\label{ch14:thm:pairb0}
Let $b_0\in(\chXIVB)_2$ be $\chXIVA$-cyclic.  If $\langle y,D_q(v,b_0)\rangle_T=0$ for all $v\in\chXIVB$, all
$q$ and all $y\in(T_{s+1}^+)_2$, then the same holds with $b_0$ replaced by any $b\in\chXIVB$.
\end{theorem}
\begin{leancode}
theorem pair_of_b0 {b₀ : VPlus} (hb₀ : b₀ ∈ frameSector 1) (hb₀2 : b₀ ∈ leechVOAPlus.grading 2)
    (hcyc : Submodule.span ℂ {y | ∃ u ∈ (Set.univ : Set V0Sub.carrier), ∃ n : ℤ, y = actV u n b₀}
      = frameSector 1)
    (h0 : ∀ v ∈ frameSector 1, ∀ q : ℤ, ∀ y ∈ TPframe (D.s + 1), y ∈ plusModule.weight (2 : ℂ) →
      TPForm y (Dq D v b₀ q) = 0)
    {b : VPlus} (hb : b ∈ frameSector 1) :
    ∀ v ∈ frameSector 1, ∀ q : ℤ, ∀ y ∈ TPframe (D.s + 1), y ∈ plusModule.weight (2 : ℂ) →
      TPForm y (Dq D v b q) = 0
\end{leancode}
File \leanpath{VOA/Moonshine/KeySkewBlock.lean}.

\begin{dictionary}
The only pairings tested are against \emph{top} vectors $y$ (weight $2$).  Hypotheses: cyclicity and
weight $2$ of $b_0$.  Same as stated.
\end{dictionary}

\begin{physics}
Inputs can be moved off the primary $b_0$ by symmetry Ward identities: every state of $\chXIVB$ is a
descendant of $b_0$ under the $\chXIVA$-algebra, and the top-level three-point functions of
descendants are fixed by those of the primary.
\end{physics}

\paragraph{Proof on paper.}  Write $b=u_{(n)}b_0$ for homogeneous $u$ (cyclicity).  By
Proposition~\ref{ch14:prop:Dqcomm}, $D_q(v,u_{(n)}b_0)=u_{(n)}D_q(v,b_0)-\sum_r\binom nr D_{n+q-r}((\sigma u)_{(r)}v,b_0)$;
the sum pairs to zero with $y$ by hypothesis (it is a defect with input $b_0$).  For the first term,
FHL invariance moves $u_{(n)}$ onto $y$, producing vectors of weight $n-\wt u+3$: if this is $2$, they are
top vectors and the hypothesis applies; if it is below $2$, they vanish.  If $u_{(n)}b_0$ itself has weight
below $2$ it is zero.  (The two lowest-weight facts used: $\chXIVB$ and $T^+$ start in weight $2$.) \qed

\begin{definition}[Top-output residual]\label{ch14:def:keytopv}
\lean{KeyTopV D t}: for every $k$ and every $v\in(\chXIVB)_k$,
$\Theta\bigl(o(v)x\bigr)=(-1)^k(\Psi x)_{(k-1)}\Phi(v)$ with $x=x_0(t_s)$.  Its $k=2$ slice
\lean{KeyTopTwo} is $\Theta(o(v)x)=o(\Psi x)\Phi(v)$.
\end{definition}
\begin{leancode}
def KeyTopV (t : ZMod 2 → TLeechT) : Prop :=
  ∀ (k : ℤ), ∀ v ∈ frameSector 1, v ∈ leechVOAPlus.grading k →
    D.Θ (mode (plusModule.YM v) (k - 1) (x0T (t D.s)))
      = negOnePow k • mode (plusModule.YM (D.Ψ (x0T (t D.s)))) (k - 1) (D.Φ v)

def KeyTopTwo (t : ZMod 2 → TLeechT) : Prop :=
  ∀ v ∈ frameSector 1, v ∈ leechVOAPlus.grading 2 →
    D.Θ (mode (plusModule.YM v) 1 (x0T (t D.s)))
      = mode (plusModule.YM (D.Ψ (x0T (t D.s)))) 1 (D.Φ v)
\end{leancode}
File \leanpath{VOA/Moonshine/KeySkewBlock.lean}.

\begin{dictionary}
\lean{mode (plusModule.YM v) (k - 1)} for $v$ of weight $k$ is the zero mode $o(v)$; applied to the weight-2
vector $x$ it gives a weight-2 vector.  The right side is the weight-2 output of $W(x)$ applied to $v$,
rewritten via Lemma~\ref{ch14:lem:Wtop}.  \lean{KeyTopV} is \emph{not finite}: $k$ is unbounded.
\end{dictionary}

\begin{physics}
\lean{KeyTopV} compares, for every operator $v$ of the sector $\chXIVB$, the \emph{primary-to-primary}
three-point coefficient $\langle\text{top}|\,v\,|x\rangle$ computed in the two pictures.  It is the
collection of all three-point functions with two legs on the lowest twisted states.
\end{physics}

\begin{theorem}[\lean{KeySkew} from top-level output data]\label{ch14:thm:keytopv}
If \lean{KeyTopV D t} holds (with $t_s\in T_s$ nonzero), then \lean{KeySkew D t}; consequently, together with
\lean{PhiTop} and \lean{ThetaTop}, all targets of the final theorem hold.
\end{theorem}
\begin{leancode}
theorem keySkew_of_keyTopV {t : ZMod 2 → TLeechT} (ht : ∀ s, t s ∈ leechFrameT s)
    (ht0 : ∀ s, t s ≠ 0) (hk : KeyTopV D t) : KeySkew D t
\end{leancode}
File \leanpath{VOA/Moonshine/KeySkewBlock.lean}; the corollary is \lean{moonshine_of_keyTopV}.

\begin{dictionary}
Hypotheses: frame membership and non-vanishing of $t_s$ (needed for cyclicity), and \lean{KeyTopV}.
\end{dictionary}

\begin{physics}
A ``bootstrap'' statement: all OPE coefficients of the defect are determined by its primary
three-point functions, so it vanishes once those do.
\end{physics}

\paragraph{Proof on paper.}  \lean{KeyTopV} gives, by linearity and a weight count,
$\langle y,D_q(v,\Psi x)\rangle=0$ for all top $y$, all $v$, $q$ (only $q=\wt v-1$ is nontrivial; the other
components have the wrong weight).  Theorem~\ref{ch14:thm:pairb0} with $b_0=\Psi x$ extends this to all
inputs $b$; in particular to $\langle x',D_q(v,b)\rangle=0$ for the cyclic $x'=x_0(t_{s+1})$; Li's lemma
gives $D\equiv0$, and in particular $D_q(\Psi x,\Psi x)=0$. \qed

\begin{status}
\lean{Dq_comm}, \lean{TPForm_W_top}, \lean{Dq_eq_zero_of_pair}, \lean{pair_of_b0}: unconditional given $D$.
\lean{keySkew_of_keyTopV}: conditional on \lean{KeyTopV}.  \texttt{\#print axioms} of \lean{Dq_eq_zero_of_pair}
(re-run): \lean{[propext, Classical.choice, Quot.sound]}.
\end{status}

\begin{warning}[title={Trap / dead end: Dead ends on the way to finiteness}]
(1) Reducing the \emph{state} slot $v$ by the same trick would need the iterate (associativity)
formula for both blocks; for $W$ that is the intertwining Jacobi identity, which the library does not
have, and even with it the recursion does not close on the top level.  (2) Induction on total weight
$\wt v+\wt b$ using skew-symmetry and \lean{Dq_comm} does not close: zero-mode reductions raise the
state weight.  (3) \lean{Modules.Fusion.Intertwiner} is commutator-only (no $L_{-1}$ or associativity
axiom), so its \lean{fusionCoeff} is \emph{not} the classical fusion rule; do not route uniqueness
through it.
\end{warning}

\section{The Frenkel--Zhu bimodule route}\label{ch14:sec:fz}

\subsection{Why bimodules}

Zhu's algebra $A(V)=V/O(V)$ (Chapter~\ref{ch:12}) captures the action of zero modes on lowest-weight
spaces.  Frenkel and Zhu (1992) extended this to a $V$-module $M$: there is an $A(V)$-\emph{bimodule}
$A(M)=M/O(M)$, and an intertwining operator of type $\binom{M_3}{M\ M_1}$ induces, on lowest-weight
spaces, a map $A(M)\otimes_{A(V)}M_1^{\mathrm{top}}\to M_3^{\mathrm{top}}$.  Li showed that the
intertwiner is \emph{determined} by this top-level map.  For us: \lean{KeyTopV} compares two such
top-level maps $\varphi_1,\varphi_2:\chXIVB\to(T_{s+1}^+)_2$.  If both factor through $A(\chXIVB)$ and are
compatible with the left and right $A(\chXIVA)$-actions, and if $A(\chXIVB)$ is generated as a bimodule by
its top level $(\chXIVB)_2$, then $\varphi_1=\varphi_2$ follows from equality on the top: a \emph{finite}
statement.

\subsection{The Frenkel--Zhu operations}

\begin{definition}[Frenkel--Zhu operations]\label{ch14:def:fz}
Let $W$ be a VOA, $M$ a weak $W$-module, $a\in W_p$ homogeneous, $m\in M$.
\begin{align*}
 a*m&=\Res_z z^{-1}(1+z)^pY(a,z)m=\sum_{i=0}^p\binom pi a_{(i-1)}m,\\
 m*a&=\Res_z z^{-1}(1+z)^{p-1}Y(a,z)m=\sum_{i=0}^p\binom{p-1}{i}a_{(i-1)}m,\\
 \operatorname{Res}_k(a,m)&=\Res_z z^{-2-k}(1+z)^{p+k}Y(a,z)m=\sum_{i=0}^{p+k}\binom{p+k}{i}a_{(i-2-k)}m\quad(k\ge0),
\end{align*}
and $O(M)$ is the span of all $\operatorname{Res}_k(a,m)$, $a$ homogeneous.
\end{definition}
\begin{leancode}
noncomputable def fzLeft (N : WeakModule W M) (p : ℕ) (a : 𝒱) (m : M) : M :=
  ∑ i ∈ Finset.range (p + 1), cbz (p : ℤ) (i : ℤ) • mode (N.YM a) ((i : ℤ) - 1) m

noncomputable def fzRight (N : WeakModule W M) (p : ℕ) (a : 𝒱) (m : M) : M :=
  ∑ i ∈ Finset.range (p + 1), cbz ((p : ℤ) - 1) (i : ℤ) • mode (N.YM a) ((i : ℤ) - 1) m

noncomputable def fzRes (N : WeakModule W M) (p k : ℕ) (a : 𝒱) (m : M) : M :=
  ∑ i ∈ Finset.range (p + k + 1), cbz ((p : ℤ) + k) (i : ℤ) • mode (N.YM a) ((i : ℤ) - 2 - k) m

def fzOgen (N : WeakModule W M) : Set M :=
  {x | ∃ (p k : ℕ) (a : 𝒱), a ∈ W.grading (p : ℤ) ∧ ∃ m : M, x = fzRes N p k a m}

def fzO (N : WeakModule W M) : Submodule ℂ M := Submodule.span ℂ (fzOgen N)
\end{leancode}
File \leanpath{VOA/Modules/FZBimodule.lean}.

\begin{dictionary}
\textbf{Decoding.}  \lean{cbz n i} is the binomial coefficient $\binom ni$ extended by $0$ for $i<0$
(\lean{cbz_apply_neg}); \lean{Finset.range (p+1)} is $\{0,\dots,p\}$.  The weight $p$ is passed as an
explicit natural number, and the formulas are meant for $a\in W_p$; \lean{fzOgen} imposes $a\in W_p$.
For $p=0$, \lean{fzRight} truncates $(1+z)^{-1}$ at $i=0$; since $W_0=\C\mathbb 1$ and
$\mathbb 1_{(i-1)}=\delta_{i0}$, this is the exact value.

\textbf{Difference from Frenkel--Zhu.}  FZ's $O(M)$ is spanned by the $k=0$ residues $a\circ m$ only.
The library spans by all $k\ge0$.  FZ's Lemma 2.1.2 says the spans agree; this is \emph{not} formalised in
\lean{FZBimodule.lean}, but the larger spanning set only makes the generation hypothesis below
\emph{weaker} (it quantifies over subspaces containing a \emph{bigger} $O(M)$), and the module form of
Lemma 2.1.2 is in fact proved in \lean{FZGenFiltration.lean} (\lean{PresM_mem}, below).  The bimodule
structure on $M/O(M)$ itself ($O(W)*M\subseteq O(M)$ etc.) is \emph{not} formalised and \emph{not
needed}: only closure of subspaces under $a*\cdot$, $\cdot*a$ is used.
\end{dictionary}

\begin{physics}
$a*m$ and $m*a$ are the zero-mode actions ``from the left and from the right'': on a correlation
function $\langle\text{top}_3|\,\mathcal Y(m,z)\,|\text{top}_1\rangle$, inserting $a$ near the outgoing
state or near the incoming state.  The bimodule $A(M)$ is the space of ``primary fields up to
descendants'', the algebraic shadow of the fusion rule $M\times M_1\to M_3$.
\end{physics}

\begin{proposition}[Exact commutator]\label{ch14:prop:fzcomm}
$a*m-m*a=\sum_{i=0}^{p-1}\binom{p-1}{i}a_{(i)}m$ exactly (no $O(M)$ correction).
\end{proposition}
\begin{leancode}
theorem fzLeft_sub_fzRight (N : WeakModule W M) (p : ℕ) (a : 𝒱) (m : M) :
    fzLeft N p a m - fzRight N p a m
      = ∑ i ∈ Finset.range p, cbz ((p : ℤ) - 1) (i : ℤ) • mode (N.YM a) (i : ℤ) m
\end{leancode}
File \leanpath{VOA/Modules/FZBimodule.lean}.

\begin{dictionary}
Holds for every $a$ (the identity is pure Pascal arithmetic, $\binom pi-\binom{p-1}i=\binom{p-1}{i-1}$).
For $a=\omega$ it is $\omega*m-m*\omega=(L_{-1}+L_0)m$.
\end{dictionary}

\begin{physics}
The left and right actions differ by the ``non-negative modes'' of $a$, i.e.\ by the singular part of the
OPE of $a$ with $m$; for the stress tensor this is translation plus scaling of the field $m$.
\end{physics}

\begin{definition}[Bimodule generation]\label{ch14:def:fzgen}
An ordinary $W$-module $M$ is \emph{FZ-generated by its weight-$h$ space} if every subspace
$S\subseteq M$ that contains $O(M)$ and $M_h$, and is closed under $a*\cdot$ and $\cdot*a$ for all
homogeneous $a\in W$, is all of $M$.
\end{definition}
\begin{leancode}
def FZGenerated (N : OrdVModule W M) (h : ℂ) : Prop :=
  ∀ S : Submodule ℂ M, fzO N.toWeakModule ≤ S → N.weight h ≤ S →
    (∀ (p : ℕ) (a : 𝒱), a ∈ W.grading (p : ℤ) → ∀ m ∈ S,
      fzLeft N.toWeakModule p a m ∈ S ∧ fzRight N.toWeakModule p a m ∈ S) → S = ⊤
\end{leancode}
File \leanpath{VOA/Modules/FZBimodule.lean}.

\begin{dictionary}
\lean{OrdVModule} is an ordinary (graded, $L_0$-semisimple) module (Definition~\ref{l12:def:ordmodule});
\lean{N.weight h} is its weight-$h$ space.  The property says $A(M)$ is generated by the image of $M_h$
as an $A(W)$-bimodule, phrased without forming the quotient.  It is a property of the pair
$(W,M)$ only.  It could be trivially true only if $M=M_h+O(M)$, which is not the case here (but
nothing is lost if it were: it is a hypothesis to be \emph{proved}).
\end{dictionary}

\begin{physics}
``Every field of the sector $M$ is, modulo descendants, obtained from the lowest-weight fields by
fusing with chiral-algebra fields on either side.''  This is what makes the fusion rule detectable on
the lowest states.
\end{physics}

\subsection{Transposed top-level identities from the commutator formula alone}

\begin{definition}[Commutator data]\label{ch14:def:commdata}
Operators $A_1(r)$ on $M_1$, $A_3(r)$ on $M_3$, $B(n):M_1\to M_3$ and $C_j(n):M_1\to M_3$ form
\emph{commutator data with bound $J$} if
$A_3(r)B(n)-B(n)A_1(r)=\sum_{j<J}\binom rj C_j(r+n-j)$ for all $r,n$.
\end{definition}
\begin{leancode}
def CommData : Prop :=
  ∀ (r n : ℤ) (m : M₁), A₃ r (B n m) - B n (A₁ r m)
    = ∑ j ∈ Finset.range J, gBinom r j • C j (r + n - j) m
\end{leancode}
File \leanpath{VOA/Modules/FZBimodule.lean}.

\begin{dictionary}
Section variables: $A_1,A_3,B,C,J$.  This is the commutator formula of an intertwining operator
$B(n)=\mathcal Y(b)_{(n)}$ with $C_j=\mathcal Y(a_{(j)}b)$, with \emph{no} associativity or
$L_{-1}$-derivative axiom.
\end{dictionary}

\begin{physics}
Only the singular part of the OPE of a current with a field is assumed (the Borcherds commutator
formula); the full OPE associativity is not.
\end{physics}

\begin{theorem}[Skew top-level identities]\label{ch14:thm:skew3}
Let the data satisfy \lean{CommData} with bound $J=p$, and let $B(r)m=0$ for $r>n$.  Then
\begin{align*}
\sum_{i=0}^{p}(-1)^{p-i}\tbinom pi\,B(n+p-i)A_1(i-1)m&=A_3(p-1)B(n)m,\\
\sum_{i=0}^{p-1}(-1)^{p-1-i}\tbinom{p-1}{i}\,B(n+p-1-i)A_1(i)m&=A_3(p-1)B(n)m-C_{p-1}(n)m\quad(p\ge1),\\
\sum_{i=0}^{p+k}(-1)^{p+k-i}\tbinom{p+k}{i}\,B(n+1+p+k-i)A_1(i-2-k)m&=0 .
\end{align*}
\end{theorem}
\begin{leancode}
theorem skew_left (p : ℕ) (hc : CommData A₁ A₃ B C p) (n : ℤ) (m : M₁)
    (hlow : ∀ r : ℤ, n < r → B r m = 0) :
    ∑ i ∈ Finset.range (p + 1), (asgn p i : ℂ) • B (n + p - i) (A₁ ((i : ℤ) - 1) m)
      = A₃ ((p : ℤ) - 1) (B n m)

theorem skew_comm (p : ℕ) (hp : 1 ≤ p) (hc : CommData A₁ A₃ B C p) (n : ℤ) (m : M₁)
    (hlow : ∀ r : ℤ, n < r → B r m = 0) :
    ∑ i ∈ Finset.range p, (asgn (p - 1) i : ℂ) • B (n + (p - 1 : ℕ) - i) (A₁ (i : ℤ) m)
      = A₃ ((p : ℤ) - 1) (B n m) - C (p - 1) n m

theorem skew_res (p k : ℕ) (hc : CommData A₁ A₃ B C p) (n : ℤ) (m : M₁)
    (hlow : ∀ r : ℤ, n < r → B r m = 0) :
    ∑ i ∈ Finset.range (p + k + 1), (asgn (p + k) i : ℂ)
      • B (n + 1 + (p + k : ℕ) - i) (A₁ ((i : ℤ) - 2 - k) m) = 0
\end{leancode}
File \leanpath{VOA/Modules/FZBimodule.lean}.

\begin{dictionary}
\lean{asgn N i} is the signed binomial $(-1)^{N-i}\binom Ni$.  \lean{hlow} says $B(n)m$ is the top
(highest-index nonzero) component.  These are the FZ identities $o(a*m)=o(a)o(m)$,
$o(m*a)=o(m)o(a)$, $o(O(M))=0$ for the \emph{transposed} operator $m\mapsto e^{zL_{-1}}Y(\cdot,-z)m$, read on
the top level.
\end{dictionary}

\begin{physics}
Skew-symmetry $\mathcal Y(m,z)m_1=e^{zL_{-1}}\mathcal Y(m_1,-z)m$ exchanges the roles of the two
inputs.  These identities say that, at the level of primary three-point functions, the transposed
operator still respects the bimodule structure — and this only needs Ward identities.
\end{physics}

\paragraph{Proof on paper.}  Sum the commutator formula with the weights $(-1)^{N-i}\binom Ni$; the
$C$-terms combine by the finite-difference identity $\sum_i(-1)^{N-i}\binom Ni\binom{s+i}j=\binom{s}{j-N}$,
leaving one $C$-term or none according to the residue taken; the $A_3B$-terms collapse to the single
top one by \lean{hlow}. \qed

\subsection{\texorpdfstring{The reduction of \lean{KeyTopV} to bimodule generation and top pairs}{The reduction of KeyTopV to bimodule generation and top pairs}}

For a top input $y\in(T_s^+)_2$ define two maps $\chXIVB\to T^+$ on homogeneous $v\in(\chXIVB)_k$:
\begin{equation*}
\varphi_1^y(v)=\Theta\bigl(o(v)y\bigr),\qquad \varphi_2^y(v)=(-1)^k(\Psi y)_{(k-1)}\Phi(v),\qquad F_y=\varphi_1^y-\varphi_2^y .
\end{equation*}
\begin{definition}\label{ch14:def:Fy}
\begin{leancode}
noncomputable def phi2Fam (y : TP) (k : ℤ) : VPlus →ₗ[ℂ] TP :=
  negOnePow k • ((mode (plusModule.YM (D.Ψ y)) (k - 1)) ∘ₗ D.Φ)

noncomputable def phi2 (y : TP) : VPlus →ₗ[ℂ] TP := gradedExt leechVOAPlus (phi2Fam D y)

noncomputable def phi1 (y : TP) : VPlus →ₗ[ℂ] TP := D.Θ ∘ₗ (oMod plusModule).flip y

noncomputable def Fy (y : TP) : VPlus →ₗ[ℂ] TP := phi1 D y - phi2 D y
\end{leancode}
File \leanpath{VOA/Moonshine/FZBimoduleKey.lean}.
\end{definition}

\begin{dictionary}
\lean{oMod plusModule v y} is the zero mode $o(v)y$, extended linearly over the grading of $v$.
\lean{gradedExt} extends a family of maps, one per weight, linearly over the direct-sum decomposition
of $V_\Leech^+$ (\lean{gradedExt_of_mem}: on $V_k$ it is the $k$-th map).  So these are exactly
$\varphi_1^y,\varphi_2^y,F_y$.
\end{dictionary}

\begin{physics}
$\varphi_1^y(v)$ is the top-level matrix element of $v$ between twisted lowest states, rotated by the
triality; $\varphi_2^y(v)$ is the matrix element of the twist field $\Phi v$ seen through skew-symmetry.
$F_y$ is their difference — the candidate violation of the triality symmetry at the level of primary
three-point couplings.
\end{physics}

\begin{theorem}[Bimodule compatibility of $F_y$]\label{ch14:thm:Fy}
For $a\in(\chXIVA)_p$, $y\in(T_s^+)_2$, $v\in(\chXIVB)_k$:
\begin{enumerate}[label=(\alph*)]
\item $F_y(a*v)=o(\sigma a)F_y(v)$;
\item for $p=q+1$: $F_y\bigl(\sum_{i\le q}\binom qi a_{(i)}v\bigr)=(\sigma a)_{(q)}F_y(v)-F_{a_{(q)}y}(v)$, hence
  $F_y(v*a)=F_{o(a)y}(v)$;
\item $F_y$ kills $O(\chXIVB)$.
\end{enumerate}
\end{theorem}
\begin{leancode}
theorem Fy_lV {p : ℕ} {a : V0Sub.carrier} (ha : (a : VPlus) ∈ leechVOAPlus.grading (p : ℤ))
    {y : TP} (hy : y ∈ OmegaT D.s) {v : VPlus} (hv : v ∈ frameSector 1) {k : ℤ}
    (hk : v ∈ leechVOAPlus.grading k) :
    Fy D y (lV p a v) = actT (sig a) ((p : ℤ) - 1) (Fy D y v)

theorem Fy_cV {q : ℕ} {a : V0Sub.carrier} (ha : (a : VPlus) ∈ leechVOAPlus.grading ((q + 1 : ℕ) : ℤ))
    {y : TP} (hy : y ∈ OmegaT D.s) {v : VPlus} (hv : v ∈ frameSector 1) {k : ℤ}
    (hk : v ∈ leechVOAPlus.grading k) :
    Fy D y (cV (q + 1) a v) = actT (sig a) (q : ℤ) (Fy D y v) - Fy D (actT a (q : ℤ) y) v

theorem Fy_resV {p k' : ℕ} {a : V0Sub.carrier} (ha : (a : VPlus) ∈ leechVOAPlus.grading (p : ℤ))
    {y : TP} (hy : y ∈ OmegaT D.s) {v : VPlus} (hv : v ∈ frameSector 1) {k : ℤ}
    (hk : v ∈ leechVOAPlus.grading k) : Fy D y (resV p k' a v) = 0
\end{leancode}
File \leanpath{VOA/Moonshine/FZBimoduleKey.lean}.

\begin{dictionary}
\lean{lV p a}, \lean{cV p a}, \lean{resV p k a} are $v\mapsto a*v$, $v\mapsto\sum_{i<p}\binom{p-1}ia_{(i)}v$ and
$v\mapsto\operatorname{Res}_k(a,v)$ computed inside $V_\Leech^+$; they agree with \lean{fzLeft}, \lean{fzLeft - fzRight},
\lean{fzRes} for the module \lean{MV3} ($\chXIVB$ as a $\chXIVA$-module), by \lean{coe_fzLeft} etc.
\lean{actT (sig a) (p - 1)} is $o(\sigma a)$ on $T^+$ (weight $p$).  Note in (b) that $a_{(q)}y=o(a)y$ is again top.
\end{dictionary}

\begin{physics}
These are the Ward identities for the primary three-point couplings: inserting a chiral-algebra
field on the left of $v$ acts on the outgoing state, on the right acts on the incoming state.
\end{physics}

\paragraph{Proof on paper.}  For $\varphi_1$: $a*v$, $\operatorname{Res}_k(a,v)$ are computed inside the VOA
$V_\Leech^+$, and on the lowest-weight space of the $V_\Leech^+$-module $T^+$ Zhu's theorem gives
$o(a*v)=o(a)o(v)$ and $o(O(V_\Leech^+))=0$ (Chapter~\ref{ch:12}); then $\Theta o(a)=o(\sigma a)\Theta$.
(b) is the module commutator formula.  For $\varphi_2$: $\Phi$ turns $a_{(i)}v$ into $(\sigma a)_{(i)}\Phi v$, and
then Theorem~\ref{ch14:thm:skew3} applies to the operators $B(n)=(\Psi y)_{(n)}$ on $T^+$, with
$A_1=A_3=(\sigma a)_{(\cdot)}$ and bound $J=p$: the commutator sum stops at $p$ because
$(\sigma a)_{(j)}\Psi y$ has weight $p+2-j-1<2$ for $j\ge p$, and $(\Psi y)_{(r)}\Phi v=0$ above the top mode
for the same reason.  The sign $(-1)^k$ in $\varphi_2$ matches the signs \lean{asgn}. \qed

\begin{definition}[The three named Props]\label{ch14:def:toppairs}
\lean{KeyTopAll D}: $F_y=0$ for every $y\in(T_s^+)_2$ (i.e.\ \lean{KeyTopV} for every top input).
\lean{BimodGenV3}: $\chXIVB$ is FZ-generated by $(\chXIVB)_2$ as a $\chXIVA$-module.
\lean{TopPairs D}: $\Theta\bigl(o(v)\Phi b\bigr)=o(b)\Phi v$ for all $v,b\in(\chXIVB)_2$.
\end{definition}
\begin{leancode}
def KeyTopAll : Prop :=
  ∀ y ∈ OmegaT D.s, ∀ (k : ℤ), ∀ v ∈ frameSector 1, v ∈ leechVOAPlus.grading k →
    D.Θ (mode (plusModule.YM v) (k - 1) y)
      = negOnePow k • mode (plusModule.YM (D.Ψ y)) (k - 1) (D.Φ v)

def BimodGenV3 : Prop := FZGenerated MV3 (2 : ℂ)

def TopPairs : Prop :=
  ∀ v ∈ OmegaV3, ∀ b ∈ OmegaV3,
    D.Θ (mode (plusModule.YM v) 1 (D.Φ b)) = mode (plusModule.YM b) 1 (D.Φ v)
\end{leancode}
File \leanpath{VOA/Moonshine/FZBimoduleKey.lean}.

\begin{dictionary}
\lean{MV3} is $\chXIVB$ as an ordinary \lean{V0P}-module (Chapter~\ref{ch:12}).  \lean{OmegaV3} $=(\chXIVB)_2$.
In \lean{TopPairs}, \lean{mode (plusModule.YM v) 1} is $o(v)$ for weight-2 $v$.  \lean{TopPairs} is a finite
identity on a $49152$-dimensional space (bilinear in $v,b$).
\end{dictionary}

\begin{physics}
\lean{TopPairs} is the statement that the primary couplings of two weight-2 operators of $\chXIVB$ to
the lowest twisted states are symmetric under exchange of the two operators, once the triality is
applied: a finite ``crossing symmetry'' of OPE coefficients.
\end{physics}

\begin{theorem}[The Frenkel--Zhu reduction and its converse]\label{ch14:thm:fzred}
\begin{enumerate}[label=(\alph*)]
\item \lean{BimodGenV3} and \lean{TopPairs D} imply \lean{KeyTopAll D}.
\item \lean{KeyTopAll D} implies \lean{TopPairs D} (so \lean{TopPairs} is also necessary).
\item Hence \lean{BimodGenV3} and \lean{TopPairs D} imply \lean{KeySkew D t}, and with \lean{PhiTop},
  \lean{ThetaTop}, all targets.
\end{enumerate}
\end{theorem}
\begin{leancode}
theorem keyTopAll_of_FZ (hgen : BimodGenV3) (htop : TopPairs D) : KeyTopAll D

theorem topPairs_of_keyTopAll (h : KeyTopAll D) : TopPairs D

theorem keySkew_of_FZ {t : ZMod 2 → TLeechT} (ht : ∀ s, t s ∈ leechFrameT s) (ht0 : ∀ s, t s ≠ 0)
    (hgen : BimodGenV3) (htop : TopPairs D) : KeySkew D t

theorem moonshine_of_FZ {β : leech.carrier} (hβ : β ∉ lambda0) {t : ZMod 2 → TLeechT}
    (ht : ∀ s, t s ∈ leechFrameT s) (ht0 : ∀ s, t s ≠ 0)
    (hΦ : PhiTop D t) (hΘ : ThetaTop D t) (hgen : BimodGenV3) (htop : TopPairs D) :
    ∃ V : VertexOperatorAlgebra Vnat,
      V.Y = Ynat VnatAxioms.wbarL ∧ V.vacuum = vacNat ∧ V.conformalVector = omegaNat ∧
      V.centralCharge = 24 ∧
      V.grading 1 = ⊥ ∧
      Module.finrank ℂ (V.grading 2) = 196884 ∧
      (∀ n : ℕ, (Module.finrank ℂ (V.grading (n : ℤ)) : ℂ) = PowerSeries.coeff n JShift) ∧
      V.IsSimple
\end{leancode}
File \leanpath{VOA/Moonshine/FZBimoduleKey.lean}.

\begin{dictionary}
All hypotheses appear explicitly.  (b) guarantees that the reduction did not replace the goal by a
\emph{stronger} unprovable statement: \lean{TopPairs} is exactly what must be true.
\end{dictionary}

\begin{physics}
Fusion-rule logic: the space of intertwiners $\chXIVB\times T_s^+\to T_{s+1}^+$ injects into top-level
data, which is generated by the lowest states; so two intertwiners agree once their lowest couplings
do.  In CFT one would say the fusion coefficient $N=1$ makes the coupling unique up to scale; here
uniqueness is replaced by the more economical statement that the \emph{difference} vanishes.
\end{physics}

\paragraph{Proof on paper.}  (a) Let $S=\{v\in\chXIVB:F_y(v)=0\ \forall y\in(T_s^+)_2\}$.  By
Theorem~\ref{ch14:thm:Fy}, $S$ contains $O(\chXIVB)$ and is closed under $a*\cdot$ and $\cdot*a$ (for the right
action use $v*a=a*v-(a*v-v*a)$, and $F_y(v*a)=F_{o(a)y}(v)$ with $o(a)y$ top).  On the top, $v\in(\chXIVB)_2$:
$F_y(v)=\Theta o(v)y-o(\Psi y)\Phi v$, which is \lean{TopPairs} at $(v,\Psi y)$ since $\Phi\Psi y=y$.  So $S$
contains $(\chXIVB)_2$, and \lean{BimodGenV3} gives $S=\chXIVB$.  (b) Take $y=\Phi b$, $k=2$.  (c) Specialise
$y=x$ and apply Theorem~\ref{ch14:thm:keytopv}. \qed

\begin{warning}[title={Trap / dead end: How the intertwining Jacobi identity of $W$ was avoided}]
The obvious route needed the three FZ top-level identities for the intertwiner $W$, i.e.\ its full
intertwining Jacobi identity, which the library lacks (\leanpath{Notes/KeySkew.md} \S3; this was the
recorded main risk).  The trick: $\varphi_2$ is the top component of $W(\Phi v,z)\Psi y$, but it is
\emph{written} through the module modes $(\Psi y)_{(n)}$ acting on $\Phi v$, and $\Phi$ turns $a*v$ into an
expression in modes of $\sigma a$.  The needed identities are then the transposed ones
(Theorem~\ref{ch14:thm:skew3}), which require only the commutator formula of $Y^T$ and the lowest
weights.  Neither the Jacobi identity of $W$, nor a transpose through \lean{TPForm}, nor
\lean{Fusion.Intertwiner} is used.
\end{warning}

\begin{status}
\lean{Fy_lV}, \lean{Fy_cV}, \lean{Fy_resV}, \lean{topPairs_of_keyTopAll}: unconditional given $D$.
\lean{keyTopAll_of_FZ}, \lean{keySkew_of_FZ}: conditional on \lean{BimodGenV3} and \lean{TopPairs D}.
\texttt{\#print axioms} \lean{VOA.Moonshine.FZKey.keySkew_of_FZ} (re-run): \lean{[propext, Classical.choice, Quot.sound]}.
\end{status}

\subsection{\texorpdfstring{Bimodule generation of $A(\chXIVB)$}{Bimodule generation of A(V3+)}}

\begin{warning}[title={Evidence: Numerical de-risking before formalisation}]
\lean{BimodGenV3} was the recorded ``real risk'': the analogous strong-generation statement had
\emph{failed} for $\chXIVA$ itself.  It was first tested (\leanpath{Notes/BimoduleGenTest.md},
\leanpath{scripts/moonshine/numerics/bimodule/}) by the filtered criterion below at weights 3, 4, 5, exact
arithmetic mod $2^{31}-1$: full rank in every excited sector; in the ground sector every coset vector
$\nu$ has a partner with $f=-1$ (counts at weight 3: $276$; weight 4: $2048$; weight 5: $11178$); the pair
lemma checked on $3937$ patterns up to weight $30$.  \emph{Control}: using only the left action
(no $0$-modes) gives rank $0$ of $24$ at weight $3$ — the right action is essential.  This is evidence;
the Lean proof below is the proof.
\end{warning}

\begin{definition}[Weight-raising products]\label{ch14:def:lowerprod}
For an ordinary $W$-module $M$ and $n\in\Z$, \lean{lowerProdM N n} is the set of $a_{(j)}x$ with $a\in W_p$,
$x\in M_q$, $q<n$, $j\le0$, $p+q-j-1=n$.
\end{definition}
\begin{leancode}
def lowerProdM (N : OrdVModule W M) (n : ℤ) : Set M :=
  {y | ∃ p q j : ℤ, q < n ∧ j ≤ 0 ∧ p + q - j - 1 = n ∧
    ∃ a ∈ W.grading p, ∃ x ∈ N.weight (q : ℂ), y = N.act a j x}
\end{leancode}
File \leanpath{VOA/Moonshine/FZGenFiltration.lean}.

\begin{dictionary}
\lean{N.act a j x} is $a_{(j)}x$.  $p+q-j-1$ is the weight of $a_{(j)}x$.  So these are products landing in
weight $n$ built from strictly lower-weight module states using non-positive modes.
\end{dictionary}

\begin{physics}
States obtained by acting with creation-type modes ($j\le0$) of chiral fields on lower states:
``descendants and composites'' of lower-weight fields.
\end{physics}

\begin{theorem}[Filtered generation criterion]\label{ch14:thm:crit}
Let $M$ be an ordinary $W$-module with integer weights bounded below by $n_0$, and $G\subseteq M$ with
$M_n\subseteq\operatorname{span}(G\cup\lean{lowerProdM}\,n)$ for all $n$.  Then every subspace $S\supseteq G$
containing all $\operatorname{Res}_0(a,x)$ and closed under $a*\cdot$, $\cdot*a$ ($p\ge1$) is $M$.  Along the way:
(module FZ Lemma 2.1.2) if a subspace contains all $\operatorname{Res}_0(a,x)$ for fixed $x$, it contains all
$\operatorname{Res}_k(a,x)$.
\end{theorem}
\begin{leancode}
theorem PresM_mem (N : OrdVModule W M) (x : M) (X : Submodule ℂ M)
    (hbase : ∀ (p : ℕ) (a : 𝒱), a ∈ W.grading (p : ℤ) → fzCirc N p a x ∈ X) (k : ℕ) :
    ∀ (p : ℕ) (a : 𝒱), a ∈ W.grading (p : ℤ) → PresM N p a x k ∈ X

theorem eq_top_of_lowerSpanM (N : OrdVModule W M) (G : Set M) (n₀ : ℤ)
    (hint : ∀ α : ℂ, N.weight α ≠ ⊥ → ∃ k : ℤ, α = k)
    (hbot : ∀ k : ℤ, k < n₀ → N.weight (k : ℂ) = ⊥)
    (hG : ∀ n : ℤ, N.weight (n : ℂ) ≤ span ℂ (G ∪ lowerProdM N n))
    (S : Submodule ℂ M)
    (hO : ∀ (p : ℕ) (a : 𝒱), a ∈ W.grading (p : ℤ) → ∀ x : M, fzCirc N p a x ∈ S)
    (hL : ∀ (p : ℕ) (a : 𝒱), a ∈ W.grading (p : ℤ) → ∀ x ∈ S, fzLeft N p a x ∈ S)
    (hR : ∀ (p : ℕ) (a : 𝒱), 1 ≤ p → a ∈ W.grading (p : ℤ) → ∀ x ∈ S, fzRight N p a x ∈ S)
    (hGS : G ⊆ S) : S = ⊤
\end{leancode}
File \leanpath{VOA/Moonshine/FZGenFiltration.lean}.

\begin{dictionary}
\lean{fzCirc N p a x} $=\operatorname{Res}_0(a,x)=a\circ x$ and \lean{PresM N p a x k} $=\operatorname{Res}_k(a,x)$;
here these are local copies in namespace \lean{FZGen}, with \lean{fzRight} summing $i<p$ (it agrees with
\lean{FZ.fzRight} for $p\ge1$, \lean{fzRight_eq}).  \lean{hint}: all weights are integers; \lean{hbot}: nothing
below $n_0$.  Hypothesis \lean{hO} uses only FZ's original $O(M)$ generators — strictly fewer than
\lean{fzO}.
\end{dictionary}

\begin{physics}
Induction on conformal weight: if every field of weight $n$ is primary-generated or a composite of
lower fields, and composites are reachable by left/right fusion and descendants, then everything is.
\end{physics}

\paragraph{Proof on paper.}  Lemma 2.1.2 by induction on $k$, using
$Y(L_{-1}a,z)=\partial_zY(a,z)$ and a Pascal identity for the coefficients.  Then induction on $n$: if
$M_{<n}\subseteq S$, each $a_{(j)}x\in\lean{lowerProdM}\,n$ lies in $S$ — for $j=-1$ because $a*x-a_{(-1)}x$
has lower weight; for $j\le-2$ because $a_{(j)}x-\operatorname{Res}_{-2-j}(a,x)$ has lower weight; for $j=0$
because $a*x-x*a-a_{(0)}x$ has lower weight (the exact commutator). \qed

\begin{theorem}[Charged excited spanning, generic form]\label{ch14:thm:excited}
For $\gamma\in\Leech\setminus\{0\}$ and any family $T(n)$ of subspaces of the Fock space with the closure
property \lean{QuadClosed γ T}, every Fock state of weight $w\ge1$ lies in $T(w+\tfrac12|\gamma|^2)$.
\end{theorem}
\begin{leancode}
def QuadClosed (γ : leech.carrier) (T : ℤ → Submodule ℂ Fk) : Prop :=
  ∀ (i j : ℕ) (a b : Fin (rank leech)) (n m : ℤ), m ≤ 0 → ∀ (w : ℕ) (p : Fk),
    p ∈ gradingWN (rank leech) w → (i : ℤ) + j + 2 < n → (w : ℤ) + LatGen.halfNorm leech γ < n →
    ((i : ℤ) + j + 2) + ((w : ℤ) + LatGen.halfNorm leech γ) - m - 1 = n →
    quadF γ a i b j (-m - 1) p ∈ T n

theorem gradingWN_le_of_quadClosed (hT : QuadClosed γ T) (hγ0 : γ ≠ 0) {w : ℕ} (hw : 1 ≤ w) :
    gradingWN (rank leech) w ≤ T ((w : ℤ) + LatGen.halfNorm leech γ)
\end{leancode}
File \leanpath{VOA/Moonshine/FZGenExcited.lean}.

\begin{dictionary}
\lean{Fk} is the 24-boson Fock space, \lean{gradingWN _ w} its weight-$w$ part; \lean{halfNorm γ} is
$\tfrac12|\gamma|^2$; \lean{quadF γ a i b j k p} is the Fock part of the action of the quadratic state
$h_a(-i-1)h_b(-j-1)\mathbb 1$ on $p\otimes e^\gamma$.  \lean{QuadClosed} says: weight-raising modes ($m\le0$) of
quadratic Heisenberg states on lower charge-$\gamma$ states lie in the target.  The charge is arbitrary,
in particular in the coset $\Leech\setminus\Leech_0$ of $\chXIVB$.
\end{dictionary}

\begin{physics}
Every oscillator excitation of a charged state can be produced by bilinear currents $\partial X\partial X$
acting on lower excitations — the free-boson statement that the $\theta$-even Heisenberg fields generate
all excited states of a charged sector of nonzero momentum.
\end{physics}

\paragraph{Proof on paper.}  The argument of Chapter~\ref{ch:12} for $\chXIVA$ (excited spanning), which used
its target only through this closure and the charge only through $|\gamma|^2\ge4$ even: induction on the
number of oscillators, with a degree-one base case solved by a small linear-algebra system in the
charge coordinates. \qed

\begin{theorem}[The pair lemma]\label{ch14:thm:pair}
(a) If $x\in\Z^{24}$ has all entries odd, $x\cdot x\ge48$ and $16\mid x\cdot x$, there are $i\ne k$ with
$|x_i|+|x_k|=6$ or $\ge10$.  (b) For $\nu\in\Leech\setminus\Leech_0$ with $|\nu|^2\ge6$ there is $\gamma\in\Leech_0$
of norm $4$ with $\langle\gamma,\nu\rangle=3$ or $\ge5$.
\end{theorem}
\begin{leancode}
theorem exists_good_pair (x : Fin 24 → ℤ) (hodd : ∀ i, Odd (x i)) (h48 : 48 ≤ x ⬝ᵥ x)
    (h16 : (16 : ℤ) ∣ x ⬝ᵥ x) :
    ∃ i k : Fin 24, i ≠ k ∧ (|x i| + |x k| = 6 ∨ 10 ≤ |x i| + |x k|)

theorem exists_pair_gamma {ν : leech.carrier} (hν : ν ∉ lambda0) (h6 : 6 ≤ leech.form ν ν) :
    ∃ γ ∈ lambda0, leech.form γ γ = 4 ∧ (leech.form γ ν = 3 ∨ 5 ≤ leech.form γ ν)
\end{leancode}
File \leanpath{VOA/Moonshine/FZGenV3.lean}.

\begin{dictionary}
\lean{x ⬝ᵥ x} is the dot product; \lean{leech.form} the Leech inner product.  In (b), $x=$ \lean{lv ν}
(coordinates with $\langle\mu,\nu\rangle=\mu\cdot\nu/8$, all entries odd for $\nu\notin\Leech_0$), and
$\gamma=4(\pm e_i\pm e_k)$; then $\langle\gamma,\nu\rangle=(|x_i|+|x_k|)/2$.  Pure arithmetic; no Golay facts.
\end{dictionary}

\begin{physics}
A momentum-space statement: every high-momentum state of the odd sector can be split as the
product of a norm-4 momentum of the even sector and a lower one, with an OPE singularity of the right
order to produce it (the exponent $\langle\gamma,\nu-\gamma\rangle$ governs $e^\gamma(z)e^{\delta}(w)\sim(z-w)^{\langle\gamma,\delta\rangle}$).
\end{physics}

\paragraph{Proof on paper.}  (a) Suppose every pair has $|x_i|+|x_k|\in\{2,4,8\}$.  If two entries
had magnitude $\ge3$, the pair sum $\in\{6,\ge10\}$ unless they are $\{3,5\}$, and no third value is then
compatible; so at most one entry exceeds $1$: $x=(m,(\pm1)^{23})$ with $m$ odd, $x\cdot x=m^2+23$.  Now $m^2+23\ge48$ forces $|m|\ge5$; $|m|=5$ has the
good pair $5+1=6$, and $|m|\ge9$ has $|m|+1\ge10$; the only remaining value $|m|=7$ gives $x\cdot x=72$,
not divisible by $16$.  Contradiction.  (b) Choose signs $s_i=\operatorname{sgn}x_i$. \qed

\begin{theorem}[$A(\chXIVB)$ is generated by its top]\label{ch14:thm:bimodgen}
(a) For $\nu\in\Leech\setminus\Leech_0$, $|\nu|^2\ge6$, the ground state $x_\nu=e^\nu+e^{-\nu}$ is in the span of
weight-raising products of lower states.  (b) For $n\ge3$, $(\chXIVB)_n\subseteq\operatorname{span}\lean{lowerProdM}\,n$.
(c) Every subspace of $\chXIVB$ containing $O(\chXIVB)$ and $(\chXIVB)_2$ and closed under the left and right
actions of homogeneous $a\in\chXIVA$ is $\chXIVB$; i.e.\ \lean{BimodGenV3} holds.
\end{theorem}
\begin{leancode}
theorem ground_mem_span {ν : leech.carrier} (hν : ν ∉ lambda0) (h6 : 6 ≤ leech.form ν ν) :
    frameX0P ν ∈ span ℂ (Lv (LatGen.halfNorm leech ν))

theorem V3_lowerSpanned (n : ℤ) (hn : 3 ≤ n) :
    MV3.weight (n : ℂ) ≤ span ℂ (lowerProdM MV3 n)

theorem bimodGen_V3 (S : Submodule ℂ (frameSector 1))
    (hO : ∀ (p : ℕ) (a : V0Sub.carrier), a ∈ V0P.grading (p : ℤ) → ∀ x, fzCirc MV3 p a x ∈ S)
    (hL : ∀ (p : ℕ) (a : V0Sub.carrier), a ∈ V0P.grading (p : ℤ) → ∀ x ∈ S, fzLeft MV3 p a x ∈ S)
    (hR : ∀ (p : ℕ) (a : V0Sub.carrier), 1 ≤ p → a ∈ V0P.grading (p : ℤ) → ∀ x ∈ S,
      fzRight MV3 p a x ∈ S)
    (htop : ∀ x : frameSector 1, x ∈ MV3.weight (2 : ℂ) → x ∈ S) : S = ⊤
\end{leancode}
File \leanpath{VOA/Moonshine/FZGenV3.lean}.
\begin{leancode}
theorem fzGenerated_MV3 : FZ.FZGenerated MV3 (2 : ℂ)
\end{leancode}
File \leanpath{VOA/Moonshine/FZGenBridge.lean}.

\begin{dictionary}
\lean{Lv n} is \lean{lowerProdM MV3 n} viewed inside $V_\Leech^+$.  \lean{bimodGen_V3} is stated with the local
operations; \lean{fzGenerated_MV3} transfers it to the vocabulary of \lean{FZBimodule.lean} (the operations agree
on homogeneous $a$, and \lean{fzO} contains the $\operatorname{Res}_0$), so it is literally \lean{BimodGenV3}.
No hypotheses: unconditional.
\end{dictionary}

\begin{physics}
Every field of the odd frame sector $\chXIVB$ is, modulo descendants, generated from its
$49152$ lowest-weight fields by fusion with the even chiral algebra $\chXIVA$ from both sides.
\end{physics}

\paragraph{Proof on paper.}  $(\chXIVB)_n$ is spanned by projected word states $P(p\otimes e^\nu)$,
$\nu\in\Leech\setminus\Leech_0$.  Excited ones ($p$ of positive Fock weight) are handled by
Theorem~\ref{ch14:thm:excited}, with the target ``Fock states whose projection is a combination of
weight-raising products''.  Ground ones: pick $\gamma$ by the pair lemma, $\delta=\nu-\gamma$,
$f=\langle\gamma,\delta\rangle=\langle\gamma,\nu\rangle-4$.  If $f=-1$, $(x_\gamma)_{(0)}x_\delta=\pm x_\nu$.  If $f=k\ge1$,
$(x_\gamma)_{(-k-1)}x_\delta=\pm(x_\nu+2P(S_{2k}(\gamma)e^{\gamma-\delta}))$, the second term excited.  Both are
weight-raising products of lower states.  Then the criterion (Theorem~\ref{ch14:thm:crit}) with $G=(\chXIVB)_2$,
$n_0=2$. \qed

\begin{status}
\lean{bimodGen_V3}, \lean{fzGenerated_MV3}: unconditional.  \texttt{\#print axioms} \lean{VOA.Moonshine.FZGen.fzGenerated_MV3} (re-run): \lean{[propext, Classical.choice, Quot.sound]}.
\end{status}

\section{The final finite identity}\label{ch14:sec:toppairs}

\subsection{\texorpdfstring{Everything reduced to \lean{TopPairs}}{Everything reduced to TopPairs}}

\begin{theorem}[$\Vnat$ from one finite identity]\label{ch14:thm:maintop}
For $\varepsilon^2=1$: if \lean{TopPairs} holds for $D(c_\star,\varepsilon)$, where $c_\star^2\kappa_0=1$ is the
fixed normalisation, then all targets of the final theorem hold.
\end{theorem}
\begin{leancode}
theorem moonshine_of_topPairs {ε : ℂ} (hε : ε * ε = 1)
    (htop : FZKey.TopPairs (mkData cStar0 ε (c_ne_of_norm cStar0_norm) hε)) :
    ∃ V : VertexOperatorAlgebra Vnat,
      V.Y = Ynat VnatAxioms.wbarL ∧ V.vacuum = vacNat ∧ V.conformalVector = omegaNat ∧
      V.centralCharge = 24 ∧
      V.grading 1 = ⊥ ∧
      Module.finrank ℂ (V.grading 2) = 196884 ∧
      (∀ n : ℕ, (Module.finrank ℂ (V.grading (n : ℤ)) : ℂ) = PowerSeries.coeff n JShift) ∧
      V.IsSimple
\end{leancode}
File \leanpath{VOA/Moonshine/MainTopPairs.lean}.

\begin{dictionary}
\lean{cStar0} is a fixed choice of $c$ with \lean{cStar0 * cStar0 * κ0 = 1} (\lean{cStar0_norm}).  The proof
term is a one-line composition: \lean{moonshine_of_keySkew cStar0_norm hε (FZKey.keySkew_of_FZ _ tF_mem tF_ne
FZGen.fzGenerated_MV3 htop)} — i.e.\ Theorems~\ref{ch14:thm:isomk}, \ref{ch14:thm:fzred},
\ref{ch14:thm:bimodgen}.  The only hypotheses are \lean{hε} and \lean{htop}.
\end{dictionary}

\begin{physics}
The entire consistency of the Moonshine orbifold — that twist fields are mutually local — has been
reduced to one finite table of OPE coefficients among the $49152$ lowest states of the odd frame sector
and the lowest twisted states.  In the bootstrap language: crossing symmetry of the full theory follows
from a finite set of primary couplings, given the symmetry structure.
\end{physics}

\begin{status}
Conditional on \lean{htop}; \texttt{\#print axioms} (re-run): \lean{[propext, Classical.choice, Quot.sound]}.
\end{status}

\subsection{Sign arithmetic of two odd norm-4 Leech vectors}

Work in the integer coordinates $n=(n_0,\dots,n_{23})$ (\lean{lv n}) with $\langle\mu,\nu\rangle=\tfrac18\mu\cdot\nu$.
A norm-4 vector $n\in\Leech\setminus\Leech_0$ has one entry $\pm3$, at position $i=j(n)$ (\lean{blk n}), and
twenty-three entries $\pm1$.  Define for odd $x$: $\chi(x)=1$ if $x\equiv1\pmod4$, else $-1$; and
$s(x)=\operatorname{sgn}x$.

\begin{definition}[Characters and the set $N$]\label{ch14:def:NS}
For $n,m\in\Leech$, $N(n,m)=\{k:\ \langle\gamma^-_{0k},n\rangle+\langle\gamma^-_{0k},m\rangle\text{ odd}\}$.
\end{definition}
\begin{leancode}
def chiZ (x : ℤ) : ℂ := if x % 4 = 1 then 1 else -1

def sgnZ (x : ℤ) : ℂ := if 0 < x then 1 else -1

noncomputable def NS (n m : leech.carrier) : Finset (Fin 24) :=
  Finset.univ.filter fun k => ¬ Even (leech.form (gM 0 k) n + leech.form (gM 0 k) m)
\end{leancode}
File \leanpath{VOA/Moonshine/TopPairsArith.lean}.

\begin{dictionary}
\texttt{x \% 4} is Lean's integer remainder, in $\{0,1,2,3\}$ for the modulus $4$ (Euclidean for positive
modulus), so for odd $x$ it is $1$ or $3$; \lean{chiZ} is $\chi$.  \lean{sgnZ} is used only on nonzero
arguments.  \lean{Finset.univ.filter} is $\{k\in\{0..23\}:\ldots\}$.  The theorem \lean{uN_NS} identifies
$N$ as the character of the joint $Z_{0k}$-eigenvector $e_ne_mt_0$.
\end{dictionary}

\begin{physics}
$N$ records which of the commuting $\Z/2$ ``frame'' symmetries $Z_{0k}$ of the twisted ground states act by
$-1$ on the state $e_ne_mt_0$: it is the charge vector of a ground state of the twisted sector under a
maximal abelian subgroup of the extraspecial group $2^{1+24}$.
\end{physics}

\begin{theorem}[The $B$-eigenvalue]\label{ch14:thm:betaN}
For $n,m\in\Leech\setminus\Leech_0$ of norm $4$, with $i=j(n)$, $j=j(m)$,
\begin{equation*}
  \omega^{|N(n,m)|}=(-1)^{\langle m,n\rangle}\,s(n_im_i)\,s(n_jm_j).
\end{equation*}
\end{theorem}
\begin{leancode}
theorem betaN_NS {n m : leech.carrier} (hn : n ∉ lambda0) (hn4 : leech.form n n = 4)
    (hm : m ∉ lambda0) (hm4 : leech.form m m = 4) :
    betaN (NS n m) = (-1 : ℂ) ^ (leech.form m n)
      * (sgnZ (lv n (blk n) * lv m (blk n)) * sgnZ (lv n (blk m) * lv m (blk m)))
\end{leancode}
File \leanpath{VOA/Moonshine/TopPairsArith.lean}.

\begin{dictionary}
\lean{betaN N} $=\prod_i(\omega\text{ if }i\in N\text{ else }1)=\omega^{|N|}$ (\lean{betaN_eq}), $\omega=e^{i\pi/4}$
(\lean{ωB}).  \lean{(-1 : ℂ) ^ (leech.form m n)} is an integer power.  Hypotheses: odd coset and norm 4
for both vectors — exactly the top basis vectors.
\end{dictionary}

\begin{physics}
The operator $B$ is the twisted-sector part of the triality, a ``discrete Fourier transform'' on the
$2^{12}$ ground states; its eigenvalue on a ground state is an eighth root of unity determined by the
charges.  The theorem says this phase is in fact a sign, computable from the momenta $n,m$.
\end{physics}

\paragraph{Proof on paper (counting $|N|$).}  Since $\langle\gamma^-_{0k},n\rangle=(n_0-n_k)/2$ and
$(-1)^{(x-y)/2}=\chi(x)\chi(y)$ for odd $x,y$,
\begin{equation*}
 2|N|=\sum_{k}\bigl(1-\chi(n_0)\chi(m_0)\,\chi(n_k)\chi(m_k)\bigr)=24-\chi(n_0)\chi(m_0)\sum_k\chi(n_k)\chi(m_k).
\end{equation*}
On the four possible values $\pm1,\pm3$ one has the polynomial identity $6\chi(x)=x(7-x^2)$:
$x=1\mapsto6$, $x=-1\mapsto-6$, $x=3\mapsto-6$, $x=-3\mapsto6$, matching $\chi(1)=1,\chi(-1)=-1,\chi(3)=-1,\chi(-3)=1$.
Substituting and using $n_k^2,m_k^2\in\{1,9\}$,
\begin{equation*}
 36\sum_k\chi(n_k)\chi(m_k)=36\,n\cdot m-48(n_im_i+n_jm_j)+64[i=j]\,n_im_i,
\end{equation*}
with $n\cdot m=8\langle n,m\rangle$.  Hence $|N|\equiv0\pmod4$ and $|N|/4$ has the parity of $\langle m,n\rangle$
corrected by the block signs; $\omega^{|N|}=(-1)^{|N|/4}$.  A finite case split over the signs finishes.
No Golay-code fact is needed beyond $\langle n,m\rangle\in\Z$. \qed

\begin{example}
Take $n=m$.  Then $N=\emptyset$ (each $\langle\gamma^-_{0k},n\rangle$ is counted twice), $\omega^0=1$; on the right,
$(-1)^{4}s(9)s(9)=1$. \qed
\end{example}

\begin{theorem}[The coefficient identity]\label{ch14:thm:coef}
For $n,m$ as above and every $l$:
\begin{equation*}
 c_\star(m)\Bigl(\delta_{jl}-\tfrac{n_jn_l}{4}\Bigr)u_l\,\omega^{|N|}
 =c_\star(n)(-1)^{\langle m,n\rangle}\Bigl(\delta_{il}-\tfrac{m_im_l}{4}\Bigr),
\end{equation*}
where $u_l=(-1)^{[l\in N]}$.
\end{theorem}
\begin{leancode}
theorem coef_identity {n m : leech.carrier} (hn : n ∉ lambda0) (hn4 : leech.form n n = 4)
    (hm : m ∉ lambda0) (hm4 : leech.form m m = 4) (l : Fin 24) :
    cStar m * ((if blk m = l then 1 else 0) - ((2 * lv n (blk m) * lv n l : ℤ) : ℂ) / 8)
        * (uN (NS n m) l * betaN (NS n m))
      = cStar n * (-1 : ℂ) ^ (leech.form m n)
        * ((if blk n = l then 1 else 0) - ((2 * lv m (blk n) * lv m l : ℤ) : ℂ) / 8)
\end{leancode}
File \leanpath{VOA/Moonshine/TopPairsArith.lean}.

\begin{dictionary}
$\tfrac{2n_jn_l}{8}=\tfrac{n_jn_l}4$; \lean{uN N l} $=-1$ if $l\in N$ else $1$; \lean{cStar n}
$=i^{n_0}\cdot n_{j(n)}/3$.  Note $\delta_{jl}-n_jn_l/4$ are the coordinates of $w_j-2\langle w_j,n\rangle n$ in the frame basis.
\end{dictionary}

\begin{physics}
The equality of two explicit OPE coefficients: the coupling of $(x_n,x_m)$ to the twisted state
$h_{w_l}(-\tfrac12)\otimes e_ne_mt_0$ in the two orders.  It is the place where the phases of the cocycle,
of the triality on ground states, and of the oscillator part must conspire.
\end{physics}

\paragraph{Proof on paper.}  Insert Theorem~\ref{ch14:thm:betaN}, $u_l=\chi(n_0)\chi(n_l)\chi(m_0)\chi(m_l)$, and
$i^{m_0}\chi(n_0)\chi(m_0)=i^{n_0}$; what remains is a local identity in the coordinates $i,j,l$, checked by
cases on $n_i,m_j\in\{\pm3\}$ and the remaining entries $\pm1$. \qed

\subsection{Both sides on the top basis}

\begin{theorem}[The top-pair identity]\label{ch14:thm:toppairs0}
(a) For $|n|^2=4$: $o(x_n)\bigl(h_v(-\tfrac12)\otimes t\bigr)=2^{-3}h_{v-2\langle v,n\rangle n}(-\tfrac12)\otimes e_nt$.
(b) $e_ne_mt_0$ is a joint $Z_{0k}$-eigenvector with character $N(n,m)$.
(c) For all $v,b\in(\chXIVB)_2$: $\Theta_0\bigl(o(v)\Phi_0b\bigr)=o(b)\Phi_0v$.
\end{theorem}
\begin{leancode}
theorem oX_hP {n : leech.carrier} (hn4 : leech.form n n = 4) (v : Fin (rank leech) → ℂ)
    (t : TLeechT) :
    mode (plusModule.YM (frameX0P n)) 1 (hP v t)
      = ((2 : ℂ) ^ (-(3 : ℤ))) • hP (latL n v) (eT leech leech_even leech_isUnit_det_gramMat n t)

theorem isEig_z (n m : leech.carrier) : IsEigN (NS n m) (zz n m)

theorem top_pairs0 {v b : VPlus} (hv : v ∈ OmegaV3) (hb : b ∈ OmegaV3) :
    Θ0 (mode (plusModule.YM v) 1 (Φ0 b)) = mode (plusModule.YM b) 1 (Φ0 v)
\end{leancode}
File \leanpath{VOA/Moonshine/TopPairsCore.lean}.

\begin{dictionary}
\lean{hP v t} is $h_v(-\tfrac12)\otimes t$ for $v\in\C^{24}$ (orthonormal coordinates); \lean{latL n v} is
$v-2\langle v,n\rangle n$; \lean{eT … n} is $e_n$ on $T$; \lean{zz n m := eL n (eL m t0)}; \lean{IsEigN N z} means
$Z_{0k}z=u_kz$ for $k\ne0$.  (c) is for the \emph{unnormalised} data.
\end{dictionary}

\begin{physics}
(a) is the twisted vertex operator of $e^n+e^{-n}$ acting on the lowest twisted states: it shifts the
ground state by $e_n$ and rotates the oscillator polarisation.  (c) is the finite crossing symmetry
of \lean{TopPairs} for the explicit triality.
\end{physics}

\paragraph{Proof on paper.}  (a) from the explicit twisted operator $Y^T(e^\alpha,z)$ of Chapter~\ref{ch:08} at its
top coefficient.  (b) $Z_{0k}$ commutes with $e_n$ up to the sign $(-1)^{\langle\gamma^-_{0k},n\rangle}$ and fixes
$t_0$ (seed).  (c) On the basis: $\Phi_0x_m=c_\star(m)h_{w_j}\otimes e_mt_0$; by (a), $o(x_n)$ of it is
$2^{-3}c_\star(m)\sum_l(\delta_{jl}-n_jn_l/4)h_{w_l}\otimes e_ne_mt_0$; $\Theta_0$ acts on block $l$ by $\psi_l$, i.e.\ by
$u_l\omega^{|N|}$ (by (b)).  The other side gives $2^{-3}c_\star(n)\sum_l(\delta_{il}-m_im_l/4)h_{w_l}\otimes e_me_nt_0$ and
$e_me_n=(-1)^{\langle m,n\rangle}e_ne_m$.  Coefficients agree by Theorem~\ref{ch14:thm:coef}.  Extend
bilinearly. \qed

\begin{theorem}[\lean{TopPairs} with $\varepsilon=+1$]\label{ch14:thm:topone}
\begin{leancode}
theorem topPairs_one :
    FZKey.TopPairs (mkData cStar0 1 (c_ne_of_norm cStar0_norm) (by norm_num))
\end{leancode}
File \leanpath{VOA/Moonshine/TopPairs.lean}.
\end{theorem}

\begin{dictionary}
\lean{(by norm_num)} is the proof of $1\cdot1=1$.  For $D(c_\star,1)$ both sides of \lean{TopPairs} are $c_\star$
times the two sides of Theorem~\ref{ch14:thm:toppairs0}(c).  Unconditional.
\end{dictionary}

\begin{physics}
The correct ``orientation'' of the triality on the self-paired twisted sector is $\Theta_0$ itself,
with no extra sign.
\end{physics}

\begin{warning}[title={Trap / dead end: The sign $\varepsilon$ and the earlier numerics}]
For $\varepsilon=-1$, \lean{TopPairs} would force $\Theta_0(o(x_n)\Phi_0x_n)=0$, which is false; this is
\emph{analysis only}, not machine-proved — and not needed.  Before the proof, the identity was tested
numerically with the library's own symbolic formulas on all $98304$ odd norm-4 vectors (via the Golay
code): $200{,}000$ random pairs and $64$ full rows, $0$ failures.  The earlier $N(A_1^{24})$ model
had prescribed ``$-\Theta_J$''; the relation of $\Theta_J$ to $\Theta_0$ was never computed, and the Lean proof
does not depend on it.
\end{warning}

\begin{theorem}[The Moonshine module]\label{ch14:thm:moonshine}
There is a vertex operator algebra structure on $\Vnat$ whose vertex operator is the block operator
\eqref{ch14:eq:block}, with vacuum $(\mathbb 1,0)$, conformal vector $(\omega,0)$, central charge $24$,
$\Vnat_1=0$, $\dim\Vnat_n=$ the coefficient of $q^n$ in the formal $q$-expansion $q\,J(q)$ for all $n\ge0$
(so $\dim\Vnat_2=196884$), and which is simple.  These data determine the structure uniquely, so the
theorem constructs a specific object.  The Lean statement and its full fidelity check are in
\S\ref{ch14:sec:decode}.
\end{theorem}
\begin{leancode}
theorem moonshine :
    ∃ V : VertexOperatorAlgebra Vnat,
      V.Y = Ynat VnatAxioms.wbarL ∧ V.vacuum = vacNat ∧ V.conformalVector = omegaNat ∧
      V.centralCharge = 24 ∧
      V.grading 1 = ⊥ ∧
      Module.finrank ℂ (V.grading 2) = 196884 ∧
      (∀ n : ℕ, (Module.finrank ℂ (V.grading (n : ℤ)) : ℂ) = PowerSeries.coeff n JShift) ∧
      V.IsSimple :=
  moonshine_of_topPairs (by norm_num) topPairs_one
\end{leancode}
File \leanpath{VOA/Moonshine/TopPairs.lean}.

\begin{physics}
The chiral algebra of the $c=24$ conformal field theory of Frenkel--Lepowsky--Meurman: a
chiral CFT with a unique vacuum, no currents (no continuous symmetry), $196884$ weight-two states
(by FLM's theorem, the stress tensor plus $196883$ primaries forming the smallest faithful
representation of the Monster), and torus partition function $J(\tau)$.  As a string
compactification it is the internal CFT of the ``Monster string''.  This is the physical picture,
not the content of the Lean statement: the Monster symmetry, the holomorphy of the theory (a single
irreducible module) and the analytic modular invariance of $J(\tau)$ are \emph{not} part of this
theorem.
\end{physics}

\section{The final theorem, read clause by clause}\label{ch14:sec:decode}

This section is the fidelity check of the headline result of the book.  We restate the theorem,
unfold \emph{every} defined term that occurs in it down to standard mathematics, examine each conjunct,
and ask how the statement could have been weaker than intended — vacuous, satisfied by an impostor, or
silently using a different convention — and why it is not.

\subsection{The statement and the objects it mentions}

\begin{leancode}
theorem moonshine :
    ∃ V : VertexOperatorAlgebra Vnat,
      V.Y = Ynat VnatAxioms.wbarL ∧ V.vacuum = vacNat ∧ V.conformalVector = omegaNat ∧
      V.centralCharge = 24 ∧
      V.grading 1 = ⊥ ∧
      Module.finrank ℂ (V.grading 2) = 196884 ∧
      (∀ n : ℕ, (Module.finrank ℂ (V.grading (n : ℤ)) : ℂ) = PowerSeries.coeff n JShift) ∧
      V.IsSimple
\end{leancode}
File \leanpath{VOA/Moonshine/TopPairs.lean} (namespace \lean{VOA.Moonshine}).

\paragraph{Shape.}  The statement is an existential: \emph{there is} a term \lean{V} of type
\lean{VertexOperatorAlgebra Vnat} — a VOA structure on the fixed vector space $\Vnat$ — satisfying eight
conjuncts joined by \lean{∧}.  There are no hypotheses, no section variables, no implicit arguments.
The space $\Vnat$ is fixed in advance; what is asserted to exist is the \emph{structure}, and the first
four conjuncts then pin that structure down completely.  So the existential is not a
loophole: it is the standard way of saying ``the following explicit data form a VOA''.

\paragraph{A construction with a unique witness.}  Every datum of a \lean{VertexOperatorAlgebra} on
$\Vnat$ is fixed by the statement: \lean{Y}, \lean{vacuum} and \lean{conformalVector} by conjuncts
(1)--(3), \lean{centralCharge} by (4); the \lean{grading} is forced, because \lean{isInternal} and
\lean{L0_grading} make \lean{V.grading n} exactly the $n$-eigenspace of the pinned operator
$L_0=\omega_{(1)}$ (argued under (2)--(3) below); and all remaining fields are propositions (the field
\lean{finrank} is the \lean{Prop} \lean{FiniteDimensional}), so two witnesses with the same data are
equal.  Hence any two witnesses of the existential coincide: the theorem cannot be satisfied by ``an
unrelated VOA structure on the same space''.  It is a \emph{construction} of one specific object,
stated in the form ``these explicit data satisfy the VOA axioms and the listed properties''.  (The
release audit suggests, as a cosmetic improvement, exposing the witness as a named definition
\lean{Classical.choose moonshine}; the mathematical content is the same.)  This uniqueness is also what
excludes the recorded trap, the square-zero extension with $\overline W=0$: it is not a witness,
because conjunct (1) pins the twisted$\times$twisted block to $\overline W$, which is proved nonzero
(\lean{leechWbar_ne_zero}), and independently because it is not simple (conjunct (8)); see the
discussion under (8).

\paragraph{The type \lean{VertexOperatorAlgebra}.}  This is the library's definition
(Definitions~\ref{l03:def:VA} and~\ref{l03:def:VOA}), which we reproduce so that the referee need not turn back:
\begin{leancode}
structure PreVertexAlgebra (𝒱 : Type u) [AddCommGroup 𝒱] [Module ℂ 𝒱] where
  /-- The internal `ℤ`-grading `𝒱 = ⊕ₙ 𝒱ₙ`, as a family of submodules of `𝒱`. -/
  grading : ℤ → Submodule ℂ 𝒱
  /-- The grading is a genuine internal direct-sum decomposition of `𝒱`. -/
  isInternal : DirectSum.IsInternal grading
  /-- Each homogeneous piece `𝒱ₙ` is finite-dimensional (Gannon Def. 5.1.3(a)). -/
  finrank : ∀ n, FiniteDimensional ℂ (grading n)
  /-- The state–field correspondence `Y(u,z) = ∑ₙ u₍ₙ₎ z^{-n-1}`, `ℂ`-linear in `u`. -/
  Y : 𝒱 →ₗ[ℂ] (End ℂ 𝒱)⟦z^±⟧
  /-- The vacuum vector `𝟙`. -/
  vacuum : 𝒱
  /-- The vacuum lies in degree `0`: `𝟙 ∈ 𝒱₀`. -/
  vacuum_mem : vacuum ∈ grading 0
  /-- **VA1 (grading).** For `u ∈ 𝒱ₖ` and `v ∈ 𝒱ₗ`, `u₍ₙ₎ v ∈ 𝒱_{k+ℓ-n-1}`. -/
  grading_mode : ∀ {k ℓ : ℤ} (u : 𝒱), u ∈ grading k → ∀ (v : 𝒱), v ∈ grading ℓ →
    ∀ (n : ℤ), mode (Y u) n v ∈ grading (k + ℓ - n - 1)
  /-- **VA2 (vacuum).** `Y(𝟙, z) = id`, the unit field. -/
  vacuum_Y : Y vacuum = (1 : (End ℂ 𝒱)⟦z^±⟧)
  /-- **VA3 (state–field correspondence).** `u₍₋₁₎ 𝟙 = u` and `u₍ₙ₎ 𝟙 = 0` for `n ≥ 0`. -/
  creation : ∀ u : 𝒱, mode (Y u) (-1) vacuum = u ∧ ∀ n : ℤ, 0 ≤ n → mode (Y u) n vacuum = 0
  /-- **VA4 (locality).** `Y(u,z₁)` and `Y(v,z₂)` are mutually local:
  `(z₁ - z₂)ᴹ [Y(u,z₁), Y(v,z₂)] = 0` for some `M = M(u,v)`. -/
  locality : ∀ u v : 𝒱, Local (Y u) (Y v)
  /-- **VA5 (regularity).** For each `u, v` there is `N` with `u₍ₙ₎ v = 0` for `n ≥ N`. -/
  regularity : ∀ u v : 𝒱, ∃ N : ℤ, ∀ n : ℤ, N ≤ n → mode (Y u) n v = 0
\end{leancode}
File \leanpath{VOA/Core/VertexAlgebra.lean}.
\begin{leancode}
structure GradedVertexOperatorAlgebra (𝒱 : Type u) [AddCommGroup 𝒱] [Module ℂ 𝒱]
    extends PreVertexAlgebra 𝒱 where
  /-- The conformal vector `ω`. -/
  conformalVector : 𝒱
  /-- The conformal vector lies in degree `2`: `ω ∈ 𝒱₂`. -/
  conformal_mem : conformalVector ∈ grading 2
  /-- The central charge / rank `c ∈ ℂ`. -/
  centralCharge : ℂ
  /-- **VOA1 (conformal symmetry).** Writing `Lₙ = ω₍ₙ₊₁₎ = mode (Y ω) (n+1)`, the modes
  obey the Virasoro relation `[Lₘ,Lₙ] = (m-n)L_{m+n} + (c/12)(m³-m)δ_{m+n,0}`. -/
  virasoro : ∀ m n : ℤ,
      ⁅mode (Y conformalVector) (m + 1), mode (Y conformalVector) (n + 1)⁆
        = ((m : ℂ) - n) • mode (Y conformalVector) (m + n + 1)
          + (if m + n = 0 then (((m : ℂ) ^ 3 - m) / 12 * centralCharge) • (1 : End ℂ 𝒱)
             else 0)
  /-- **VOA2 (conformal weight).** `L₀ = mode (Y ω) 1` acts on `𝒱ₙ` as multiplication by
  `n`: `L₀ v = n • v` for `v ∈ 𝒱ₙ`. -/
  L0_grading : ∀ {n : ℤ} (v : 𝒱), v ∈ grading n → mode (Y conformalVector) 1 v = (n : ℂ) • v
  /-- **VOA3 (translation generator).** `Y(L₋₁ v, z) = ∂_z Y(v,z)`, where
  `L₋₁ = mode (Y ω) 0`. -/
  translation : ∀ v : 𝒱, Y (mode (Y conformalVector) 0 v) = deriv (Y v)
  /-- **Grading truncation** (FLM/LL): `𝒱ₙ = 0` for all sufficiently negative `n`. -/
  grading_bounded_below : ∃ N : ℤ, ∀ n : ℤ, n < N → grading n = ⊥

structure VertexOperatorAlgebra (𝒱 : Type u) [AddCommGroup 𝒱] [Module ℂ 𝒱]
    extends GradedVertexOperatorAlgebra 𝒱 where
  /-- **VOA4 (CFT type), part 1.** `𝒱₀ = ℂ𝟙`. -/
  vacuum_grade : grading 0 = ℂ ∙ vacuum
  /-- **VOA4 (CFT type), part 2.** `𝒱ₙ = 0` for `n < 0`. -/
  positivity : ∀ n : ℤ, n < 0 → grading n = ⊥
  grading_bounded_below := ⟨0, positivity⟩
\end{leancode}
File \leanpath{VOA/Core/VertexOperatorAlgebra.lean}.

\begin{dictionary}
\textbf{The definition is the textbook one, of CFT type} (Frenkel--Lepowsky--Meurman, Gannon Def.~5.1.3): an internal
$\Z$-graded space with finite-dimensional pieces (\lean{DirectSum.IsInternal grading}: the canonical map
$\bigoplus_n\mathcal V_n\to\mathcal V$ is an isomorphism), a linear state--field map
$Y:\mathcal V\to\End(\mathcal V)[[z^{\pm1}]]$, a vacuum, grading compatibility, vacuum and creation axioms,
locality, truncation, a conformal vector of weight $2$ whose modes $L_n=\omega_{(n+1)}$ satisfy Virasoro with
central charge $c$, $L_0$ acting by the grading, $L_{-1}$-translation, and CFT type
($\mathcal V_0=\C\mathbb 1$, $\mathcal V_{n<0}=0$).  The two blocks above are the library's two layers:
\lean{GradedVertexOperatorAlgebra} is the general FLM/Lepowsky--Li notion (grading truncated from below), and
\lean{VertexOperatorAlgebra} adds CFT type, its truncation field being filled automatically by $N=0$; a
\lean{VertexOperatorAlgebra} is CFT type by definition.  Locality plus the other axioms imply the Jacobi
identity (Chapter~\ref{ch:L03}).  It is a demanding definition: the library has a recorded trap where an
over-quantified axiom made a related structure (\lean{TwistedModule}) unsatisfiable
(Chapter~\ref{ch:06}); the \emph{existence} of the concrete VOAs of Chapters~\ref{ch:03}--\ref{ch:04}, and
here of $\Vnat$, shows this definition is satisfiable and non-degenerate.  Note that CFT type and
finite-dimensionality are \emph{part of} the definition, so they are asserted for $\Vnat$ too.
\end{dictionary}

\begin{physics}
These are the axioms of the chiral algebra of a 2d CFT: states $\leftrightarrow$ fields, a unique
$SL_2$-invariant vacuum, mutually local fields, a stress tensor $T(z)=Y(\omega,z)=\sum L_nz^{-n-2}$ with
central charge $c$, and $L_0$ = energy.  CFT type means a unique vacuum and positive energies.
\end{physics}

\subsection{Conjunct by conjunct}

\paragraph{(1) \lean{V.Y = Ynat VnatAxioms.wbarL}: the vertex map is the genuine block operator.}
\begin{leancode}
abbrev Vnat : Type := SkewVP × TP
\end{leancode}
File \leanpath{VOA/Moonshine/VnatCarrier.lean}.
\begin{leancode}
noncomputable def Ynat (Wbar : TP →ₗ[ℂ] FormalDist (TP →ₗ[ℂ] SkewVP)) :
    Vnat →ₗ[ℂ] FormalDist (End ℂ Vnat) where
  toFun x := blockF (leechVOAPlus.Y x.1) (Wbar x.2) (leechWskLin x.2) (plusModule.YM x.1)
\end{leancode}
File \leanpath{VOA/Moonshine/VnatFields.lean}.
\begin{leancode}
noncomputable def wbarL : TP →ₗ[ℂ] FormalDist (TP →ₗ[ℂ] SkewVP) where
  toFun := leechWbar
\end{leancode}
File \leanpath{VOA/Moonshine/VnatAxioms.lean}.
\begin{leancode}
noncomputable def leechWbar (t : TP) : FormalDist (TP →ₗ[ℂ] SkewVP) :=
  Wbar leechBV plusModule TPForm TPForm_weightOrth t
\end{leancode}
File \leanpath{VOA/Moonshine/WbarLeech.lean}.
\begin{leancode}
noncomputable def leechWsk (t : TP) : FormalDist (SkewVP →ₗ[ℂ] TP) :=
  Wsk plusModule.toWeakModule t
\end{leancode}
File \leanpath{VOA/Moonshine/SkewLeech.lean}.

\begin{dictionary}
\textbf{Decoding.}  For $x=(u,t)\in\Vnat$, \lean{Ynat wbarL x} is the $2\times2$ block field
\eqref{ch14:eq:block}: top-left $Y(u,z)$, the vertex operator of the VOA $V_\Leech^+$ (\lean{leechVOAPlus},
Chapter~\ref{ch:04}); bottom-right $Y^T(u,z)$, the module vertex operator of $V_\Leech^+$ on
$(V_\Leech^T)^+$ (\lean{plusModule}, Chapter~\ref{ch:08}); bottom-left $W(t,z)$, the skew intertwiner
$W(t,z)\varphi=e^{zL_{-1}}Y^T(\varphi,-z)t$ (\lean{Wsk}, Chapter~\ref{ch:09}); top-right $\overline W(t,z)$, the
adjoint of $W$ through the canonical forms \lean{leechBV} and \lean{TPForm} (\lean{Wbar}, Chapter~\ref{ch:09}).
\lean{wbarL} is merely \lean{leechWbar} packaged as a linear map.

\textbf{Why this conjunct is essential.}  Without it, the theorem would assert the existence of
\emph{some} VOA on the vector space $\Vnat$ with the listed numerical properties, and the space
$\Vnat$ would be playing no role beyond its graded dimension.  With it, the vertex operator is
\emph{the} FLM/DGM vertex operator on $V_\Leech^+\oplus(V_\Leech^T)^+$: every field is the explicit
untwisted, twisted, or intertwining field built in Chapters~\ref{ch:03}--\ref{ch:09}.

\textbf{The adjoint is not zero.}  The theorem
\begin{leancode}
theorem leechWbar_ne_zero : leechWbar ≠ 0
\end{leancode}
File \leanpath{VOA/Moonshine/VnatGenuine.lean}.
shows that the top-right block is genuinely present: for a nonzero weight-2 twisted $t$ there is $t'$ with
$\overline W(t)_{(3)}t'=c\,\mathbb 1$, $c\ne0$.
\end{dictionary}

\begin{physics}
The OPE of two twist fields contains the identity operator: $\sigma_t(z)\sigma_{t'}(w)\sim
\langle t,t'\rangle(z-w)^{-4}\mathbb 1+\cdots$ for weight-2 twist fields.  This is what
$\overline W\neq0$ expresses, and it is what makes the twisted sector part of one theory rather than a module
glued on with no interaction.
\end{physics}

\paragraph{(2)--(3) \lean{V.vacuum = vacNat}, \lean{V.conformalVector = omegaNat}.}
\begin{leancode}
noncomputable def vacNat : Vnat := (leechVOAPlus.vacuum, 0)

noncomputable def omegaNat : Vnat := (leechVOAPlus.conformalVector, 0)
\end{leancode}
File \leanpath{VOA/Moonshine/VnatCarrier.lean}.

\begin{dictionary}
Vacuum and conformal vector are those of $V_\Leech^+$, i.e.\ of the Leech lattice VOA: $\mathbb 1=e^0$ and
$\omega=\tfrac12\sum_{a=1}^{24}h_a(-1)^2\mathbb 1$ (Chapter~\ref{ch:03}).

\textbf{The grading is determined by (1)--(3).}  The theorem does not say \lean{V.grading = gradingNat}, but
it need not: by \lean{L0_grading} each \lean{V.grading n} is contained in the $n$-eigenspace of
$L_0=\omega_{(1)}$, and because the grading is an internal direct sum, it \emph{equals} that eigenspace
(decompose an eigenvector into homogeneous components; each component is an eigenvector of its own
weight, so only the weight-$n$ component survives).  Since $Y$ and $\omega$ are pinned, $L_0$ is the explicit
operator $L_0\oplus L_0^T$, and \lean{V.grading n} is the explicit space $(V_\Leech^+)_n\times(T^+)_n$
(the library's \lean{gradingNat n}).  Thus the dimension statements (5)--(7) refer to the intended spaces,
not to an arbitrary regrading.  Likewise \lean{V.L n} are the explicit Virasoro operators.
\end{dictionary}

\begin{physics}
The stress tensor of the orbifold is the stress tensor of the 24 free bosons, $T=-\tfrac12\sum:\partial X^a\partial X^a:$,
which is $\theta$-invariant; the orbifold does not change the conformal structure.
\end{physics}

\paragraph{(4) \lean{V.centralCharge = 24}.}  With $\omega$ pinned, $c$ is determined by
$[L_2,L_{-2}]=4L_0+\tfrac c2$; a wrong value would be refutable.  It is the rank of the Leech lattice.

\begin{physics}
$c=24$: 24 free bosons; the critical dimension of the bosonic string after light-cone gauge.
Holomorphic $c=24$ CFTs are exactly those whose characters are (essentially) modular functions,
$Z(\tau)=J(\tau)+\text{const}$.
\end{physics}

\paragraph{(5) \lean{V.grading 1 = ⊥}.}  $\Vnat_1=0$: \lean{⊥} is the zero subspace.  This holds because
$(V_\Leech^+)_1=0$ ($\theta$ kills the 24 currents $h_a(-1)\mathbb 1$, and $\Leech$ has no roots) and $T^+$ starts in
weight $2$.

\begin{physics}
No weight-one states means no conserved currents: the Moonshine CFT has no continuous symmetry —
its symmetry group is finite (by FLM's theorem it is the Monster; $\Aut\Vnat=\mathbb M$ is not formalized).  This distinguishes it from the Leech lattice theory, whose 24
currents generate a $U(1)^{24}$.
\end{physics}

\paragraph{(6) \lean{Module.finrank ℂ (V.grading 2) = 196884}.}  \lean{Module.finrank ℂ S} is the dimension
of the subspace $S$ (a natural number; it would be $0$ for an infinite-dimensional space, but every
\lean{V.grading n} is finite-dimensional by the axiom \lean{finrank}).  $196884=98580+98304$
($(V_\Leech^+)_2$ plus $(T^+)_2=24\cdot4096$).

\begin{physics}
$196884=1+196883$: the stress tensor and the smallest nontrivial representation of the Monster
(a fact from the literature; no Monster action is formalized);
McKay's observation that started Monstrous Moonshine.  The weight-two space carries the
commutative non-associative Griess algebra via $a\cdot b=a_{(1)}b$.
\end{physics}

\paragraph{(7) the full graded dimension.}
\begin{leancode}
noncomputable def jShift : PowerSeries ℂ :=
  E4Series ^ 3 * ((deltaOverQUnit⁻¹ : (PowerSeries ℂ)ˣ) : PowerSeries ℂ)

noncomputable def JShift : PowerSeries ℂ := jShift - (744 : ℂ) • PowerSeries.X
\end{leancode}
File \leanpath{VOA/Modularity/JFunction.lean}.

\begin{dictionary}
\lean{E4Series} is the $q$-expansion $1+240\sum\sigma_3(n)q^n$ of $E_4$; \lean{deltaOverQUnit} is $\Delta/q=\prod(1-q^n)^{24}$,
invertible as a power series (constant term $1$).  So \lean{jShift} $=E_4^3q/\Delta$ is the formal
$q$-expansion of $q\,j$, and \lean{JShift} $=q(j-744)=q\,J$ is the formal $q$-expansion of $qJ$:
\begin{equation*}
 \lean{JShift}=1+0\cdot q+196884\,q^2+21493760\,q^3+\cdots .
\end{equation*}
\lean{PowerSeries.coeff n} is the coefficient of $q^n$; the cast to $\C$ compares a natural-number dimension with
a complex coefficient.  The conjunct says $\sum_n\dim\Vnat_n\,q^n=qJ(q)$, i.e.\ $\operatorname{ch}\Vnat=q^{-1}\sum\dim\Vnat_nq^n=J$.
It is \emph{independent} of (5) and (6) in principle, consistent with them ($[q^1]=0$, $[q^2]=196884$),
and was proved in Chapter~\ref{ch:05} from $\Theta_\Leech=E_4^3-720\Delta$ and the twisted-sector character.
The quantifier is over all $n\in\N$; negative $n$ are covered by \lean{positivity}.

\textbf{Formal, not analytic.}  \lean{JShift} is a formal power series, defined from divisor sums; it is
not a function on the upper half plane.  That it is the $q$-expansion of Klein's analytic
$j(\tau)-744$ is classical (and follows on paper from Mathlib's $q$-expansions of $E_4$ and $\Delta$,
Chapter~\ref{ch:05}), but no Lean theorem identifies \lean{JShift} with the analytic function, and no
modular-invariance statement about the character of $\Vnat$ is formalized.  The conjunct is therefore
``graded dimension $=J$'' in the standard sense of formal $q$-series, which is how the statement of
FLM is usually read.
\end{dictionary}

\begin{physics}
The torus partition function of the chiral CFT, $Z(\tau)=\tr_{\Vnat}q^{L_0-c/24}$, equals $J(\tau)$: it is
modular invariant, as a holomorphic $c=24$ theory must be, with no constant term (no currents).
(Heuristic: the Lean conjunct equates formal $q$-series coefficients; neither the analytic modular
invariance nor holomorphy of $\Vnat$ is formalized.)
\end{physics}

\paragraph{(8) \lean{V.IsSimple}.}
\begin{leancode}
def IsIdeal (W : VertexOperatorAlgebra 𝒱) (I : Submodule ℂ 𝒱) : Prop :=
  ∀ (u v : 𝒱) (n : ℤ), v ∈ I → W.smul u n v ∈ I

def IsSimple (W : VertexOperatorAlgebra 𝒱) : Prop :=
  (⊤ : Submodule ℂ 𝒱) ≠ ⊥ ∧ ∀ I : Submodule ℂ 𝒱, W.IsIdeal I → I = ⊥ ∨ I = ⊤
\end{leancode}
File \leanpath{VOA/Core/BilinearForm.lean}.

\begin{dictionary}
An ideal is a subspace stable under all modes $u_{(n)}$ (a left ideal; for VOAs $L_{-1}$-stable left ideals
are two-sided by skew-symmetry, and $L_{-1}=\omega_{(0)}$ is among the modes).  Simple: nonzero
(\lean{⊤ ≠ ⊥}), and the only ideals are $0$ and everything.  Since \emph{all} left ideals are excluded, not
only two-sided ones, this notion is at least as strong as the textbook notion of a simple VOA.

\textbf{Why this conjunct cannot be dropped — the square-zero impostor.}  Recorded trap (Section~\ref{ch15:sec:trap-sqzero},
Chapter~\ref{ch:10}): replacing $\overline W$ by $0$ in \eqref{ch14:eq:block} gives the square-zero
extension $V_\Leech^+\ltimes(V_\Leech^T)^+$.  It \emph{is} a VOA, with the same vacuum, conformal vector,
$c=24$, $V_1=0$, $\dim V_2=196884$ and graded dimension $J$ — it satisfies conjuncts (2)--(7).  It is
\emph{not} $\Vnat$, and it is excluded twice: (i) by conjunct (1), since \lean{leechWbar ≠ 0}; (ii) by simplicity,
since $0\oplus(V_\Leech^T)^+$ is a proper nonzero ideal there (untwisted fields preserve it, twisted fields
send untwisted states into it and kill it).  Either exclusion alone suffices; the theorem has both.
In the library simplicity was proved via nondegeneracy of an invariant form on $\Vnat$ that pairs the twisted
sector with itself through $\overline W$ — exactly what fails for the impostor.
\end{dictionary}

\begin{physics}
Simplicity is the algebraic form of ``a unique vacuum and a non-degenerate two-point function'': no
decoupled sector, every state has a nonzero overlap with some other state.  A unitary CFT with a
unique vacuum has a simple chiral algebra.
\end{physics}

\subsection{\texorpdfstring{How the constructed object relates to FLM's $\Vnat$}{How the constructed object relates to FLM's V-natural}}

Conjuncts (1)--(4) fix the data $(\Vnat,Y,\mathbb 1,\omega,c)$ completely: the carrier is FLM's space
$V_\Leech^+\oplus(V_\Leech^T)^+$, and the vertex operator is the explicit block operator
\eqref{ch14:eq:block}.  Given those, (5)--(8) are \emph{properties} of a fixed object, not choices; they are
what the literature states about $\Vnat$.  What the Lean theorem does \emph{not} contain is a comparison
with FLM's own vertex operator algebra, and the reader should be clear about where the two presentations
differ.
\begin{itemize}
\item The carrier, the untwisted block $Y$, the twisted-module block $Y^T$ and the intertwiner block $W$
  are FLM's (up to the conventions recorded in Chapters~\ref{ch:03}--\ref{ch:09}).
\item The twisted$\times$twisted block $\overline W$ is \emph{defined} as the adjoint of the skew intertwiner
  $W$ through the invariant forms \lean{leechBV} and \lean{TPForm} (Chapter~\ref{ch:09}).  This is the
  presentation of Huang and of Dolan--Goddard--Montague, \emph{not} FLM's explicit formula for the twisted
  vertex operators on $T\times T$; likewise the locality relation (4.23) is proved here by the
  Frenkel--Zhu / triality reduction of this chapter, not by FLM's own route.
\item No Lean theorem states that the constructed VOA is isomorphic to FLM's $\Vnat$.  On paper the
  identification is standard: $(V_\Leech^T)^+$ is a simple current of $V_\Leech^+$, and a simple-current
  extension of $V_\Leech^+$ by it carries at most one VOA structure up to isomorphism (Dong--Mason, Prop.~5.3; see the
  notes), so any two such structures — FLM's and the one constructed here — are isomorphic.  That
  uniqueness theorem is not formalized.
\end{itemize}
One normalisation subtlety remains: the library's $\overline W$ is fixed by the canonical forms
(\lean{leechBV} normalised by $(\mathbb 1|\mathbb 1)=1$, and \lean{TPForm}).  Rescaling the twisted sector by
$\lambda$ rescales the twisted$\times$twisted block by a nonzero scalar and gives an isomorphic VOA, so
normalisations of this kind do not affect the isomorphism class.  It is therefore fair to call the
constructed object ``the Moonshine module $\Vnat$ (the FLM Leech-lattice orbifold)''; it is \emph{not} fair
to say that FLM's construction is formalized verbatim, or that the isomorphism with FLM's $\Vnat$ is proved.

\subsection{What the theorem does not say}

\begin{itemize}
\item \textbf{No uniqueness.}  It does not say that $\Vnat$ is the unique holomorphic $c=24$ VOA with
  $V_1=0$ (the FLM uniqueness conjecture; proved under additional hypotheses, e.g.\ by
  Dong--Griess--Lam, see the notes).  It does not say any VOA with these numerical properties is
  isomorphic to $\Vnat$ — the square-zero impostor shows the numerical properties alone do not suffice.
\item \textbf{No Monster.}  It does not construct $\Aut\Vnat$ or identify it with the Monster group; the
  Lean library contains no Monster group.  The triality $\chXIVsig$ used in the proof is
  \emph{not} asserted to be an automorphism in the final statement (it is one, by the proof, but no Lean
  theorem states it).
\item \textbf{No McKay--Thompson series}, no Borcherds' Monstrous Moonshine theorem, no modular
  invariance of trace functions (only the graded dimension), no rationality or $C_2$-cofiniteness of $\Vnat$.
\item \textbf{No module category}: holomorphicity (that $\Vnat$ is its only irreducible module) is not
  stated.
\item \textbf{No analytic $j$ in the statement.}  The headline's $J$ is the formal $q$-series \lean{JShift}; that it is
  the $q$-expansion of the analytic $j(\tau)-744$ on the upper half plane is a separate Lean theorem
  (\lean{Moonshine.qExpansion_qJ}, \lean{Moonshine.hasSum_JShift}), not a clause of \lean{moonshine}.
\item \textbf{No isomorphism with FLM's explicit $\Vnat$.}  See the previous subsection: the
  twisted$\times$twisted block is the form-adjoint $\overline W$, and the identification with FLM's own
  vertex operators rests on the (unformalized) uniqueness of simple-current extensions.
\end{itemize}
Accordingly, none of the following may be said on the strength of this theorem: ``$\Vnat$ is characterised
/ unique'', ``$\Aut(\Vnat)=\mathbb M$'', ``$\Vnat$ is holomorphic'', or ``FLM's construction is formalized
verbatim''.

\subsection{Axioms and the trust model}

\begin{plaincode}
'VOA.Moonshine.moonshine' depends on axioms: [propext, Classical.choice, Quot.sound]
\end{plaincode}
This was run for this chapter at an earlier commit (when the chapter was first written; not re-run for this edition) with
\lean{lake env lean -DmaxSynthPendingDepth=3} on a file importing \lean{VOA.Moonshine.TopPairs}
(the flag is required: a bare \lean{lake env lean} can silently drop it).  The three axioms are Lean's
standard ones (propositional extensionality, choice, quotient soundness).  \lean{VOA.pbw} (the PBW theorem of
Chapter~\ref{ch:L02}; formerly the library's only project axiom, now a theorem) does \emph{not} appear: the Moonshine construction does not use it.
\texttt{\#print axioms} does not report hypotheses — but \lean{moonshine} has none.  So the theorem is
\emph{unconditional}.  The project's build record for that earlier commit reports an aggregate
\lean{lake build} of $10343$ jobs with exit $0$ and \lean{scripts/check_repository.py} passing; the full
build was not re-run for this chapter.  One upstream definition was changed during the Moonshine work:
the repair of \lean{TwistedModule.jacobi} (Chapter~\ref{ch:06}); the core definitions of
VOA, module, $J$ used in this statement are otherwise unchanged in content.  (For the release the vertex-algebra
structures were renamed and layered --- \lean{PreVertexAlgebra}, and \lean{VertexOperatorAlgebra} extending
\lean{GradedVertexOperatorAlgebra} --- without changing the fields a VOA must supply; the truncation field of the new
middle layer is filled automatically from \lean{positivity}.)

\paragraph{Release audit 2026-10-27: \texttt{RELEASE\_AUDIT\_MOONSHINE.md}.}  An independent release audit
(\texttt{Notes/Textbook/referee/RELEASE\_AUDIT\_MOONSHINE.md}, at an earlier commit, not re-run for this edition) re-ran
\texttt{\#print axioms} for \lean{VOA.Moonshine.moonshine} and \lean{moonshine_of_topPairs}
(result: the three standard axioms), confirmed \lean{leechWbar_ne_zero} and two non-vacuity examples
(the witness's grading pieces are $L_0$-eigenspaces of the pinned vertex map; the twisted$\times$twisted
block of the pinned map is nonzero), and scanned the $619$ modules of the import closure for
\lean{axiom}, \lean{sorry}, \lean{native_decide} and similar escape hatches (none outside comments; the
closure does not contain \leanpath{VOA/Foundations/AffineLie.lean}, where \lean{VOA.pbw} lives).  Its
verdict is ``announceable with caveats'': no mathematical or formal defect, and caveats on wording only
— the unique-witness reading of the existential, $J$ as a formal $q$-series, the informal identification
with FLM's explicit $\Vnat$, and the absence of the Monster, holomorphy and uniqueness.  This chapter has
been worded accordingly.  The audit did not re-run the aggregate \lean{lake build}; the release commit
still requires one.

\subsection{The dependency chain through the book}

\begin{center}
\begin{tabular}{p{0.35\textwidth}p{0.55\textwidth}}
\toprule
ingredient & where \\ \midrule
VOA axioms, locality, forms & Chapters~\ref{ch:L03}--\ref{ch:L04}\\
Leech lattice, Golay code & Chapter~\ref{ch:02}\\
$V_\Leech$ (\lean{leechVOA}), $c=24$ & Chapter~\ref{ch:03}\\
$\theta$, $V_\Leech^+$, $(V_\Leech^+)_1=0$, \lean{leechBV} & Chapter~\ref{ch:04}\\
graded dimension $=J$ & Chapter~\ref{ch:05}\\
twisted modules (repair), $T$ & Chapters~\ref{ch:06}--\ref{ch:07}\\
$V_\Leech^T$, $(V_\Leech^T)^+$ (\lean{plusModule}) & Chapter~\ref{ch:08}\\
$W$, $\overline W$, \lean{TPForm} & Chapter~\ref{ch:09}\\
block VOA modulo (4.23); simplicity & Chapter~\ref{ch:10}\\
frame, $\sigma$ on $\chXIVA$ & Chapter~\ref{ch:11}\\
twisted triality: Zhu generation, existence & Chapters~\ref{ch:12}--\ref{ch:13}\\
(4.23) via conjugation, FZ, \lean{TopPairs} & this chapter\\
\bottomrule
\end{tabular}
\end{center}

\begin{status}
\lean{VOA.Moonshine.moonshine}: \textbf{unconditional} (no hypotheses); a construction with a unique
witness; \texttt{\#print axioms} $=$
\lean{[propext, Classical.choice, Quot.sound]} (re-run for this chapter, and again by the release audit of
2026-10-27).  Not formalized: $\Aut\Vnat=\mathbb M$, holomorphy, uniqueness of $\Vnat$, the analytic $j$, and
the isomorphism with FLM's explicit $\Vnat$.
\end{status}

\section{Referee's checklist}\label{ch14:sec:checklist}

\begin{longtable}{p{0.27\textwidth}p{0.33\textwidth}p{0.2\textwidth}p{0.1\textwidth}}
\toprule
declaration & intended meaning & hypotheses & check \\ \midrule
\endhead
\lean{modeSpan_x0T} & one twisted vector generates $T^+$ over $V_\Leech^+$ & $t_0\ne0$ & Thm~\ref{ch14:thm:modespan}\\
\lean{SigmaData} & triality data $\Phi,\Psi,\Theta$ & — (structure) & Def~\ref{ch14:def:sigmadata}\\
\lean{sigNat_sq}, \lean{conjF_YnatU_V0} & $\chXIVsig$ involution, $\chXIVA$-compatible & $D$ & Thm~\ref{ch14:thm:signat}\\
\lean{eq_zero_of_V0local} & uniqueness with $\chXIVA$-locality only & cyclic data & Thm~\ref{ch14:thm:V0local}\\
\lean{moonshine_of_sigma} & $\Vnat$ from \lean{SigmaKey} & \lean{SigmaKey} & Thm~\ref{ch14:thm:tpblock}\\
\lean{aut_isometry} & automorphisms preserve the invariant form & CFT type, $L_1V_1=0$ & Thm~\ref{ch14:thm:autiso}\\
\lean{sigmaKey_of_top} & \lean{SigmaKey} from finite isometries and \lean{KeySkew} & three Props & Thm~\ref{ch14:thm:top}\\
\lean{mkData} & explicit triality data & $c\ne0$, $\varepsilon^2=1$ & Def~\ref{ch14:def:mkData}\\
\lean{Φ0_simil}, \lean{κ0_ne} & $\Phi_0$ nonzero similitude & none & Thm~\ref{ch14:thm:simil}\\
\lean{Θ0_Θ0}, \lean{Θ0_isom} & $\Theta_0$ involutive isometry ($\mu=+1$) & none & Thm~\ref{ch14:thm:psi}\\
\lean{moonshine_of_keySkew0} & $\Vnat$ from \lean{KeySkew0} for one sign & \lean{KeySkew0} & Def~\ref{ch14:def:keyskew0}\\
\lean{Dq_eq_zero_of_pair} & Li's lemma for the defect & cyclic $x'$ & Thm~\ref{ch14:thm:li}\\
\lean{keySkew_of_keyTopV} & \lean{KeySkew} from top outputs & \lean{KeyTopV} & Thm~\ref{ch14:thm:keytopv}\\
\lean{keyTopAll_of_FZ} & FZ reduction & \lean{BimodGenV3}, \lean{TopPairs} & Thm~\ref{ch14:thm:fzred}\\
\lean{topPairs_of_keyTopAll} & \lean{TopPairs} necessary & \lean{KeyTopAll} & Thm~\ref{ch14:thm:fzred}\\
\lean{fzGenerated_MV3} & $A(\chXIVB)$ generated by top & none & Thm~\ref{ch14:thm:bimodgen}\\
\lean{betaN_NS}, \lean{coef_identity} & sign arithmetic & odd norm-4 & Thms~\ref{ch14:thm:betaN}, \ref{ch14:thm:coef}\\
\lean{top_pairs0} & top-pair identity for $\Theta_0,\Phi_0$ & none & Thm~\ref{ch14:thm:toppairs0}\\
\lean{topPairs_one} & \lean{TopPairs} with $\varepsilon=+1$ & none & Thm~\ref{ch14:thm:topone}\\
\lean{moonshine_of_topPairs} & $\Vnat$ from \lean{TopPairs} & \lean{TopPairs} & Thm~\ref{ch14:thm:maintop}\\
\lean{moonshine} & construction: the block VOA on the FLM Leech orbifold space is a simple VOA with $c=24$, $V_1=0$, formal character $J$ (unique witness) & \textbf{none} & \S\ref{ch14:sec:decode}\\
\bottomrule
\end{longtable}

All rows except the conditional reductions are unconditional; every conditional hypothesis listed is
discharged further down the table, ending in \lean{moonshine}, which has none.

\section{Notes and references}\label{ch14:sec:notes}

The triality argument is that of I.~Frenkel, J.~Lepowsky and A.~Meurman, \emph{Vertex Operator
Algebras and the Monster}, Pure and Applied Mathematics 134, Academic Press, 1988, Chapters~12--13
(the moonshine module, the triality $\sigma$, and the proof that the twisted vertex operators satisfy
the Jacobi identity); their announcement is \emph{A natural representation of the Fischer--Griess
Monster with the modular function $J$ as character}, Proc.\ Natl.\ Acad.\ Sci.\ USA 81 (1984),
3256--3260.  The block presentation (4.17), (4.22), (4.23) is from L.~Dolan, P.~Goddard and P.~Montague,
\emph{Conformal field theory of twisted vertex operators}, Nucl.\ Phys.\ B 338 (1990) 529--601, and
\emph{Conformal field theories, representations and lattice constructions}, Comm.\ Math.\ Phys.\ 179
(1996) 61--120 (arXiv:hep-th/9410029).  The bimodule $A(M)$ is due to I.~B.~Frenkel and Y.~Zhu,
\emph{Vertex operator algebras associated to representations of affine and Virasoro algebras}, Duke
Math.\ J.\ 66 (1992) 123--168 (Lemma 2.1.2 and Theorem 1.5.3 are the facts used here); Zhu's algebra is
from Y.~Zhu, \emph{Modular invariance of characters of vertex operator algebras}, J.\ Amer.\ Math.\ Soc.\
9 (1996) 237--302.  Uniqueness of intertwining operators from top-level data, and the correction of
the Frenkel--Zhu fusion theorem, are due to H.~Li, \emph{Determining fusion rules by $A(V)$-modules and
bimodules}, J.\ Algebra 212 (1999) 515--556 (arXiv:q-alg/9706012).  Intertwining operators and their
skew-symmetry: I.~B.~Frenkel, Y.-Z.~Huang and J.~Lepowsky, \emph{On axiomatic approaches to vertex
operator algebras and modules}, Mem.\ Amer.\ Math.\ Soc.\ 104 (1993), no.\ 494; Y.-Z.~Huang, \emph{A
nonmeromorphic extension of the Moonshine module vertex operator algebra}, Contemp.\ Math.\ 193 (1996)
123--148 (arXiv:hep-th/9406190), which gives an alternative proof of the $\Vnat$ construction via
intertwining operator algebras.  Fusion rules for $V_L^+$: T.~Abe, C.~Dong and H.~Li, \emph{Fusion rules
for the vertex operator algebras $M(1)^+$ and $V_L^+$}, Comm.\ Math.\ Phys.\ 253 (2005) 171--219.  The
uniqueness results mentioned in \S\ref{ch14:sec:decode}: C.~Dong, R.~L.~Griess Jr.\ and C.~H.~Lam,
\emph{Uniqueness results for the moonshine vertex operator algebra}, Amer.\ J.\ Math.\ 129 (2007)
583--609.  Uniqueness of the vertex operator algebra structure on a simple-current extension, which
identifies the object constructed here with FLM's $\Vnat$ on paper: C.~Dong and G.~Mason,
\emph{Rational vertex operator algebras and the effective central charge}, Int.\ Math.\ Res.\ Not.\ 2004,
no.~56, 2989--3008 (arXiv:math/0201318), Proposition~5.3, applied with $G=\mathbb Z/2$, $V^0=V_\Leech^+$,
$V^1=(V_\Leech^T)^+$ and the simple-current fusion rules of Abe--Dong--Li; see also C.~Dong, H.~Li and G.~Mason, \emph{Simple currents and extensions of vertex operator
algebras}, Comm.\ Math.\ Phys.\ 180 (1996) 671--707.  For physics background on orbifolds and twist fields, L.~Dixon, D.~Friedan, E.~Martinec and
S.~Shenker, \emph{The conformal field theory of orbifolds}, Nucl.\ Phys.\ B 282 (1987) 13--73.  The
project's own route notes are \leanpath{Notes/M8GluingPlan.md}, \leanpath{Notes/SigmaKey.md},
\leanpath{Notes/KeySkew.md} and \leanpath{Notes/BimoduleGenTest.md}.
